\kefalaio{Βασικές γεωμετρικές έννοιες}{Πριν χτίσουμε, πρέπει να ονομάσουμε τα υλικά}{%
\begin{tabular}{@{}l@{\hspace{1em}}p{0.82\linewidth}@{}}
1.1 & Σημείο, ευθύγραμμο τμήμα, ευθεία, ημιευθεία, επίπεδο, ημιεπίπεδο\\[2pt]
1.2 & Γωνία, γραμμή, επίπεδα και ευθύγραμμα σχήματα, ίσα σχήματα\\[2pt]
1.3 & Μέτρηση, σύγκριση και ισότητα τμημάτων, απόσταση, μέσο\\[2pt]
1.4 & Πρόσθεση και αφαίρεση ευθυγράμμων τμημάτων\\[2pt]
1.5 & Μέτρηση, σύγκριση και ισότητα γωνιών, διχοτόμος\\[2pt]
1.6 & Είδη γωνιών, κάθετες ευθείες\\[2pt]
1.7 & Εφεξής και διαδοχικές γωνίες, άθροισμα γωνιών\\[2pt]
1.8 & Παραπληρωματικές, συμπληρωματικές και κατακορυφήν γωνίες\\[2pt]
1.9 & Θέσεις ευθειών στο επίπεδο\\[2pt]
1.10 & Απόσταση σημείου από ευθεία, απόσταση παραλλήλων\\[2pt]
1.11 & Κύκλος και στοιχεία του κύκλου\\[2pt]
1.12 & Επίκεντρη γωνία, τόξο, μέτρηση τόξου\\[2pt]
1.13 & Θέσεις ευθείας και κύκλου\\[2pt]
 & Επανάληψη του κεφαλαίου\\
\end{tabular}}
\paragrafos{Βασικές έννοιες: σημείο, ευθεία, ημιευθεία, ευθύγραμμο τμήμα, γωνία, κύκλος}{Β.1.1, «Βασικές έννοιες: σημείο, ευθεία, ημιευθεία, ευθύγραμμο τμήμα, γωνία, κύκλος»}

\begin{anagke}
Η κυρία Αλεξάνδρα ανακοινώνει ένα μεγάλο εγχείρημα για τη Γεωμετρία της χρονιάς: «Ο κήπος του σχολείου θα ανανεωθεί και η Α3 θα τον σχεδιάσει. Θα φέρουμε τα δέντρα, τα μονοπάτια, τον φράχτη, ένα στρογγυλό παρτέρι. Αλλά πρώτα πρέπει να μάθουμε να μιλάμε για αυτά που θα σχεδιάσουμε». Η Αγγελική, που βλέπει σχήματα παντού, χαίρεται: «Επιτέλους ένα μάθημα για αυτό που κάνω όλη μέρα!».

Η Αγγελική παίρνει ένα χαρτί και κάνει μια κάτοψη του κήπου. Κάθε δέντρο το σημειώνει με μια τελεία. Το μονοπάτι ανάμεσα σε δύο δέντρα το ενώνει με ένα τεντωμένο νήμα. Ο παππούς Γιάννης, που έρχεται να βοηθήσει, λέει: «Για να φυτέψετε σε ευθεία γραμμή, τεντώστε σχοινί και προχωρήστε όσο θέλετε προς τις δύο μεριές». Ο Στέλιος ρωτάει: «Μέχρι πού; Ο φράχτης τελειώνει στον τοίχο!». Η Ρήνη δείχνει τον φακό της: η δέσμη του φωτός ξεκινά από ένα σημείο και συνεχίζει προς μία μεριά χωρίς να τελειώνει. Για το στρογγυλό παρτέρι ο παππούς καρφώνει ένα παλούκι, δένει ένα σχοινί και γυρίζει τεντωμένο γύρω του. «Όλα τα σημεία της γραμμής απέχουν το ίδιο από το παλούκι», παρατηρεί η Φανή και το γράφει στο τετράδιό της.

Κάθε μία από αυτές τις εικόνες είναι μια γεωμετρική έννοια: σημείο, ευθύγραμμο τμήμα, ευθεία, ημιευθεία, γωνία, κύκλος. Σε αυτή την παράγραφο τις ονομάζουμε προσεκτικά και μαθαίνουμε να τις γράφουμε με σύμβολα.
\end{anagke}

\begin{lexi}{σημείο, ευθεία, επίπεδο, κύκλος}{σημείο, ευθεία, επίπεδο, κύκλος}
Το \lexiword{σημείο} είναι «αυτό που σημαδεύεται»: ένα σημάδι χωρίς διαστάσεις. Η \lexiword{ευθεία} πηγαίνει «ευθύ», χωρίς να στρίβει. Το \lexiword{επίπεδο} είναι ό,τι είναι «επί του πέδου», ίσιο σαν το δάπεδο. Το \lexiword{ευθύγραμμο τμήμα} είναι ένα «κομμένο» κομμάτι ευθείας. Το \lexiword{ημι-} σημαίνει «μισό»: ημιευθεία, ημιεπίπεδο. Ο \lexiword{κύκλος} είναι ό,τι «κυλάει» και γυρίζει. Η \lexiword{ακτίνα} είναι η «αχτίδα» που ξεκινά από το κέντρο.
\end{lexi}

\begin{metafrasi}
Το ίχνος της μύτης του μολυβιού. & σημείο $\mathrm{Α}$\\
Τεντωμένη κλωστή με άκρα $\mathrm{Α}$ και $\mathrm{Β}$. & ευθύγραμμο τμήμα $\mathrm{ΑΒ}$\\
Το ίδιο τμήμα προεκτεινόμενο απεριόριστα και προς τις δύο μεριές. & ευθεία $\mathrm{ΑΒ}$ ή $\varepsilon$ ή $x'x$\\
Φως φακού που ξεκινά από ένα σημείο. & ημιευθεία $\mathrm{Αx}$\\
Δύο ημιευθείες με κοινή αρχή $\mathrm{Ο}$. & γωνία $x\hat{\mathrm{Ο}}y$\\
Τα σημεία που απέχουν $\rho$ από το $\mathrm{Ο}$. & κύκλος $(\mathrm{Ο},\rho)$\\
\end{metafrasi}

\tmima{Η θεωρία}

\tmimasmall{Σημείο και ευθεία}

Ένα \textbf{σημείο} παριστάνεται με μία τελεία και δεν έχει διαστάσεις (ούτε μήκος ούτε πλάτος). Το ίχνος της μύτης του μολυβιού μας δίνει την έννοια του σημείου. Το συμβολίζουμε με ένα κεφαλαίο γράμμα: $\mathrm{Α}$, $\mathrm{Β}$, $\Gamma$.

\textbf{Ευθεία} είναι η γραμμή που εκτείνεται απεριόριστα, χωρίς αρχή και τέλος, και είναι απολύτως ίσια. Τη συμβολίζουμε με ένα μικρό γράμμα, π.χ. $\varepsilon$, ή με δύο γράμματα, π.χ. $x'x$.

\begin{idiotita}{ευθεία και σημεία}
Από ένα σημείο διέρχονται \textbf{άπειρες} ευθείες. Από δύο διαφορετικά σημεία διέρχεται \textbf{μία και μοναδική} ευθεία.
\end{idiotita}

\begin{center}
\begin{tikzpicture}[scale=0.9]
\begin{scope}
\foreach \a in {0,35,70,105,140,-35} {\draw[gblue] ({-1.5*cos(\a)},{-1.5*sin(\a)})--({1.5*cos(\a)},{1.5*sin(\a)});}
\fill (0,0) circle (2pt) node[below right] {$\mathrm{Α}$};
\node at (0,-1.9) {\small άπειρες ευθείες};
\end{scope}
\begin{scope}[xshift=5.5cm]
\draw[gblue,{Stealth}-{Stealth}] (-1.7,0)--(1.7,0);
\fill (-0.8,0) circle (2pt) node[above] {$\mathrm{Α}$}; \fill (0.8,0) circle (2pt) node[above] {$\mathrm{Β}$};
\node at (0,-0.7) {\small μία ευθεία};
\end{scope}
\end{tikzpicture}
\end{center}

\tmimasmall{Ευθύγραμμο τμήμα και ημιευθεία}

Αν σε μια ευθεία πάρουμε δύο σημεία $\mathrm{Α}$ και $\mathrm{Β}$, το σχήμα που αποτελείται από τα $\mathrm{Α}$, $\mathrm{Β}$ και όλα τα σημεία ανάμεσά τους λέγεται \textbf{ευθύγραμμο τμήμα} $\mathrm{ΑΒ}$ (ή $\mathrm{ΒΑ}$). Τα $\mathrm{Α}$ και $\mathrm{Β}$ είναι τα \textbf{άκρα} του. Το τμήμα έχει αρχή και τέλος, ενώ η ευθεία όχι.

Αν προεκτείνουμε το τμήμα $\mathrm{ΑΒ}$ και προς τα δύο άκρα του, σχηματίζεται ευθεία. Αν το προεκτείνουμε \emph{μόνο προς το ένα} άκρο, το σχήμα έχει αρχή αλλά δεν έχει τέλος και λέγεται \textbf{ημιευθεία}. Τη συμβολίζουμε $\mathrm{Αx}$, με πρώτο το γράμμα της αρχής.

\begin{center}
\begin{tikzpicture}[scale=0.95]
\draw[thick] (0,0)--(3,0); \fill (0,0) circle (2pt) node[below] {$\mathrm{Α}$}; \fill (3,0) circle (2pt) node[below] {$\mathrm{Β}$};
\node at (1.5,-0.9) {τμήμα $\mathrm{ΑΒ}$};
\begin{scope}[xshift=4.8cm]
\draw[thick,{Stealth}-{Stealth}] (-0.5,0)--(3.5,0); \fill (0.5,0) circle (2pt) node[below] {$\mathrm{Α}$}; \fill (2.7,0) circle (2pt) node[below] {$\mathrm{Β}$}; \node[above] at (3.2,0) {$\varepsilon$};
\node at (1.5,-0.9) {ευθεία $\varepsilon$};
\end{scope}
\begin{scope}[xshift=9.8cm]
\draw[thick,-{Stealth}] (0,0)--(3.2,0) node[above] {$x$}; \fill (0,0) circle (2pt) node[below] {$\mathrm{Α}$};
\node at (1.5,-0.9) {ημιευθεία $\mathrm{Αx}$};
\end{scope}
\end{tikzpicture}
\end{center}

Ένα σημείο $\mathrm{Α}$ μιας ευθείας $x'x$ τη χωρίζει σε δύο ημιευθείες, τις $\mathrm{Αx}$ και $\mathrm{Αx'}$, που έχουν κοινή αρχή το $\mathrm{Α}$. Οι ημιευθείες αυτές λέγονται \textbf{αντικείμενες}.

\begin{center}
\begin{tikzpicture}
\draw[thick,{Stealth}-{Stealth}] (-3,0)--(3,0); \fill (0,0) circle (2pt) node[above] {$\mathrm{Α}$};
\node[above] at (-2.8,0) {$x'$}; \node[above] at (2.8,0) {$x$};
\node[below,gblue] at (-1.5,0) {\small $\mathrm{Αx'}$}; \node[below,gblue] at (1.5,0) {\small $\mathrm{Αx}$};
\end{tikzpicture}
\end{center}

\tmimasmall{Γωνία}

Σχεδιάζουμε δύο ημιευθείες $\mathrm{Οx}$ και $\mathrm{Οy}$ με κοινή αρχή το σημείο $\mathrm{Ο}$. Οι ημιευθείες χωρίζουν το επίπεδο σε δύο περιοχές. Καθεμία από αυτές τις περιοχές, μαζί με τις δύο ημιευθείες, ονομάζεται \textbf{γωνία}. Το $\mathrm{Ο}$ είναι η \textbf{κορυφή} και οι ημιευθείες οι \textbf{πλευρές} της γωνίας.
\begin{itemize}[leftmargin=1.2em,itemsep=0.3ex,topsep=2pt]
\item Η μικρότερη περιοχή δίνει την \textbf{κυρτή} γωνία.
\item Η μεγαλύτερη περιοχή δίνει τη \textbf{μη κυρτή} γωνία.
\end{itemize}

\begin{center}
\begin{tikzpicture}[scale=1.1]
\fill[gblueL] (0,0)--(2.5,0) arc(0:38:2.5)--cycle;
\draw[thick] (2.9,0)--(0,0)--(2.3,1.8); \fill (0,0) circle (2pt) node[below left] {$\mathrm{Ο}$};
\node[right] at (2.9,0) {$x$}; \node[above right] at (2.3,1.8) {$y$};
\draw[gblue] (0.8,0) arc(0:38:0.8); \node[gblue] at (1.15,0.33) {\small κυρτή};
\draw[ggray] (-0.6,0) arc(180:322:0.6); \node[ggray] at (-1.15,-0.6) {\small μη κυρτή};
\end{tikzpicture}
\end{center}

Συμβολίζουμε τη γωνία με τρία γράμματα, γράφοντας στη μέση το γράμμα της κορυφής: $x\hat{\mathrm{Ο}}y$ ή $y\hat{\mathrm{Ο}}x$. Μπορούμε επίσης να τη συμβολίσουμε $\hat{\mathrm{Ο}}$ (με το γράμμα της κορυφής) ή με ένα μικρό ελληνικό γράμμα, π.χ. $\hat{\omega}$.

\textbf{Σημείωση.} Με τον γνώμονα κατασκευάζουμε ευθείες που λέγονται \textbf{κάθετες}. Η γωνία που σχηματίζεται λέγεται \textbf{ορθή}.

\tmimasmall{Επίπεδο και ημιεπίπεδο}

\textbf{Επίπεδο} είναι η επιφάνεια πάνω στην οποία εφαρμόζει παντού η ευθεία γραμμή. Το επίπεδο επεκτείνεται απεριόριστα και το ονομάζουμε με ένα κεφαλαίο γράμμα, π.χ. $\Pi$. Κάθε ευθεία $\varepsilon$ ενός επιπέδου το χωρίζει σε δύο \textbf{ημιεπίπεδα}.

\begin{center}
\begin{tikzpicture}[scale=0.95]
\fill[gblueL] (0,0)--(4.4,0)--(5.4,1.8)--(1,1.8)--cycle;
\fill[gwarm] (0,0)--(4.4,0)--(3.9,-0.9)--(-0.5,-0.9)--cycle;
\draw (-0.5,-0.9)--(3.9,-0.9)--(5.4,1.8)--(1,1.8)--cycle;
\draw[thick,gblue] (-0.25,-0.45)--(5.15,1.35) coordinate (e);
\node[gblue] at (5.5,1.5) {$\varepsilon$};
\node at (1.4,1.1) {$\Pi_1$}; \node at (3.5,-0.45) {$\Pi_2$};
\end{tikzpicture}
\end{center}

\tmimasmall{Κύκλος}

\begin{orismos}{κύκλος}
\textbf{Κύκλος} λέγεται το σύνολο όλων των σημείων του επιπέδου που απέχουν την ίδια απόσταση από ένα σταθερό σημείο $\mathrm{Ο}$.
\end{orismos}

Το σημείο $\mathrm{Ο}$ λέγεται \textbf{κέντρο} του κύκλου. Η σταθερή απόσταση λέγεται \textbf{ακτίνα} του κύκλου και συμβολίζεται με $\rho$. Έναν κύκλο με κέντρο $\mathrm{Ο}$ και ακτίνα $\rho$ τον γράφουμε $(\mathrm{Ο},\rho)$. Τον σχεδιάζουμε με τον \textbf{διαβήτη}. Οι κύκλοι που έχουν το ίδιο κέντρο λέγονται \textbf{ομόκεντροι}.

\begin{center}
\begin{tikzpicture}
\draw[thick] (0,0) circle (1.2); \draw[thick] (0,0) circle (1.9);
\fill (0,0) circle (2pt) node[below left] {$\mathrm{Ο}$};
\draw[gblue,thick] (0,0)--(0.85,0.85) node[midway,above left] {$\rho$};
\fill (0.85,0.85) circle (2pt) node[above right] {$\mathrm{Α}$};
\node at (0,-2.4) {\small δύο ομόκεντροι κύκλοι};
\end{tikzpicture}
\end{center}

\begin{prosoxi}
\begin{itemize}[leftmargin=1.2em,itemsep=0.3ex,topsep=2pt]
\item Το τμήμα $\mathrm{ΑΒ}$ και το $\mathrm{ΒΑ}$ είναι το ίδιο τμήμα. Αλλά οι ημιευθείες $\mathrm{Αx}$ και $\mathrm{Βx}$ είναι διαφορετικές, γιατί έχουν διαφορετική αρχή.
\item Το τμήμα έχει άκρα και μήκος. Η ευθεία και η ημιευθεία δεν έχουν δύο άκρα και δεν έχουν μήκος.
\item Στο όνομα της γωνίας το γράμμα της κορυφής γράφεται πάντα στη μέση: $x\hat{\mathrm{Ο}}y$, όχι $\mathrm{Ο}xy$.
\item Ο κύκλος είναι μόνο η γραμμή, όχι το εσωτερικό του. Όλες οι ακτίνες του ίδιου κύκλου είναι ίσες.
\end{itemize}
\end{prosoxi}

\tmima{Λυμένα παραδείγματα}

\paradeigma{Πόσες ευθείες;}
\tixero Τέσσερα δέντρα $\mathrm{Α},\,\mathrm{Β},\,\Gamma,\,\Delta$ του κήπου, από τα οποία τρία οποιαδήποτε δεν βρίσκονται στην ίδια ευθεία.\par
\tithelo Πόσες ευθείες ορίζουν.\par
\possindeo Κάθε ζεύγος σημείων ορίζει μία μοναδική ευθεία. Τα ζεύγη είναι $\mathrm{ΑΒ},\,\mathrm{ΑΓ},\,\mathrm{ΑΔ},\,\mathrm{ΒΓ},\,\mathrm{ΒΔ},\,\Gamma\Delta$. Με τον τρόπο των χειραψιών: $\dfrac{4\cdot3}{2}=6$.\par
\apantisi $6$ ευθείες.\par
\elegxos Αν το $\mathrm{Δ}$ ήταν πάνω στην ευθεία $\mathrm{ΑΒ}$, οι ευθείες θα ήταν λιγότερες.

\paradeigma{Ονόματα ημιευθειών}
\tixero Το σημείο $\mathrm{Ο}$ βρίσκεται πάνω στην ευθεία $x'x$. Τα $\mathrm{Α}$ και $\mathrm{Β}$ είναι δεξιά και αριστερά του $\mathrm{Ο}$ αντίστοιχα.\par
\tithelo Ποιες είναι οι δύο ημιευθείες που ορίζει το $\mathrm{Ο}$ και ποια περιέχει το $\mathrm{Α}$.\par
\possindeo Το $\mathrm{Ο}$ χωρίζει την ευθεία σε δύο ημιευθείες με αρχή το $\mathrm{Ο}$. Η μία πηγαίνει προς το $x$ και η άλλη προς το $x'$.\par
\apantisi Οι $\mathrm{Οx}$ (περιέχει το $\mathrm{Α}$) και $\mathrm{Οx'}$ (περιέχει το $\mathrm{Β}$), αντικείμενες.\par
\elegxos Αν τις ενώσω παίρνω ξανά ολόκληρη την ευθεία.

\paradeigma{Το όνομα μιας γωνίας}
\tixero Οι ημιευθείες $\mathrm{Οx}$ και $\mathrm{Οy}$ με κοινή αρχή $\mathrm{Ο}$ σχηματίζουν τη μικρότερη από τις δύο περιοχές του επιπέδου.\par
\tithelo Να ονομάσω τη γωνία με δύο τρόπους και να πω αν είναι κυρτή.\par
\possindeo Η κορυφή είναι το $\mathrm{Ο}$ και μπαίνει στη μέση: $x\hat{\mathrm{Ο}}y$ ή $y\hat{\mathrm{Ο}}x$. Επίσης μπορώ να γράψω απλώς $\hat{\mathrm{Ο}}$.\par
\apantisi Η γωνία $x\hat{\mathrm{Ο}}y$ είναι κυρτή, γιατί αντιστοιχεί στη μικρότερη περιοχή.\par
\elegxos Η μεγαλύτερη περιοχή θα έδινε τη μη κυρτή γωνία.

\paradeigma{Το μονοπάτι χωρίζει τον κήπο}
\tixero Ένα ευθύγραμμο μονοπάτι $\varepsilon$ χωρίζει τον κήπο σε δύο ημιεπίπεδα. Τα δέντρα $\mathrm{Κ}$ και $\Lambda$ είναι στο ένα και το $\mathrm{Μ}$ στο άλλο.\par
\tithelo Ποιες ευθύγραμμες διαδρομές ανάμεσα σε δύο δέντρα διασχίζουν το μονοπάτι.\par
\possindeo Για να πάω από το ένα ημιεπίπεδο στο άλλο πρέπει να περάσω από την ευθεία $\varepsilon$.\par
\apantisi Τα τμήματα $\mathrm{ΚΜ}$ και $\Lambda\mathrm{Μ}$ διασχίζουν το μονοπάτι, το $\mathrm{ΚΛ}$ όχι.\par
\elegxos Τα $\mathrm{Κ}$ και $\Lambda$ είναι στην ίδια πλευρά.

\paradeigma{Το στρογγυλό παρτέρι}
\tixero Το σχοινί του παππού είναι $1{,}5$ μ. και ο παππούς το κρατά τεντωμένο γύρω από το παλούκι $\mathrm{Ο}$. Ένα δέντρο $\mathrm{Δ}$ βρίσκεται $2$ μ. από το παλούκι.\par
\tithelo Ποιος κύκλος χαράσσεται και αν το δέντρο είναι πάνω στη γραμμή του.\par
\possindeo Το παλούκι είναι το κέντρο και το σχοινί η ακτίνα: κύκλος $(\mathrm{Ο},1{,}5)$ μ. Το δέντρο θα ήταν πάνω στον κύκλο αν $\mathrm{ΟΔ}=1{,}5$ μ.\par
\apantisi Ο κύκλος είναι ο $(\mathrm{Ο},1{,}5)$. Το δέντρο έχει $\mathrm{ΟΔ}=2\neq1{,}5$, άρα δεν είναι πάνω στη γραμμή.\par
\elegxos Το δέντρο απέχει περισσότερο από το σχοινί, άρα είναι έξω από το παρτέρι.

\tmima{Ασκήσεις}

\zoni{Το κατάλαβα}{σημεία, ευθείες, ημιευθείες, γωνίες}

\ex Συμπλήρωσε τα κενά. α) Το ίχνος της μύτης του μολυβιού μας δίνει την έννοια του \dots. β) Αν προεκτείνουμε ένα τμήμα και προς τα δύο άκρα του, προκύπτει \dots. γ) Από ένα σημείο διέρχονται \dots ευθείες, ενώ από δύο σημεία διέρχεται \dots ευθεία. δ) Αν προεκτείνουμε ένα τμήμα μόνο προς το ένα άκρο του, προκύπτει \dots. ε) Οι ημιευθείες που έχουν κοινή αρχή και σχηματίζουν ευθεία λέγονται \dots.

\ex Σχεδίασε τρία σημεία $\mathrm{Α},\,\mathrm{Β},\,\Gamma$ που δεν βρίσκονται στην ίδια ευθεία. Φέρε και ονόμασε τις ευθείες που ορίζουν και γράψε όλα τα ευθύγραμμα τμήματα που βλέπεις. Πόσες ευθείες και πόσα τμήματα είναι;

\ex Γράψε με σύμβολα: α) το τμήμα με άκρα $\mathrm{Κ}$ και $\Lambda$ \quad β) η ημιευθεία με αρχή το $\mathrm{Κ}$ που περνά από το $\Lambda$ \quad γ) η ευθεία που ορίζουν τα $\mathrm{Κ}$ και $\Lambda$.

\ex Σχεδίασε ένα σημείο $\mathrm{Ο}$ και δύο διαφορετικές ευθείες που να διέρχονται από αυτό. Ονόμασε τις αντικείμενες ημιευθείες με αρχή το $\mathrm{Ο}$. Πόσες ημιευθείες βλέπεις συνολικά;

\ex Σχεδίασε μια γωνία $x\hat{\mathrm{Ο}}y$ και γράψε το όνομά της με δύο διαφορετικούς τρόπους. Σημείωσε με χρώμα την κυρτή και με άλλο χρώμα τη μη κυρτή γωνία.

\ex Σχεδίασε με διαβήτη τον κύκλο $(\mathrm{Κ},2{,}5)$ cm και τρεις διαφορετικές ακτίνες του. Μέτρησέ τις. Τι παρατηρείς;

\ex Σωστό ή λάθος; α) Η ευθεία έχει δύο άκρα. \quad β) Τα τμήματα $\mathrm{ΑΒ}$ και $\mathrm{ΒΑ}$ είναι ίδια. \quad γ) Οι ημιευθείες $\mathrm{Αx}$ και $\mathrm{Βx}$ είναι ίδιες. \quad δ) Από δύο σημεία περνά μόνο μία ευθεία. \quad ε) Η γωνία έχει δύο ημιευθείες με κοινή αρχή. \quad στ) Οι ομόκεντροι κύκλοι έχουν την ίδια ακτίνα.

\zoni{Το χρησιμοποιώ}{ο κήπος, ο φακός, το στρογγυλό παρτέρι}

\ex \textbf{Η σειρά των δέντρων.} Ο παππούς Γιάννης θέλει να φυτέψει δέντρα σε ευθεία γραμμή. Πόσα παλούκια αρκεί να βάλει για να ορίσει την ευθεία και πώς θα την ελέγξει με ένα σχοινί;

\ex \textbf{Οι ακτίνες.} Ένας φακός φωτίζει με δέσμη που ξεκινά από το σημείο $\mathrm{Α}$ και περνά από το $\mathrm{Β}$. Ποια γεωμετρική έννοια περιγράφει η δέσμη; Αν δύο φακοί στέκονται αντικριστά και φωτίζουν ο ένας τον άλλον, τι σχέση έχουν οι δύο ημιευθείες;

\ex \textbf{Τα σχοινιά.} Πέντε παιδιά στέκονται σε σημεία του κήπου, από τα οποία τρία οποιαδήποτε δεν είναι στην ίδια ευθεία, και ενώνουν κάθε ζευγάρι με σχοινί. Πόσα σχοινιά χρειάζονται;

\ex \textbf{Το μονοπάτι.} Ένα ευθύγραμμο μονοπάτι $\varepsilon$ χωρίζει τον κήπο σε δύο ημιεπίπεδα. Τα δέντρα $\mathrm{Κ}$, $\Lambda$ είναι στην ίδια πλευρά και το $\mathrm{Μ}$ στην άλλη. Ποια από τα τμήματα $\mathrm{ΚΛ}$, $\mathrm{ΚΜ}$, $\Lambda\mathrm{Μ}$ τέμνουν το μονοπάτι;

\ex \textbf{Το παρτέρι.} Η Ρήνη χαράσσει με σχοινί $1{,}2$ μ. ένα στρογγυλό παρτέρι με κέντρο το παλούκι $\mathrm{Ο}$. Τρία φυτά απέχουν από το παλούκι $0{,}9$ μ., $1{,}2$ μ. και $1{,}6$ μ. Ποιο είναι πάνω στη γραμμή του κύκλου και ποιο εκτός του κύκλου;

\zoni{Το σκέφτομαι}{ευθεία ή τμήμα;}

\ex Γιατί η ευθεία δεν έχει μήκος αλλά το τμήμα έχει; Μπορείς να μετρήσεις το μήκος μιας ημιευθείας;

\ex Τρία σημεία $\mathrm{Α},\,\mathrm{Β},\,\Gamma$ βρίσκονται στην ίδια ευθεία με αυτή τη σειρά. Πόσες διαφορετικές ημιευθείες με αρχή ένα από αυτά τα σημεία ορίζονται πάνω στην ευθεία; (Μέτρα και τις δύο κατευθύνσεις.)

\ex Μπορούν δύο διαφορετικές ευθείες να περνούν και οι δύο από τα ίδια δύο σημεία $\mathrm{Α}$ και $\mathrm{Β}$; Εξήγησε με την ιδιότητα της ευθείας.

\ex Δύο ομόκεντροι κύκλοι έχουν ακτίνες $3$ cm και $5$ cm. Μπορεί να υπάρχει σημείο που ανήκει και στους δύο; Γιατί;

\zoni{Για το σπίτι}{γεωμετρία στο δωμάτιο}

\ex Βρες στο δωμάτιό σου πέντε παραδείγματα σημείου, ευθύγραμμου τμήματος και επιπέδου. Πού μοιάζει με ημιευθεία ένα αντικείμενο της καθημερινότητας;

\ex Με νήμα και δύο καρφίτσες σε ένα χαρτόνι δείξε ότι από δύο σημεία περνά μία ευθεία. Μετά σχεδίασε με διαβήτη έναν κύκλο που να περνά από ένα σημείο της επιλογής σου.

\paragrafos{Μήκος ευθύγραμμου τμήματος}{Β.1.2, «Μήκος ευθύγραμμου τμήματος»}

\begin{anagke}
Για να σχεδιάσουν τον κήπο, η παρέα μετράει πρώτα το οικόπεδο. Ο Στέλιος μετράει το μήκος του μονοπατιού με βήματα και βρίσκει $14$. Ο Ηλίας το μετράει με τα δικά του βήματα και βρίσκει $16$. Ποιος έχει δίκιο; «Και οι δύο», λέει η Πολυξένη: «το βήμα του Στέλιου είναι μεγαλύτερο. Έχουμε διαφορετικές μονάδες». Η Αγγελική παίρνει τη μετροταινία του παππού: $9{,}8$ μέτρα. «Αυτό είναι σίγουρο: όλοι θα βρούμε το ίδιο νούμερο».

Μετά η Ρήνη ρωτάει: «Θέλουμε το συντριβάνι να είναι ακριβώς στη μέση του μονοπατιού. Πώς βρίσκουμε τη μέση;». Ο Νίκος πιάνει ένα σχοινί με το μήκος του μονοπατιού και το διπλώνει στα δύο: «Εκεί που διπλώνει είναι η μέση!». Ο Βασίλης προσθέτει: «Και για τον φράχτη γύρω από το παρτέρι πρέπει να προσθέσουμε όλες τις πλευρές του». Η Φανή τα σημειώνει όλα στο τετράδιό της και ρωτάει: «Μα πόσα μέτρα φράχτη είναι αυτά, αν κάποιες πλευρές τις μετρήσαμε σε εκατοστά;».

Σε αυτή την παράγραφο μαθαίνουμε να μετράμε το μήκος ενός τμήματος με σωστές μονάδες, να μετατρέπουμε μονάδες, να βρίσκουμε το μέσο ενός τμήματος και να προσθέτουμε και να αφαιρούμε τμήματα. Έτσι θα υπολογίζουμε και την περίμετρο ενός σχήματος.
\end{anagke}

\begin{lexi}{απόσταση, μέσο, περίμετρος}{απόσταση, μέσο, περίμετρος}
Η \lexiword{απόσταση} δύο σημείων είναι το πόσο «απέχουν». Το \lexiword{μέσο} ενός τμήματος είναι το σημείο που βρίσκεται «στη μέση». Η \lexiword{περίμετρος} είναι το «γύρω μέτρημα» (περί + μέτρο): το συνολικό μήκος του περιγράμματος ενός σχήματος. Το \lexiword{μέτρο} στα αρχαία ελληνικά είναι το «μέτρον», η ποσότητα με την οποία μετράμε.
\end{lexi}

\begin{metafrasi}
Το μήκος του τμήματος $\mathrm{ΑΒ}$ είναι $6{,}4$ εκατοστά. & $\mathrm{ΑΒ}=6{,}4\,\text{cm}$\\
Τα δύο τμήματα έχουν το ίδιο μήκος. & $\mathrm{ΑΒ}=\Gamma\Delta$\\
Η απόσταση των $\mathrm{Α}$ και $\mathrm{Β}$. & το μήκος του $\mathrm{ΑΒ}$\\
Το $\mathrm{Μ}$ είναι το μέσο του $\mathrm{ΑΒ}$. & $\mathrm{ΑΜ}=\mathrm{ΜΒ}$\\
Το $\mathrm{Β}$ είναι ανάμεσα στα $\mathrm{Α}$ και $\Gamma$. & $\mathrm{ΑΓ}=\mathrm{ΑΒ}+\mathrm{ΒΓ}$\\
Το άθροισμα των πλευρών. & περίμετρος\\
\end{metafrasi}

\tmima{Η θεωρία}

\tmimasmall{Μέτρηση μήκους}

Για να συγκρίνουμε δύο μεγέθη, δηλαδή για να δούμε αν το ένα είναι μεγαλύτερο, μικρότερο ή ίσο με το άλλο, τα συγκρίνουμε με ένα σταθερό μέγεθος που λέγεται \textbf{μονάδα μέτρησης}.

\begin{kanonas}{μήκος ευθύγραμμου τμήματος}
Βρίσκουμε το \textbf{μήκος} ενός ευθύγραμμου τμήματος συγκρίνοντάς το με ένα άλλο ευθύγραμμο τμήμα που το λέμε \textbf{μονάδα μήκους}. Η βασική μονάδα μήκους που χρησιμοποιούμε είναι το \textbf{μέτρο} (m).
\end{kanonas}

Για μικρά μήκη χρησιμοποιούμε συνήθως το \textbf{υποδεκάμετρο}, που είναι το ένα δέκατο του μέτρου. Αν το τμήμα είναι πολύ μεγαλύτερο ή πολύ μικρότερο από το μέτρο, διαλέγουμε ένα πολλαπλάσιο ή μια υποδιαίρεση του μέτρου.

\tmimasmall{Μονάδες μέτρησης μήκους}

\renewcommand{\arraystretch}{1.25}
\begin{center}
\begin{tabular}{@{}lll@{}}
\toprule
 & \textbf{Ονομασία} & \textbf{Σύμβολο}\\
\midrule
Πολλαπλάσιο του μέτρου & Χιλιόμετρο & km\\
Βασική μονάδα & Μέτρο & m\\
Υποδιαιρέσεις του μέτρου & Δεκατόμετρο & dm\\
 & Εκατοστόμετρο & cm\\
 & Χιλιοστόμετρο & mm\\
\bottomrule
\end{tabular}
\end{center}

\begin{center}
\begin{tikzpicture}[x=1cm,y=1cm]
\foreach \x/\n in {0/km,2.6/m,5.2/dm,7.8/cm,10.4/mm}{\node[draw,rounded corners=3pt,minimum width=1.5cm,fill=gblueL] (\n) at (\x,0) {\n};}
\draw[-{Stealth},gblue] (km.north) .. controls +(0.4,0.6) and +(-0.4,0.6) .. node[above] {\small $\times1000$} (m.north);
\draw[-{Stealth},gblue] (m.north) .. controls +(0.4,0.6) and +(-0.4,0.6) .. node[above] {\small $\times10$} (dm.north);
\draw[-{Stealth},gblue] (dm.north) .. controls +(0.4,0.6) and +(-0.4,0.6) .. node[above] {\small $\times10$} (cm.north);
\draw[-{Stealth},gblue] (cm.north) .. controls +(0.4,0.6) and +(-0.4,0.6) .. node[above] {\small $\times10$} (mm.north);
\draw[-{Stealth},gwarmD] (mm.south) .. controls +(-0.4,-0.6) and +(0.4,-0.6) .. node[below] {\small $:10$} (cm.south);
\draw[-{Stealth},gwarmD] (cm.south) .. controls +(-0.4,-0.6) and +(0.4,-0.6) .. node[below] {\small $:10$} (dm.south);
\draw[-{Stealth},gwarmD] (dm.south) .. controls +(-0.4,-0.6) and +(0.4,-0.6) .. node[below] {\small $:10$} (m.south);
\draw[-{Stealth},gwarmD] (m.south) .. controls +(-0.4,-0.6) and +(0.4,-0.6) .. node[below] {\small $:1000$} (km.south);
\end{tikzpicture}
\end{center}

Από τη μεγαλύτερη προς τη μικρότερη μονάδα \textbf{πολλαπλασιάζουμε}: $1\,\text{m}=10\,\text{dm}=100\,\text{cm}=1000\,\text{mm}$ και $1\,\text{km}=1000\,\text{m}$. Από τη μικρότερη προς τη μεγαλύτερη \textbf{διαιρούμε}: $1\,\text{mm}=0{,}1\,\text{cm}=0{,}001\,\text{m}$ και $1\,\text{m}=0{,}001\,\text{km}$.

Στην Αγγλία, την Αμερική και σε μερικές ακόμη χώρες η βασική μονάδα μήκους είναι η \textbf{υάρδα} ή γιάρδα (yd). Η $1$ yd διαιρείται σε $3$ πόδια (ft) και ισούται περίπου με $0{,}9144$ m. Το $1$ πόδι διαιρείται σε $12$ ίντσες (in) και ισούται περίπου με $0{,}3048$ m.

\tmimasmall{Ίσα τμήματα, απόσταση, μέσο}

Δύο τμήματα είναι \textbf{ίσα} όταν έχουν το ίδιο μήκος, όταν δηλαδή το ένα εφαρμόζει ακριβώς στο άλλο (π.χ. με τον διαβήτη). Για να συγκρίνουμε δύο μήκη πρέπει να τα εκφράσουμε στην ίδια μονάδα.

\begin{orismos}{απόσταση δύο σημείων}
\textbf{Απόσταση} δύο σημείων $\mathrm{Α}$ και $\mathrm{Β}$ λέγεται το μήκος του ευθύγραμμου τμήματος $\mathrm{ΑΒ}$ που τα ενώνει.
\end{orismos}

\begin{orismos}{μέσο ευθύγραμμου τμήματος}
\textbf{Μέσο} ενός ευθύγραμμου τμήματος $\mathrm{ΑΒ}$ ονομάζουμε το σημείο $\mathrm{Μ}$ του τμήματος που απέχει εξίσου από τα άκρα του: $\mathrm{ΑΜ}=\mathrm{ΜΒ}$.
\end{orismos}

\begin{center}
\begin{tikzpicture}
\draw[thick] (0,0)--(6,0);
\foreach \x/\l in {0/Α,3/Μ,6/Β}{\fill (\x,0) circle (2pt); \node[below] at (\x,-0.1) {$\mathrm{\l}$};}
\node[gblue,above] at (1.5,0.05) {\small $2$ cm}; \node[gblue,above] at (4.5,0.05) {\small $2$ cm};
\end{tikzpicture}
\end{center}

Κάθε τμήμα έχει ένα μοναδικό μέσο, και $\mathrm{ΑΜ}=\mathrm{ΜΒ}=\dfrac{\mathrm{ΑΒ}}{2}$.

\tmimasmall{Άθροισμα και διαφορά τμημάτων}

Αν πάρουμε σε μια ευθεία με τη σειρά τα σημεία $\mathrm{Α}$, $\mathrm{Β}$ και $\Gamma$, τότε:
\begin{itemize}[leftmargin=1.2em,itemsep=0.3ex,topsep=2pt]
\item το τμήμα $\mathrm{ΑΓ}$ είναι το \textbf{άθροισμα} των $\mathrm{ΑΒ}$ και $\mathrm{ΒΓ}$: $\mathrm{ΑΓ}=\mathrm{ΑΒ}+\mathrm{ΒΓ}$.
\item το τμήμα $\mathrm{ΑΒ}$ είναι η \textbf{διαφορά} των $\mathrm{ΑΓ}$ και $\mathrm{ΒΓ}$: $\mathrm{ΑΒ}=\mathrm{ΑΓ}-\mathrm{ΒΓ}$.
\end{itemize}

\begin{center}
\begin{tikzpicture}
\draw[thick] (0,0)--(7,0);
\foreach \x/\l in {0/Α,2.5/Β,7/Γ}{\fill (\x,0) circle (2pt); \node[below] at (\x,-0.1) {$\mathrm{\l}$};}
\draw[gblue,thick] (0,0.4)--(2.5,0.4); \draw[gwarmD,thick] (2.5,0.4)--(7,0.4);
\node[gblue,above] at (1.25,0.45) {$\mathrm{ΑΒ}$}; \node[gwarmD,above] at (4.75,0.45) {$\mathrm{ΒΓ}$};
\draw[thick,ggray] (0,-0.9)--(7,-0.9); \node[ggray,below] at (3.5,-0.95) {$\mathrm{ΑΓ}=\mathrm{ΑΒ}+\mathrm{ΒΓ}$};
\end{tikzpicture}
\end{center}

Για να κατασκευάσουμε με τον διαβήτη ένα τμήμα με μήκος $\mathrm{ΑΒ}+\Gamma\Delta$, σχεδιάζουμε μια ευθεία, παίρνουμε πάνω της ένα τμήμα ίσο με το $\mathrm{ΑΒ}$ και στη συνέχεια, από το άκρο του, ένα τμήμα ίσο με το $\Gamma\Delta$. Για τη διαφορά $\mathrm{ΑΒ}-\Gamma\Delta$ (όταν $\mathrm{ΑΒ}>\Gamma\Delta$) παίρνουμε πάνω στο $\mathrm{ΑΒ}$, από το ένα άκρο του, τμήμα ίσο με το $\Gamma\Delta$ και το κομμάτι που μένει είναι η διαφορά.

\tmimasmall{Περίμετρος ευθύγραμμου σχήματος}

\textbf{Κυρτό} ονομάζεται ένα ευθύγραμμο σχήμα όταν η προέκταση κάθε πλευράς του αφήνει το σχήμα στο ίδιο ημιεπίπεδο. Διαφορετικά ονομάζεται \textbf{μη κυρτό}.

\begin{center}
\begin{tikzpicture}[scale=0.9]
\coordinate (A) at (0,0); \coordinate (B) at (2.5,0); \coordinate (C) at (3.1,1.8); \coordinate (D) at (1.2,2.8); \coordinate (E) at (-0.5,1.5);
\draw[thick,fill=gblueL] (A)--(B)--(C)--(D)--(E)--cycle;
\draw[dashed,ggray] ($(A)!-0.5!(B)$)--($(A)!1.7!(B)$);
\node at (1.3,-0.8) {\small κυρτό};
\begin{scope}[xshift=6cm]
\coordinate (P) at (0,0); \coordinate (Q) at (3,0); \coordinate (R) at (3,2.2); \coordinate (S) at (1.6,0.9); \coordinate (T) at (0,2.2);
\draw[thick,fill=gwarm] (P)--(Q)--(R)--(S)--(T)--cycle;
\draw[dashed,ggray] ($(S)!-0.7!(R)$)--($(S)!2.2!(R)$);
\node at (1.5,-0.8) {\small μη κυρτό};
\end{scope}
\end{tikzpicture}
\end{center}

\begin{orismos}{περίμετρος}
Το άθροισμα των πλευρών ενός ευθύγραμμου σχήματος λέγεται \textbf{περίμετρος} του σχήματος.
\end{orismos}

Για παράδειγμα, αν ένα τετράπλευρο έχει πλευρές $4$ cm, $2{,}6$ cm, $3{,}1$ cm και $3$ cm, η περίμετρός του είναι $4+2{,}6+3{,}1+3=12{,}7$ cm.

\begin{prosoxi}
\begin{itemize}[leftmargin=1.2em,itemsep=0.3ex,topsep=2pt]
\item Για να συγκρίνω ή να προσθέσω μήκη πρέπει πρώτα να τα φέρω στην ίδια μονάδα: τα $0{,}8$ m δεν είναι μικρότερα από τα $75$ cm, γιατί $0{,}8$ m $=80$ cm.
\item Στις μετατροπές μεγάλη $\to$ μικρή μονάδα πολλαπλασιάζω, και μικρή $\to$ μεγάλη διαιρώ. Αν ο αριθμός που βγάζω είναι μεγαλύτερος για τη μεγαλύτερη μονάδα, κάτι πήγε στραβά.
\item Το μέσο \emph{ανήκει} στο τμήμα: δεν αρκεί να απέχει εξίσου από τα άκρα.
\item Η σχέση $\mathrm{ΑΓ}=\mathrm{ΑΒ}+\mathrm{ΒΓ}$ ισχύει όταν το $\mathrm{Β}$ βρίσκεται \emph{ανάμεσα} στα $\mathrm{Α}$ και $\Gamma$ πάνω στην ίδια ευθεία.
\item Η περίμετρος μετράει μόνο το περίγραμμα, όχι την επιφάνεια του σχήματος.
\end{itemize}
\end{prosoxi}

\tmima{Λυμένα παραδείγματα}

\paradeigma{Μετατροπές μονάδων}
\tixero $2{,}4$ km, $35$ cm και $7$ dm.\par
\tithelo Να τα μετατρέψω σε m, m και mm αντίστοιχα.\par
\possindeo Από km σε m πολλαπλασιάζω επί $1000$. Από cm σε m διαιρώ με $100$. Από dm σε mm πολλαπλασιάζω επί $100$.\par
\apantisi $2{,}4$ km $=2400$ m, $35$ cm $=0{,}35$ m, $7$ dm $=700$ mm.\par
\elegxos Το $0{,}35$ m είναι λιγότερο από ένα μέτρο, όπως τα $35$ cm.

\paradeigma{Ποιο είναι μεγαλύτερο;}
\tixero Δύο ράφια έχουν μήκη $0{,}8$ m και $75$ cm.\par
\tithelo Ποιο είναι μεγαλύτερο;\par
\possindeo Τα φέρνω στην ίδια μονάδα: $0{,}8$ m $=80$ cm.\par
\apantisi Μεγαλύτερο είναι το πρώτο ($80$ cm $>75$ cm).\par
\elegxos Σε mm: $800>750$.

\paradeigma{Το συντριβάνι}
\tixero Το μονοπάτι είναι $9{,}8$ μ. και το συντριβάνι θα μπει στο μέσο του.\par
\tithelo Πόσο απέχει το συντριβάνι από κάθε άκρο.\par
\possindeo Το μέσο $\mathrm{Μ}$ έχει $\mathrm{ΑΜ}=\dfrac{\mathrm{ΑΒ}}{2}=\dfrac{9{,}8}{2}$.\par
\apantisi $4{,}9$ μ. από κάθε άκρο.\par
\elegxos $4{,}9+4{,}9=9{,}8$.

\paradeigma{Τα δέντρα στη σειρά}
\tixero Τα δέντρα $\mathrm{Α},\,\mathrm{Β},\,\Gamma,\,\Delta$ είναι με αυτή τη σειρά σε ευθεία και $\mathrm{ΑΒ}=3$ μ., $\mathrm{ΒΓ}=4$ μ., $\Gamma\Delta=2{,}5$ μ.\par
\tithelo Τα $\mathrm{ΑΓ}$, $\mathrm{ΑΔ}$ και $\mathrm{ΒΔ}$.\par
\possindeo $\mathrm{ΑΓ}=\mathrm{ΑΒ}+\mathrm{ΒΓ}=3+4=7$. $\mathrm{ΑΔ}=\mathrm{ΑΓ}+\Gamma\Delta=7+2{,}5=9{,}5$. $\mathrm{ΒΔ}=\mathrm{ΑΔ}-\mathrm{ΑΒ}=9{,}5-3=6{,}5$.\par
\apantisi $\mathrm{ΑΓ}=7$ μ., $\mathrm{ΑΔ}=9{,}5$ μ. και $\mathrm{ΒΔ}=6{,}5$ μ.\par
\elegxos $\mathrm{ΒΔ}=\mathrm{ΒΓ}+\Gamma\Delta=4+2{,}5=6{,}5$.

\paradeigma{Περίμετρος με μικτές μονάδες}
\tixero Ένα τρίγωνο έχει πλευρές $4$ cm, $35$ mm και $0{,}06$ m.\par
\tithelo Η περίμετρός του σε cm.\par
\possindeo Τα φέρνω όλα σε cm: $35$ mm $=3{,}5$ cm και $0{,}06$ m $=6$ cm. Προσθέτω: $4+3{,}5+6$.\par
\apantisi $13{,}5$ cm.\par
\elegxos Αν είχα προσθέσει $4+35+0{,}06$ χωρίς μετατροπή, θα έβγαινε παράλογος αριθμός.

\paradeigma{Η περίμετρος του παρτεριού}
\tixero Ορθογώνιο παρτέρι με πλευρές $4{,}5$ μ. και $3{,}2$ μ.\par
\tithelo Πόσος φράχτης χρειάζεται για να το περιβάλει, και τι γίνεται αν διπλασιάσουμε κάθε πλευρά.\par
\possindeo Η περίμετρος είναι το άθροισμα των πλευρών: $4{,}5+3{,}2+4{,}5+3{,}2=15{,}4$. Με διπλάσιες πλευρές: $9+6{,}4+9+6{,}4=30{,}8$.\par
\apantisi $15{,}4$ μ. και, με διπλάσιες πλευρές, $30{,}8$ μ., δηλαδή διπλάσια περίμετρος.\par
\elegxos $2\cdot15{,}4=30{,}8$.

\tmima{Ασκήσεις}

\zoni{Το κατάλαβα}{μονάδες, μέσο, άθροισμα, περίμετρος}

\ex Ποια μονάδα (km, m, dm, cm, mm) είναι προτιμότερη για τις μετρήσεις; α) απόσταση δύο πόλεων \quad β) μήκος κρεβατιού \quad γ) πλάτος παραθύρου \quad δ) πάχος τετραδίου \quad ε) μήκος ποδηλατόδρομου.

\ex Μετάτρεψε σε m: α) $1{,}78$ km \quad β) $0{,}12$ dm \quad γ) $21{,}3$ cm \quad δ) $0{,}03$ km \quad ε) $2$ dm \quad στ) $10$ mm.

\ex Μετάτρεψε σε cm: α) $0{,}7$ m \quad β) $45$ mm \quad γ) $1{,}2$ dm \quad δ) $0{,}03$ km.

\ex Σύγκρινε και βάλε σε αύξουσα σειρά τα μήκη: $0{,}35$ m, $3{,}4$ dm, $36$ cm, $345$ mm.

\ex Σχεδίασε ένα ευθύγραμμο τμήμα $\mathrm{ΑΒ}=6{,}5$ cm και βρες το μέσο του $\mathrm{Μ}$. Μέτρησε τα $\mathrm{ΑΜ}$ και $\mathrm{ΜΒ}$ και σύγκρινέ τα.

\ex Τα $\mathrm{Α},\,\mathrm{Β},\,\Gamma$ είναι με αυτή τη σειρά σε μια ευθεία. α) Αν $\mathrm{ΑΒ}=4{,}2$ cm και $\mathrm{ΒΓ}=3{,}5$ cm, βρες το $\mathrm{ΑΓ}$. β) Αν $\mathrm{ΑΓ}=10$ cm και $\mathrm{ΑΒ}=6{,}5$ cm, βρες το $\mathrm{ΒΓ}$.

\ex Σχεδίασε τμήματα $\mathrm{ΑΒ}=5$ cm, $\Gamma\Delta=3$ cm και $\mathrm{ΕΖ}=2$ cm. Κατασκεύασε με τον διαβήτη τμήμα με μήκος: α) $\mathrm{ΑΒ}+\Gamma\Delta$ \quad β) $\mathrm{ΑΒ}-\mathrm{ΕΖ}$ \quad γ) $\mathrm{ΑΒ}+\Gamma\Delta-\mathrm{ΕΖ}$. Έλεγξε με τον χάρακα.

\ex Βρες την περίμετρο: α) ορθογωνίου με πλευρές $6$ και $2{,}5$ \quad β) τριγώνου με πλευρές $5$, $6$, $7{,}5$ \quad γ) τετραγώνου με πλευρά $2{,}8$ \quad δ) πενταγώνου με όλες τις πλευρές ίσες με $4$.

\zoni{Το χρησιμοποιώ}{βήματα, σχοινιά, φράχτες, δρόμοι}

\ex \textbf{Τα βήματα.} Ο Στέλιος μετράει το μονοπάτι με $14$ βήματα και ο Ηλίας με $16$. Γιατί βγαίνουν διαφορετικοί αριθμοί; Αν το μονοπάτι είναι $9{,}8$ μ., πόσο είναι το βήμα του καθενός;

\ex \textbf{Το συντριβάνι.} Ένα άλλο μονοπάτι είναι $12{,}6$ μ. και το συντριβάνι μπαίνει στο μέσο του. Πόσο απέχει από κάθε άκρο; Ένα παγκάκι βρίσκεται πάνω στο μονοπάτι, $1{,}5$ μ. από το συντριβάνι προς το ένα άκρο. Πόσο απέχει το παγκάκι από αυτό το άκρο;

\ex \textbf{Το κορδόνι.} Η θεία Μαρία πουλά κορδόνι για κήπο προς $1{,}80$ € το μέτρο. Η Ρήνη αγοράζει $3$ μ. και $45$ cm. Πόσα μέτρα αγόρασε και πόσα ευρώ πλήρωσε;

\ex \textbf{Ο φράχτης.} Το ορθογώνιο παρτέρι έχει πλευρές $4{,}5$ μ. και $3{,}2$ μ. Ο φράχτης κοστίζει $3$ € το μέτρο και ένα μέτρο μένει ανοιχτό για την είσοδο. Πόσο θα κοστίσει;

\ex \textbf{Η διαδρομή.} Ο Ηλίας πηγαίνει από το σπίτι στο σχολείο με στάσεις: $250$ μ., $180$ μ. και $320$ μ. Πόσα μέτρα περπάτησε και πόσα χιλιόμετρα είναι;

\ex \textbf{Αμερικανικά μέτρα.} Ένα ράφι σε αμερικανικό κατάστημα έχει μήκος $4$ ft. Πόσες ίντσες είναι; Πόσα μέτρα είναι περίπου (με $1$ ft $\approx0{,}3048$ m);

\zoni{Το σκέφτομαι}{ίδια μονάδα, διπλασιάζω, αληθές ή ψευδές}

\ex Ο Βασίλης λέει ότι $0{,}8$ m $<75$ cm, επειδή το $0{,}8$ είναι μικρότερο από το $75$. Έχει δίκιο; Πώς θα του το εξηγούσες;

\ex Το $\mathrm{Μ}$ είναι το μέσο του $\mathrm{ΑΒ}=12$ cm και το $\mathrm{Ν}$ το μέσο του $\mathrm{ΑΜ}$. Πόσο είναι τα $\mathrm{ΑΝ}$ και $\mathrm{ΝΒ}$;

\ex Σωστό ή λάθος; α) Το χιλιοστόμετρο είναι μεγαλύτερη μονάδα από το εκατοστόμετρο. \quad β) Αν $\mathrm{Μ}$ μέσο του $\mathrm{ΑΒ}$ τότε $\mathrm{ΑΜ}=\mathrm{ΜΒ}$. \quad γ) Αν $\mathrm{ΜΑ}=\mathrm{ΜΒ}$ τότε το $\mathrm{Μ}$ είναι πάντα το μέσο του $\mathrm{ΑΒ}$. \quad δ) Η περίμετρος ενός σχήματος αλλάζει αν μετρήσω τις πλευρές με άλλη σειρά. \quad ε) $1$ m $=0{,}001$ km.

\ex Αν διπλασιάσουμε κάθε πλευρά ενός ορθογωνίου, πόσο γίνεται η περίμετρος; Θα ίσχυε το ίδιο και για ένα τρίγωνο; Δοκίμασε με παράδειγμα.

\zoni{Για το σπίτι}{μέτρησε και πρόσθεσε}

\ex Μέτρησε με τα βήματά σου και με μετροταινία το ίδιο μήκος στο σπίτι (π.χ. το διάδρομο) και υπολόγισε το μήκος του βήματός σου.

\ex Μέτρησε τις πλευρές του δωματίου σου και υπολόγισε την περίμετρο της κάτοψής του. Βρες και το μέσο ενός σχοινιού διπλώνοντάς το, και έλεγξε με μετροταινία.

\paragrafos{Μέτρο γωνίας}{Β.1.3, «Μέτρο γωνίας»}

\begin{anagke}
Ο παππούς Γιάννης, που ταξίδευε παλιά με καράβια, φέρνει στην τάξη μια πυξίδα και τη δείχνει στην παρέα. «Οι ναυτικοί βρίσκουν την πορεία τους με την πυξίδα», λέει. «Όταν το καράβι αλλάζει πορεία, ο καπετάνιος δεν λέει “στρίψε λίγο”. Λέει πόσο ακριβώς». Η κυρία Αλεξάνδρα ρωτάει: «Αν το πλοίο στρίψει προς τα βορειοανατολικά ή προς τα βορειοδυτικά, πώς ξέρουμε σε ποια περίπτωση η στροφή είναι μεγαλύτερη;».

Στο μεταξύ στον κήπο η Αγγελική σχεδιάζει μια γωνία φράχτη με πολύ μακριές πλευρές, και ο Νίκος μια ίδια γωνία με κοντές πλευρές. «Η δική μου είναι μεγαλύτερη, έχει μεγαλύτερες πλευρές!», λέει η Αγγελική. «Όχι», απαντά η Ρήνη. «Εγώ θα τις μετρήσω». Παίρνει ένα μοιρογνωμόνιο και μετράει και τις δύο: βγαίνει το ίδιο νούμερο.

Για να συγκρίνουμε, να προσθέτουμε και να αφαιρούμε γωνίες χρειαζόμαστε, όπως και στα μήκη, μια μονάδα και ένα όργανο μέτρησης. Σε αυτή την παράγραφο γνωρίζουμε τη μοίρα και το μοιρογνωμόνιο και μαθαίνουμε να υπολογίζουμε με μέτρα γωνιών.
\end{anagke}

\begin{lexi}{μοίρα, μοιρογνωμόνιο}{μοίρα, μοιρογνωμόνιο}
Η \lexiword{μοίρα} είναι το «μέρος», το μερίδιο που παίρνει ο καθένας. Μία μοίρα είναι ένα πολύ μικρό μέρος του πλήρους κύκλου. Το \lexiword{μοιρογνωμόνιο} είναι το όργανο που «γνωρίζει», δηλαδή μετράει, τις μοίρες.
\end{lexi}

\begin{metafrasi}
Μονάδα μέτρησης των γωνιών. & μοίρα, $1^\circ$\\
Η γωνία $\hat{\mathrm{Α}}$ ανοίγει $40$ μοίρες. & $\hat{\mathrm{Α}}=40^\circ$\\
Υποδιαίρεση της μοίρας. & $1^\circ=60'$ (πρώτα λεπτά)\\
Υποδιαίρεση του λεπτού. & $1'=60''$ (δεύτερα λεπτά)\\
Δύο γωνίες με το ίδιο μέτρο. & $\hat{\mathrm{Κ}}=\hat{\Lambda}$\\
\end{metafrasi}

\tmima{Η θεωρία}

\tmimasmall{Η μοίρα και το μοιρογνωμόνιο}

\begin{orismos}{μοίρα}
\textbf{Μονάδα μέτρησης των γωνιών} είναι η \textbf{μοίρα}, που γράφεται $1^\circ$.
\end{orismos}

Η μέτρηση των γωνιών γίνεται με το \textbf{μοιρογνωμόνιο}. Ο αριθμός που βρίσκουμε ονομάζεται \textbf{μέτρο} της γωνίας. Τοποθετούμε το κέντρο του μοιρογνωμονίου στην κορυφή της γωνίας και τη μία πλευρά πάνω στο μηδέν. Διαβάζουμε το μέτρο στο σημείο όπου η άλλη πλευρά διασχίζει την κλίμακα, και προσέχουμε να διαβάζουμε τη σωστή από τις δύο κλίμακες.

\begin{center}
\begin{tikzpicture}[scale=1.15]
\draw[ggray] (-2.2,0) arc(180:0:2.2) -- cycle;
\foreach \a in {0,10,...,180}{\draw[ggray] ({1.95*cos(\a)},{1.95*sin(\a)})--({2.2*cos(\a)},{2.2*sin(\a)});}
\foreach \a in {0,45,90,135,180}{\node[ggray,font=\scriptsize] at ({1.65*cos(\a)},{1.65*sin(\a)}) {$\a^\circ$};}
\draw[thick] (-2.6,0)--(2.6,0); \draw[thick,gblue] (0,0)--({2.6*cos(40)},{2.6*sin(40)});
\fill (0,0) circle (2pt) node[below] {$\mathrm{Α}$};
\draw[gblue,thick] (0.7,0) arc(0:40:0.7); \node[gblue] at (0.95,0.3) {\small $40^\circ$};
\end{tikzpicture}
\end{center}

Παραδείγματα: το μέτρο της γωνίας $\hat{\mathrm{Α}}$ είναι $40^\circ$ και το μέτρο της $\hat{\mathrm{Β}}$ είναι $125^\circ$. Μια γωνία που ανοίγει όσο η ορθή έχει μέτρο $90^\circ$. Μια γωνία με πλευρές αντικείμενες ημιευθείες έχει μέτρο $180^\circ$ και μια μη κυρτή γωνία έχει μέτρο μεγαλύτερο από $180^\circ$.

Η μοίρα έχει υποδιαίρεση, το \textbf{λεπτό}: $1^\circ=60'$ (πρώτα λεπτά) και $1'=60''$ (δεύτερα λεπτά). Για παράδειγμα, $25^\circ30'=25{,}5^\circ$, γιατί $30'$ είναι η μισή μοίρα. Και $48{,}25^\circ=48^\circ15'$, γιατί $0{,}25\cdot60=15$.

\begin{prosoxi}[Προσοχή: δεν είναι δεκαδικά]
Το $1^\circ=60'$ και όχι $100'$. Γι' αυτό το $25^\circ30'$ δεν είναι $25{,}30^\circ$ αλλά $25{,}5^\circ$.
\end{prosoxi}

\begin{kanonas}{το μέτρο εξαρτάται μόνο από το «άνοιγμα»}
Το μέτρο μιας γωνίας είναι μοναδικό και εξαρτάται μόνο από το \textbf{άνοιγμα} των πλευρών της, όχι από το μήκος τους. Οι πλευρές είναι ημιευθείες και δεν έχουν μήκος.
\end{kanonas}

\tmimasmall{Ίσες γωνίες και σύγκριση γωνιών}

Δύο γωνίες είναι \textbf{ίσες} αν έχουν το ίδιο μέτρο. Τότε η μία εφαρμόζει ακριβώς πάνω στην άλλη.

Μπορούμε να συγκρίνουμε δύο γωνίες και με διαφανές χαρτί, χωρίς μοιρογνωμόνιο:
\begin{itemize}[leftmargin=1.2em,itemsep=0.3ex,topsep=2pt]
\item Αντιγράφουμε τη μία από τις δύο γωνίες σε διαφανές χαρτί.
\item Τοποθετούμε το διαφανές χαρτί πάνω στην άλλη γωνία ώστε να έχουν κοινή κορυφή και μία κοινή πλευρά.
\end{itemize}
Η γωνία που μένει μέσα στην άλλη είναι η μικρότερη.

\begin{center}
\begin{tikzpicture}[scale=1.1]
\begin{scope}
\draw[thick] (3,0)--(0,0)--(2.6,1.5); \draw[gblue] (0.7,0) arc(0:30:0.7); \fill (0,0) circle (2pt);
\node[below left] at (0,0) {$\mathrm{Κ}$}; \node[right] at (1.6,0.3) {\small $30^\circ$};
\end{scope}
\begin{scope}[xshift=4.5cm]
\draw[thick] (1.5,0)--(0,0)--(1.3,0.75); \draw[gblue] (0.7,0) arc(0:30:0.7); \fill (0,0) circle (2pt);
\node[below left] at (0,0) {$\Lambda$};
\end{scope}
\begin{scope}[xshift=8cm]
\draw[thick] (3,0)--(0,0)--(2.2,2); \fill[gblueL] (0,0)--(2.2,0) arc(0:42:2.2)--cycle;
\fill[gwarm] (0,0)--(1.7,0) arc(0:28:1.7)--cycle; \fill (0,0) circle (2pt);
\node[below left] at (0,0) {$\mathrm{Ο}$};
\end{scope}
\end{tikzpicture}
\end{center}

Στα δύο πρώτα σχήματα οι γωνίες έχουν το ίδιο μέτρο ($30^\circ$) παρά τα διαφορετικά μήκη των πλευρών τους, άρα είναι ίσες. Στο τρίτο σχήμα η μικρότερη γωνία (στο εσωτερικό της άλλης) έχει μικρότερο μέτρο.

\tmimasmall{Πρόσθεση και αφαίρεση γωνιών}

Δύο γωνίες μπορούν να προστεθούν με \textbf{επίθεση}: τις τοποθετούμε δίπλα δίπλα, με κοινή κορυφή και μία κοινή πλευρά, χωρίς να επικαλύπτονται. Η γωνία που σχηματίζεται από τις δύο εξωτερικές πλευρές έχει μέτρο το \textbf{άθροισμα} των μέτρων τους. Αντίστροφα, αν από μια γωνία αφαιρέσουμε ένα μέρος της, το υπόλοιπο έχει μέτρο τη \textbf{διαφορά}.

\begin{center}
\begin{tikzpicture}[scale=1.1]
\draw[thick] (3,0)--(0,0); \draw[thick] (0,0)--(2.4,1.7); \draw[thick] (0,0)--(-0.9,2.6);
\fill (0,0) circle (2pt) node[below left] {$\mathrm{Ο}$};
\node[right] at (3,0) {$x$}; \node[right] at (2.4,1.7) {$y$}; \node[above] at (-0.9,2.6) {$z$};
\draw[gblue] (0.8,0) arc(0:35:0.8); \node[gblue] at (1.2,0.18) {\small $\alpha$};
\draw[gwarmD] (0.55,0.38) arc(35:110:0.55) ; \node[gwarmD] at (0.12,0.85) {\small $\beta$};
\end{tikzpicture}
\end{center}

Αν οι γωνίες $x\hat{\mathrm{Ο}}y$ και $y\hat{\mathrm{Ο}}z$ τοποθετηθούν δίπλα δίπλα με αυτό τον τρόπο, τότε $x\hat{\mathrm{Ο}}z=x\hat{\mathrm{Ο}}y+y\hat{\mathrm{Ο}}z$ και $y\hat{\mathrm{Ο}}z=x\hat{\mathrm{Ο}}z-x\hat{\mathrm{Ο}}y$.

\begin{prosoxi}
\begin{itemize}[leftmargin=1.2em,itemsep=0.3ex,topsep=2pt]
\item Μια γωνία με μακρύτερες πλευρές δεν είναι μεγαλύτερη: το μέτρο εξαρτάται μόνο από το άνοιγμα.
\item Στο μοιρογνωμόνιο διαβάζω την κλίμακα που ξεκινά από το μηδέν στην πλευρά που έβαλα πάνω στο $0$. Αν η γωνία δείχνει μικρότερη από την ορθή, δεν γράφω αριθμό μεγαλύτερο από $90$.
\item Για να προσθέσω δύο γωνίες πρέπει να έχουν κοινή κορυφή και κοινή πλευρά και να μην επικαλύπτονται.
\item Το $25^\circ30'$ δεν ισούται με $25{,}30^\circ$.
\end{itemize}
\end{prosoxi}

\tmima{Λυμένα παραδείγματα}

\paradeigma{Η σωστή κλίμακα}
\tixero Το μοιρογνωμόνιο έχει δύο κλίμακες. Η μία πλευρά μιας γωνίας δείχνει στο $0$ και η άλλη δείχνει $40$ στη μία κλίμακα και $140$ στην άλλη. Η γωνία φαίνεται μικρότερη από την ορθή.\par
\tithelo Το μέτρο της γωνίας.\par
\possindeo Αφού η γωνία είναι μικρότερη από την ορθή, το μέτρο της είναι μικρότερο από $90^\circ$.\par
\apantisi $40^\circ$.\par
\elegxos $40^\circ<90^\circ$ και $40^\circ+140^\circ=180^\circ$, όπως οι δύο κλίμακες.

\paradeigma{Μοίρες και λεπτά}
\tixero Η γωνία $\hat{\mathrm{Α}}=36^\circ30'$ και η γωνία $\hat{\mathrm{Β}}=48{,}25^\circ$.\par
\tithelo Να γράψω την $\hat{\mathrm{Α}}$ με δεκαδικό και την $\hat{\mathrm{Β}}$ με λεπτά.\par
\possindeo Επειδή $1^\circ=60'$, τα $30'$ είναι $30:60=0{,}5^\circ$. Και $0{,}25^\circ=0{,}25\cdot60'=15'$.\par
\apantisi $\hat{\mathrm{Α}}=36{,}5^\circ$ και $\hat{\mathrm{Β}}=48^\circ15'$.\par
\elegxos $0{,}5^\circ$ είναι η μισή μοίρα, δηλαδή $30'$.

\paradeigma{Ίσες ή όχι;}
\tixero Η Αγγελική έχει γωνία με πλευρές $12$ cm και ο Νίκος γωνία με πλευρές $4$ cm. Και οι δύο μετρούν $\hat{\omega}=62^\circ$.\par
\tithelo Ποια γωνία είναι μεγαλύτερη;\par
\possindeo Το μέτρο μιας γωνίας εξαρτάται από το άνοιγμα και όχι από τα μήκη των πλευρών.\par
\apantisi Είναι ίσες.\par
\elegxos Αν αντιγράψω τη μία σε διαφανές χαρτί, εφαρμόζει πάνω στην άλλη.

\paradeigma{Πρόσθεση και αφαίρεση}
\tixero Η γωνία του φράχτη είναι $x\hat{\mathrm{Ο}}z=84^\circ$ και ένα μονοπάτι $\mathrm{Οy}$ την χωρίζει ώστε $x\hat{\mathrm{Ο}}y=50^\circ$.\par
\tithelo Η γωνία $y\hat{\mathrm{Ο}}z$.\par
\possindeo Η $x\hat{\mathrm{Ο}}z$ είναι το άθροισμα των $x\hat{\mathrm{Ο}}y$ και $y\hat{\mathrm{Ο}}z$. Άρα $y\hat{\mathrm{Ο}}z=84^\circ-50^\circ$.\par
\apantisi $34^\circ$.\par
\elegxos $50^\circ+34^\circ=84^\circ$.

\paradeigma{Οι δείκτες του ρολογιού}
\tixero Ο κύκλος γύρω από το κέντρο του ρολογιού έχει $360^\circ$ και το ρολόι έχει $12$ ώρες.\par
\tithelo Η γωνία ανάμεσα σε δύο διαδοχικές ώρες και η γωνία των δεικτών στις $4{:}00$.\par
\possindeo Ανάμεσα σε δύο διαδοχικές ώρες: $360^\circ:12=30^\circ$. Στις $4{:}00$ ο λεπτοδείκτης δείχνει το $12$ και ο ωροδείκτης το $4$, δηλαδή $4$ διαστήματα ωρών.\par
\apantisi $30^\circ$ και $4\cdot30^\circ=120^\circ$.\par
\elegxos Στις $3{:}00$ η γωνία είναι $3\cdot30^\circ=90^\circ$, δηλαδή ορθή.

\tmima{Ασκήσεις}

\zoni{Το κατάλαβα}{μέτρο, σύγκριση, λεπτά, πρόσθεση}

\ex Σχεδίασε με το μοιρογνωμόνιο τις γωνίες: $\hat{\alpha}=30^\circ$, $\hat{\beta}=58^\circ$, $\hat{\gamma}=90^\circ$, $\hat{\delta}=135^\circ$, $\hat{\varepsilon}=180^\circ$.

\ex Σχεδίασε δύο γωνίες με μέτρο $40^\circ$: τη μία με πλευρές μήκους $3$ cm και την άλλη με πλευρές μήκους $8$ cm. Ποια είναι μεγαλύτερη; Εξήγησε.

\ex Μετάτρεψε: α) $12^\circ30'$ σε μοίρες με δεκαδικό \quad β) $40^\circ15'$ σε μοίρες με δεκαδικό \quad γ) $18{,}75^\circ$ σε μοίρες και λεπτά \quad δ) $1^\circ$ σε δεύτερα λεπτά.

\ex Βάλε σε αύξουσα σειρά τις γωνίες: $45^\circ$, $40^\circ30'$, $40{,}25^\circ$, $45^\circ10'$.

\ex Οι γωνίες $x\hat{\mathrm{Ο}}y$ και $y\hat{\mathrm{Ο}}z$ έχουν κοινή κορυφή, κοινή πλευρά $\mathrm{Οy}$ και δεν επικαλύπτονται. α) Αν $x\hat{\mathrm{Ο}}y=35^\circ$ και $y\hat{\mathrm{Ο}}z=48^\circ$, βρες την $x\hat{\mathrm{Ο}}z$. β) Αν $x\hat{\mathrm{Ο}}z=120^\circ$ και $x\hat{\mathrm{Ο}}y=75^\circ$, βρες την $y\hat{\mathrm{Ο}}z$.

\ex Σχεδίασε ένα τρίγωνο, μέτρησε τις τρεις γωνίες του με το μοιρογνωμόνιο και βρες το άθροισμά τους. Επανάλαβε με άλλο τρίγωνο. Τι παρατηρείς;

\ex Σωστό ή λάθος; α) Μονάδα μέτρησης των γωνιών είναι η μοίρα. \quad β) Όσο μεγαλύτερο μήκος έχουν οι πλευρές μιας γωνίας, τόσο μεγαλύτερη είναι. \quad γ) Δύο γωνίες είναι ίσες όταν έχουν το ίδιο μέτρο. \quad δ) $1^\circ=100'$. \quad ε) Το μέτρο της ορθής γωνίας είναι $90^\circ$. \quad στ) Το μοιρογνωμόνιο μετράει γωνίες.

\zoni{Το χρησιμοποιώ}{πυξίδα, ρολόι, φράχτης, πόρτα}

\ex \textbf{Το πλοίο.} Ένα πλοίο στρίβει πρώτα $45^\circ$ και μετά ακόμη $30^\circ$ προς την ίδια μεριά. Πόση είναι η συνολική στροφή; Αν μετά στρίψει $20^\circ$ προς την αντίθετη μεριά, πόση στροφή απομένει από την αρχική πορεία;

\ex \textbf{Το ρολόι.} Πόσες μοίρες είναι η γωνία ανάμεσα στους δείκτες στις $2{:}00$, στις $5{:}00$ και στις $6{:}00$;

\ex \textbf{Το παρτέρι.} Η γωνία του φράχτη στον κήπο είναι $84^\circ$. Η Αγγελική σχεδιάζει ένα μονοπάτι που σχηματίζει $50^\circ$ με τη μία πλευρά και χωρίζει το παρτέρι σε δύο μέρη. Πόση γωνία σχηματίζει το μονοπάτι με την άλλη πλευρά;

\ex \textbf{Η πόρτα.} Η πόρτα της τάξης ανοίγει πρώτα $35^\circ$ και μετά ανοίγει ακόμη $50^\circ$. Στη συνέχεια η Φανή τη σπρώχνει και κλείνει $20^\circ$. Πόσο ανοιχτή είναι τώρα;

\ex \textbf{Η κλίμακα.} Η Φανή μετράει μια γωνία που φαίνεται μικρότερη από την ορθή και το μοιρογνωμόνιο δείχνει $130$ στη μία κλίμακα και $50$ στην άλλη. Ποιο είναι το σωστό μέτρο και γιατί;

\zoni{Το σκέφτομαι}{τι εξαρτάται από τι}

\ex Γιατί το μέτρο μιας γωνίας δεν εξαρτάται από το μήκος των πλευρών της; Σκέψου τι είναι οι πλευρές μιας γωνίας.

\ex Οι γωνίες $x\hat{\mathrm{Ο}}y=70^\circ$, $y\hat{\mathrm{Ο}}z=50^\circ$ και $z\hat{\mathrm{Ο}}t=40^\circ$ είναι διαδοχικές, χωρίς να επικαλύπτονται. Πόση είναι η $x\hat{\mathrm{Ο}}t$;

\ex Πώς μπορείς να συγκρίνεις δύο γωνίες όταν δεν έχεις μοιρογνωμόνιο; Περιέγραψε τα βήματα.

\ex Σχεδίασε μια γωνία $x\hat{\mathrm{Ο}}y=30^\circ$. Με κοινή κορυφή και κοινή πλευρά $\mathrm{Οx}$ σχεδίασε γωνία $x\hat{\mathrm{Ο}}t$ α) διπλάσιου μέτρου \quad β) μισού μέτρου.

\zoni{Για το σπίτι}{γωνίες γύρω μας}

\ex Βρες στο σπίτι πέντε γωνίες (ανοιχτό βιβλίο, ψαλίδι, ανοιχτή πόρτα, γωνία τραπεζιού, ανοιχτός φορητός υπολογιστής) και μέτρησε το άνοιγμά τους με το μοιρογνωμόνιο.

\ex Άνοιξε έναν διαβήτη σε τρία διαφορετικά ανοίγματα και σχεδίασε στο τετράδιό σου τη γωνία που σχηματίζουν τα σκέλη του κάθε φορά. Μέτρησέ τες.

\paragrafos{Άθροισμα και διαφορά τμημάτων: από την ευθεία στην τεθλασμένη}{Β.1.4, «Πρόσθεση και αφαίρεση ευθυγράμμων τμημάτων»}

\begin{anagke}
Ο παππούς Γιάννης φέρνει τα δέντρα που θα φυτευτούν σε μία σειρά, και ο Ηλίας μετράει τις αποστάσεις: το δεύτερο δέντρο απέχει από το πρώτο $3$ μέτρα, το τρίτο $4$ μέτρα από το δεύτερο και το τέταρτο $2{,}5$ από το τρίτο. «Πόσο απέχει το τέταρτο από το πρώτο;» ρωτάει η Πολυξένη. Ο Ηλίας το βρίσκει εύκολα: $3+4+2{,}5=9{,}5$ μέτρα. «Και το τρίτο από το πρώτο; $3+4=7$. Και το τέταρτο από το δεύτερο;». Εδώ ο Στέλιος προσπαθεί να προσθέσει, αλλά η Αγγελική λέει: «Αφαίρεσε: από την απόσταση του τέταρτου από το πρώτο, αφαιρώ την απόσταση του δεύτερου από το πρώτο: $9{,}5-3=6{,}5$».

Η παρέα ανακαλύπτει ότι οι αποστάσεις προσθέτονται και αφαιρούνται όταν τα σημεία βρίσκονται στην ίδια ευθεία με τη σειρά. Μετά ο Νίκος γυρίζει στο μονοπάτι που στρίβει: «Αν πάω από το πρώτο στο τρίτο δέντρο περνώντας από ένα άλλο σημείο που δεν είναι στην ευθεία, κάνω περισσότερο δρόμο;». Η Ρήνη το δοκιμάζει με σχοινί και δείχνει ότι το ευθύγραμμο τμήμα είναι πάντα το συντομότερο μονοπάτι ανάμεσα σε δύο σημεία. Και η κυρία Αλεξάνδρα προσθέτει: «Το μήκος του περιγράμματος ενός παρτεριού, για το φράχτη του, είναι η περίμετρός του».
\end{anagke}

\begin{lexi}{άθροισμα τμημάτων, περίμετρος}{άθροισμα τμημάτων, περίμετρος}
Η \lexiword{περίμετρος} είναι το «γύρω μέτρημα» (περί + μέτρο): το συνολικό μήκος του περιγράμματος ενός σχήματος. Η \lexiword{τεθλασμένη} γραμμή έχει μήκος το άθροισμα των μηκών των ευθυγράμμων τμημάτων της.
\end{lexi}

\begin{metafrasi}
Τα $\mathrm{Α},\,\mathrm{Β},\,\Gamma$ στη σειρά σε ευθεία. & $\mathrm{ΑΓ}=\mathrm{ΑΒ}+\mathrm{ΒΓ}$\\
Ό,τι περισσεύει μετά το $\mathrm{ΑΒ}$. & $\mathrm{ΒΓ}=\mathrm{ΑΓ}-\mathrm{ΑΒ}$\\
Τεθλασμένη $\mathrm{ΑΒΓΔ}$. & μήκος $=\mathrm{ΑΒ}+\mathrm{ΒΓ}+\Gamma\Delta$\\
Ο συντομότερος δρόμος. & $\mathrm{ΑΒ}<$ μήκος κάθε τεθλασμένης με άκρα $\mathrm{Α},\,\mathrm{Β}$\\
Το άθροισμα των πλευρών. & περίμετρος\\
\end{metafrasi}

\tmima{Η θεωρία}

Αν τα σημεία $\mathrm{Α},\,\mathrm{Β},\,\Gamma$ βρίσκονται με αυτή τη σειρά σε μια ευθεία, το τμήμα $\mathrm{ΑΓ}$ λέγεται \textbf{άθροισμα} των τμημάτων $\mathrm{ΑΒ}$ και $\mathrm{ΒΓ}$:
$$\mathrm{ΑΓ}=\mathrm{ΑΒ}+\mathrm{ΒΓ}.$$
Και το $\mathrm{ΒΓ}$ είναι η \textbf{διαφορά} των $\mathrm{ΑΓ}$ και $\mathrm{ΑΒ}$: $\mathrm{ΒΓ}=\mathrm{ΑΓ}-\mathrm{ΑΒ}$.

\begin{center}
\begin{tikzpicture}
\draw[thick] (0,0)--(7,0);
\foreach \x/\l in {0/Α,2.5/Β,7/Γ}{\fill (\x,0) circle (2pt); \node[below] at (\x,-0.1) {$\mathrm{\l}$};}
\draw[gblue,thick] (0,0.4)--(2.5,0.4); \draw[gblue,thick,dashed] (2.5,0.4)--(7,0.4);
\node[gblue,above] at (1.25,0.45) {$\mathrm{ΑΒ}$}; \node[gblue,above] at (4.75,0.45) {$\mathrm{ΒΓ}$};
\end{tikzpicture}
\end{center}

\tmimasmall{Τεθλασμένη και περίμετρος}

Το μήκος μιας \textbf{τεθλασμένης γραμμής} είναι το άθροισμα των μηκών των ευθυγράμμων τμημάτων της. Το μήκος του ευθυγράμμου τμήματος $\mathrm{ΑΒ}$ είναι μικρότερο από το μήκος κάθε τεθλασμένης γραμμής με τα ίδια άκρα $\mathrm{Α}$, $\mathrm{Β}$ (εκτός αν η τεθλασμένη είναι το ίδιο το τμήμα). Δηλαδή το ευθύγραμμο τμήμα είναι ο συντομότερος δρόμος μεταξύ δύο σημείων.

\begin{center}
\begin{tikzpicture}
\draw[thick] (0,0)--(6,0); \draw[thick,gblue] (0,0)--(2,1.6)--(4,1)--(6,0);
\foreach \p/\l in {(0,0)/Α,(6,0)/Β,(2,1.6)/Γ,(4,1)/Δ}{\fill \p circle (2pt);}
\node[below left] at (0,0) {$\mathrm{Α}$}; \node[below right] at (6,0) {$\mathrm{Β}$}; \node[above] at (2,1.6) {$\Gamma$}; \node[above right] at (4,1) {$\Delta$};
\end{tikzpicture}
\end{center}

Το άθροισμα των πλευρών ενός ευθύγραμμου σχήματος (πολυγώνου) λέγεται \textbf{περίμετρος} του σχήματος.

\begin{prosoxi}
\begin{itemize}[leftmargin=1.2em,itemsep=0.3ex,topsep=2pt]
\item Η σχέση $\mathrm{ΑΓ}=\mathrm{ΑΒ}+\mathrm{ΒΓ}$ ισχύει μόνο όταν το $\mathrm{Β}$ βρίσκεται \emph{ανάμεσα} στα $\mathrm{Α}$ και $\mathrm{Γ}$ πάνω σε ευθεία. Αν δεν είναι συνευθειακά, το $\mathrm{ΑΓ}$ είναι μικρότερο από το $\mathrm{ΑΒ}+\mathrm{ΒΓ}$.
\item Η περίμετρος μετράει μόνο το μήκος του περιγράμματος, όχι την επιφάνεια.
\item Για την αφαίρεση προσέχω τη σειρά: $\mathrm{ΒΓ}=\mathrm{ΑΓ}-\mathrm{ΑΒ}$ και όχι $\mathrm{ΑΒ}-\mathrm{ΑΓ}$ (αρνητικό).
\item Για να συγκρίνω δύο διαδρομές πρέπει να έχουν την ίδια μονάδα μέτρησης.
\end{itemize}
\end{prosoxi}

\tmima{Λυμένα παραδείγματα}

\paradeigma{Τα δέντρα στη σειρά}
\tixero Τα δέντρα $\mathrm{Α},\,\mathrm{Β},\,\Gamma,\,\Delta$ είναι στη σειρά και $\mathrm{ΑΒ}=3$, $\mathrm{ΒΓ}=4$, $\Gamma\Delta=2{,}5$ μ.\par
\tithelo Τα $\mathrm{ΑΓ}$, $\mathrm{ΑΔ}$ και $\mathrm{ΒΔ}$.\par
\possindeo $\mathrm{ΑΓ}=\mathrm{ΑΒ}+\mathrm{ΒΓ}=3+4=7$. $\mathrm{ΑΔ}=\mathrm{ΑΓ}+\Gamma\Delta=7+2{,}5=9{,}5$. $\mathrm{ΒΔ}=\mathrm{ΑΔ}-\mathrm{ΑΒ}=9{,}5-3=6{,}5$.\par
\apantisi $7$, $9{,}5$ και $6{,}5$ μ.\par
\elegxos $\mathrm{ΒΔ}=\mathrm{ΒΓ}+\Gamma\Delta=4+2{,}5=6{,}5$.

\paradeigma{Με αφαίρεση}
\tixero Το $\mathrm{Β}$ είναι ανάμεσα στα $\mathrm{Α}$ και $\mathrm{Γ}$, $\mathrm{ΑΓ}=9$ και $\mathrm{ΑΒ}=4{,}5$.\par
\tithelo Το $\mathrm{ΒΓ}$.\par
\possindeo $\mathrm{ΒΓ}=\mathrm{ΑΓ}-\mathrm{ΑΒ}=9-4{,}5=4{,}5$.\par
\apantisi $4{,}5$.\par
\elegxos Άρα το $\mathrm{Β}$ είναι το μέσο του $\mathrm{ΑΓ}$.

\paradeigma{Ο κοντινός δρόμος}
\tixero Μονοπάτι $\mathrm{ΑΒΓ}$ με $\mathrm{ΑΒ}=6$ μ. και $\mathrm{ΒΓ}=8$ μ. (στρίβει στο $\mathrm{Β}$).\par
\tithelo Πόσο είναι το μήκος του και τι ξέρουμε για το $\mathrm{ΑΓ}$;\par
\possindeo Το μήκος της τεθλασμένης είναι $6+8=14$ μ. Το ευθύγραμμο τμήμα $\mathrm{ΑΓ}$ είναι συντομότερο.\par
\apantisi Τεθλασμένη $14$ μ. και $\mathrm{ΑΓ}<14$ μ.\par
\elegxos Αν το $\mathrm{Β}$ ήταν πάνω στο $\mathrm{ΑΓ}$, θα ήταν ίσα με $14$.

\paradeigma{Η περίμετρος}
\tixero Ένα ορθογώνιο παρτέρι με πλευρές $4{,}5$ μ. και $3{,}2$ μ.\par
\tithelo Το μήκος του φράχτη που το περιβάλλει.\par
\possindeo Η περίμετρος είναι το άθροισμα των πλευρών: $4{,}5+3{,}2+4{,}5+3{,}2=15{,}4$ μ.\par
\apantisi $15{,}4$ μ.\par
\elegxos $2\cdot(4{,}5+3{,}2)=2\cdot7{,}7=15{,}4$.

\paradeigma{Το μέσο και το άθροισμα}
\tixero Το $\mathrm{Γ}$ βρίσκεται πάνω στο $\mathrm{ΑΒ}$ με $\mathrm{ΑΓ}=3$ και $\mathrm{ΓΒ}=7$ και το $\mathrm{Μ}$ είναι το μέσο του $\mathrm{ΑΒ}$.\par
\tithelo Το $\mathrm{ΓΜ}$.\par
\possindeo $\mathrm{ΑΒ}=3+7=10$, άρα $\mathrm{ΑΜ}=5$. Το $\mathrm{Μ}$ είναι πέρα από το $\mathrm{Γ}$ και $\mathrm{ΓΜ}=\mathrm{ΑΜ}-\mathrm{ΑΓ}=5-3=2$.\par
\apantisi $\mathrm{ΓΜ}=2$.\par
\elegxos $\mathrm{ΓΜ}+\mathrm{ΜΒ}=2+5=7=\mathrm{ΓΒ}$.

\tmima{Ασκήσεις}

\zoni{Το κατάλαβα}{άθροισμα, διαφορά, περίμετρος}

\ex Τα $\mathrm{Α},\,\mathrm{Β},\,\Gamma,\,\Delta$ είναι στη σειρά σε μια ευθεία με $\mathrm{ΑΒ}=3$, $\mathrm{ΒΓ}=4$, $\Gamma\Delta=2{,}5$. Βρες τα $\mathrm{ΑΓ}$, $\mathrm{ΑΔ}$, $\mathrm{ΒΔ}$.

\ex Το $\mathrm{Β}$ είναι ανάμεσα στα $\mathrm{Α}$ και $\mathrm{Γ}$. Αν $\mathrm{ΑΓ}=9$ και $\mathrm{ΑΒ}=4{,}5$, βρες το $\mathrm{ΒΓ}$.

\ex Βρες το μήκος της τεθλασμένης $\mathrm{ΑΒΓΔ}$ με $\mathrm{ΑΒ}=5$, $\mathrm{ΒΓ}=3{,}5$, $\Gamma\Delta=6$.

\ex Βρες την περίμετρο: α) ορθογωνίου $4{,}5$ και $3{,}2$ \quad β) τριγώνου με πλευρές $5$, $6$, $7{,}5$ \quad γ) τετραγώνου πλευράς $2{,}8$ \quad δ) κανονικού εξαγώνου πλευράς $4$.

\ex Στην ευθεία $\mathrm{Α},\,\mathrm{Β},\,\Gamma$ με $\mathrm{ΑΒ}=2\mathrm{ΒΓ}$ και $\mathrm{ΑΓ}=12$ βρες τα $\mathrm{ΑΒ}$ και $\mathrm{ΒΓ}$.

\ex Σωστό ή λάθος; α) Αν το $\mathrm{Β}$ είναι ανάμεσα στα $\mathrm{Α}$ και $\Gamma$ τότε $\mathrm{ΑΓ}=\mathrm{ΑΒ}+\mathrm{ΒΓ}$. \quad β) Η τεθλασμένη μεταξύ δύο σημείων είναι πάντα μεγαλύτερη από το τμήμα που τα ενώνει (ή ίση). \quad γ) Η περίμετρος ενός τετραγώνου πλευράς $\alpha$ είναι $\alpha^2$. \quad δ) $\mathrm{ΒΓ}=\mathrm{ΑΒ}-\mathrm{ΑΓ}$ όταν το $\mathrm{Β}$ είναι ανάμεσα στα $\mathrm{Α}$, $\mathrm{Γ}$.

\zoni{Το χρησιμοποιώ}{φράχτες, δρόμοι, κήποι}

\ex \textbf{Ο φράχτης.} Ένα ορθογώνιο παρτέρι $4{,}5$ μ. επί $3{,}2$ μ. θα περιφραχθεί με δίχτυ που κοστίζει $3$ € το μέτρο. Πόσο θα κοστίσει; Μην ξεχάσεις ότι ένα μέτρο το χρειάζεται για την είσοδο που μένει ανοιχτή.

\ex \textbf{Ο κοντινός δρόμος.} Το μονοπάτι $\mathrm{ΑΒΓ}$ έχει $\mathrm{ΑΒ}=6$ μ., $\mathrm{ΒΓ}=8$ μ. Ο Νίκος κόβει δρόμο πηγαίνοντας ευθεία από το $\mathrm{Α}$ στο $\mathrm{Γ}$. Έκανε περισσότερο ή λιγότερο δρόμο από $14$ μ.; Γιατί;

\ex \textbf{Η απόσταση.} Τρία σπίτια $\mathrm{Α}$, $\mathrm{Β}$, $\Gamma$ βρίσκονται στον ίδιο δρόμο, με το $\mathrm{Β}$ στη μέση. Το $\mathrm{Α}$ απέχει από το $\mathrm{Β}$ $350$ μ. και το $\mathrm{Β}$ από το $\Gamma$ $480$ μ. Πόσο απέχουν τα $\mathrm{Α}$ και $\Gamma$;

\ex \textbf{Η τεθλασμένη διαδρομή.} Η Πολυξένη πηγαίνει από το σχολείο στο σπίτι με στάσεις: $250$ μ., $180$ μ. και $320$ μ. Αν πήγαινε ευθεία θα έκανε $600$ μ. Πόσο περισσότερο περπάτησε;

\ex \textbf{Ο φράχτης του κήπου.} Ο κήπος έχει σχήμα πενταγώνου με πλευρές $12$, $9$, $14$, $10$ και $8$ μ. Πόσος φράχτης χρειάζεται;

\zoni{Το σκέφτομαι}{ο ευθύς είναι ο συντομότερος}

\ex Εξήγησε γιατί το ευθύγραμμο τμήμα είναι ο συντομότερος δρόμος ανάμεσα σε δύο σημεία (σκέψου ένα τεντωμένο νήμα).

\ex Αν $\mathrm{ΑΒ}=5$ και $\mathrm{ΒΓ}=7$ και τα $\mathrm{Α},\,\mathrm{Β},\,\Gamma$ δεν είναι συνευθειακά, ανάμεσα σε ποιες τιμές βρίσκεται το $\mathrm{ΑΓ}$;

\ex Ένα σχοινί $\mathrm{ΑΒ}$ μήκους $10$ μ. διπλώνεται ώστε το $\mathrm{Α}$ να πέσει πάνω στο $\mathrm{Β}$. Πού βρίσκεται η πτυχή;

\ex Δείξε με παράδειγμα ότι η περίμετρος δύο σχημάτων μπορεί να είναι ίση χωρίς τα σχήματα να είναι ίσα.

\zoni{Για το σπίτι}{μέτρησε και πρόσθεσε}

\ex Μέτρησε τις πλευρές του δωματίου σου και υπολόγισε την περίμετρο της κάτοψής του.

\ex Φτιάξε μια τεθλασμένη με τρία τμήματα από σπάγκο και σύγκρινε το μήκος της με την απόσταση των άκρων.

\paragrafos{Μέτρηση γωνιών, σύγκριση και διχοτόμος}{Β.1.5, «Μέτρηση, σύγκριση και ισότητα γωνιών – Διχοτόμος γωνίας»}

\begin{anagke}
Στον κήπο υπάρχει μια γωνία φράχτη, και η Α3 θέλει να τη χωρίσει σε δύο ίσα παρτέρια με ένα μονοπάτι. «Πού βάζουμε το μονοπάτι;» ρωτάει ο Βασίλης. Η Αγγελική παίρνει το μοιρογνωμόνιο: «Πρώτα πρέπει να μετρήσουμε πόσο ανοίγει η γωνία». Βρίσκει $84$ μοίρες. «Το μισό είναι $42$», υπολογίζει η Πολυξένη. «Το μονοπάτι θα ξεκινήσει από την κορυφή και θα σχηματίζει $42$ μοίρες με κάθε πλευρά».

Ο Ηλίας, που φοράει ρολόι με δείκτες, παρατηρεί ότι οι δείκτες σχηματίζουν γωνία: «Στις $3{:}00$ οι δείκτες είναι κάθετοι, στις $6{:}00$ ευθεία». Η Ρήνη ρωτάει: «Μπορούμε να ξέρουμε πόσες μοίρες είναι η γωνία των δεικτών κάθε ώρα;». Ο Στέλιος δουλεύει με τον κύκλο: «Ολόκληρος ο κύκλος έχει $360$ μοίρες. Το ρολόι έχει $12$ ώρες. Κάθε ώρα είναι $360:12=30$ μοίρες».

Σε αυτή την παράγραφο μαθαίνουμε να μετράμε γωνίες με μοιρογνωμόνιο, να τις συγκρίνουμε και να τις χωρίζουμε σε δύο ίσα μέρη με τη διχοτόμο.
\end{anagke}

\begin{lexi}{μοίρα, μοιρογνωμόνιο, διχοτόμος}{μοίρα, μοιρογνωμόνιο, διχοτόμος}
Η \lexiword{μοίρα} είναι το «μέρος» (από το «μοίρα», μερίδιο): η μοίρα ($1^\circ$) είναι το $\frac{1}{360}$ του πλήρους κύκλου. Το \lexiword{μοιρογνωμόνιο} «γνωρίζει» (μετράει) τις μοίρες. Η \lexiword{διχοτόμος} «χωρίζει στα δύο» (δίχα + τέμνω) μια γωνία σε δύο ίσες γωνίες.
\end{lexi}

\begin{metafrasi}
Το $\frac{1}{360}$ του κύκλου. & $1^\circ$\\
Η υποδιαίρεση της μοίρας. & $1^\circ=60'$ και $1'=60''$\\
Η γωνία ανοίγει $84$ μοίρες. & $x\mathrm{Ο}y=84^\circ$\\
Η $\mathrm{Οδ}$ είναι διχοτόμος. & $x\mathrm{Ο}\delta=\delta\mathrm{Ο}y=\dfrac{x\mathrm{Ο}y}{2}$\\
Κάθε ώρα στο ρολόι. & $360^\circ:12=30^\circ$\\
\end{metafrasi}

\tmima{Η θεωρία}

\tmimasmall{Μέτρηση γωνιών}

Για να \textbf{μετρήσουμε} μια γωνία τη συγκρίνουμε με μια μονάδα, τη \textbf{μοίρα}. Μία μοίρα ($1^\circ$) είναι το $\frac{1}{360}$ του πλήρους κύκλου. Υποδιαιρείται σε \textbf{πρώτα λεπτά} και \textbf{δεύτερα λεπτά}: $1^\circ=60'$ και $1'=60''$. Για παράδειγμα $25^\circ30'=25{,}5^\circ$ και $48{,}25^\circ=48^\circ15'$, γιατί $0{,}25\cdot60=15$. Μετράμε με το \textbf{μοιρογνωμόνιο}: τοποθετούμε το κέντρο του στην κορυφή και τη μία πλευρά στο μηδέν, και διαβάζουμε πού περνά η άλλη πλευρά. Προσέχουμε σε ποια από τις δύο κλίμακες ξεκινά το μηδέν.

\tmimasmall{Σύγκριση και ισότητα}

Δύο γωνίες είναι \textbf{ίσες} όταν εφαρμόζουν η μία στην άλλη (έχουν το ίδιο μέτρο). Η μία είναι \textbf{μεγαλύτερη} από την άλλη όταν το μέτρο της είναι μεγαλύτερο, ή όταν η μικρή χωράει στο εσωτερικό της μεγάλης με κοινή κορυφή και μία κοινή πλευρά. Το μέτρο μιας γωνίας δεν εξαρτάται από το μήκος των πλευρών της (οι πλευρές είναι ημιευθείες).

\tmimasmall{Διχοτόμος}

\begin{orismos}{διχοτόμος γωνίας}
\textbf{Διχοτόμος} της γωνίας $x\mathrm{Ο}y$ είναι η ημιευθεία $\mathrm{Οδ}$ που βρίσκεται στο εσωτερικό της και τη χωρίζει σε δύο ίσες γωνίες: $x\mathrm{Ο}\delta=\delta\mathrm{Ο}y$. Τότε κάθε μία είναι το μισό της $x\mathrm{Ο}y$.
\end{orismos}

\begin{center}
\begin{tikzpicture}[scale=1.1]
\draw[thick] (3.5,0)--(0,0)--(2.6,2.2); \draw[gblue,thick] (0,0)--(3.2,1.1); \fill (0,0) circle (2pt);
\node[right] at (3.5,0) {$x$}; \node[above] at (2.6,2.2) {$y$}; \node[right,gblue] at (3.2,1.1) {$\delta$}; \node[below left] at (0,0) {$\mathrm{Ο}$};
\draw[gblue] (0.9,0) arc(0:19:0.9); \draw[gblue] (1.1,0.32) arc(19:40:1.1);
\draw (0.6,0.3) -- (0.6,0.3);
\end{tikzpicture}
\end{center}

Κάθε γωνία έχει μία μοναδική διχοτόμο. Την κατασκευάζουμε με μοιρογνωμόνιο (μετράμε, διαιρούμε με το $2$ και σημειώνουμε), ή με διαφανές χαρτί διπλώνοντας ώστε οι δύο πλευρές να συμπέσουν: η πτυχή είναι η διχοτόμος.

\begin{prosoxi}
\begin{itemize}[leftmargin=1.2em,itemsep=0.3ex,topsep=2pt]
\item Το μέτρο της γωνίας δεν εξαρτάται από το μήκος των πλευρών της: μια γωνία με μεγαλύτερο σχέδιο δεν είναι μεγαλύτερη.
\item Διαβάζω τη σωστή κλίμακα του μοιρογνωμονίου: αν η γωνία είναι οξεία το μέτρο της είναι μικρότερο από $90^\circ$.
\item Το $1^\circ=60'$ και όχι $100'$: $25^\circ30'=25{,}5^\circ$, όχι $25{,}30^\circ$.
\item Η διχοτόμος είναι ημιευθεία μέσα στη γωνία, όχι τμήμα που ενώνει δύο σημεία των πλευρών.
\end{itemize}
\end{prosoxi}

\tmima{Λυμένα παραδείγματα}

\paradeigma{Το μονοπάτι του κήπου}
\tixero Η γωνία του φράχτη είναι $84^\circ$.\par
\tithelo Η γωνία που σχηματίζει το μονοπάτι με κάθε πλευρά, αν ήθελαν να τη χωρίσουν σε δύο ίσα παρτέρια.\par
\possindeo Το μονοπάτι θα είναι η διχοτόμος. Κάθε γωνία είναι $84^\circ:2=42^\circ$.\par
\apantisi $42^\circ$.\par
\elegxos $42^\circ+42^\circ=84^\circ$.

\paradeigma{Οι δείκτες του ρολογιού}
\tixero Ο ωροδείκτης και ο λεπτοδείκτης στις $4{:}00$ και στις $5{:}00$.\par
\tithelo Η γωνία που σχηματίζουν.\par
\possindeo Κάθε ώρα αντιστοιχεί σε $360^\circ:12=30^\circ$. Στις $4{:}00$ ο λεπτοδείκτης δείχνει το $12$ και ο ωροδείκτης το $4$: $4\cdot30^\circ=120^\circ$. Στις $5{:}00$: $5\cdot30^\circ=150^\circ$.\par
\apantisi $120^\circ$ και $150^\circ$.\par
\elegxos Στις $6{:}00$ θα είναι $180^\circ$ (ευθεία).

\paradeigma{Μοίρες και λεπτά}
\tixero Οι γωνίες $25^\circ30'$ και $48{,}25^\circ$.\par
\tithelo Να τις γράψω στην άλλη μορφή.\par
\possindeo $30'=\frac{30}{60}^\circ=0{,}5^\circ$, άρα $25^\circ30'=25{,}5^\circ$. Και $0{,}25^\circ=0{,}25\cdot60'=15'$, άρα $48{,}25^\circ=48^\circ15'$.\par
\apantisi $25{,}5^\circ$ και $48^\circ15'$.\par
\elegxos $15'=\frac{15}{60}=0{,}25^\circ$.

\paradeigma{Σύγκριση γωνιών}
\tixero Οι γωνίες $42^\circ$, $4{,}5^\circ$, $38^\circ30'$, $38^\circ$.\par
\tithelo Η σειρά.\par
\possindeo Σε μοίρες: $42$, $4{,}5$, $38{,}5$, $38$.\par
\apantisi $4{,}5^\circ<38^\circ<38^\circ30'<42^\circ$.\par
\elegxos $30'=0{,}5^\circ$.

\paradeigma{Η διχοτόμος και το μέτρο}
\tixero Η $\mathrm{Οδ}$ είναι διχοτόμος της $x\mathrm{Ο}y$ και $x\mathrm{Ο}\delta=28^\circ$.\par
\tithelo Η $x\mathrm{Ο}y$.\par
\possindeo $\delta\mathrm{Ο}y=x\mathrm{Ο}\delta=28^\circ$, άρα $x\mathrm{Ο}y=28^\circ+28^\circ=56^\circ$.\par
\apantisi $56^\circ$.\par
\elegxos Η διχοτόμος της $56^\circ$ δίνει $28^\circ$.

\tmima{Ασκήσεις}

\zoni{Το κατάλαβα}{μοίρες, ρολόι, διχοτόμος}

\ex Σχεδίασε με μοιρογνωμόνιο γωνίες $40^\circ$, $75^\circ$ και $130^\circ$ και έλεγξε με το μέτρημα της διχοτόμου.

\ex Γράψε στην άλλη μορφή: α) $25^\circ30'$ σε μοίρες \quad β) $48{,}25^\circ$ σε μοίρες και λεπτά \quad γ) $3{,}5^\circ$ σε μοίρες και λεπτά \quad δ) $90^\circ$ σε λεπτά.

\ex Βάλε σε αύξουσα σειρά: $42^\circ$, $38^\circ30'$, $38^\circ$, $4{,}5^\circ$, $0{,}7^\circ$.

\ex Βρες τη διχοτόμο: α) $x\mathrm{Ο}y=70^\circ$ \quad β) $x\mathrm{Ο}y=124^\circ$ \quad γ) $x\mathrm{Ο}y=8^\circ30'$. Ποιο είναι το μέτρο κάθε μισής γωνίας;

\ex Η $\mathrm{Οδ}$ είναι διχοτόμος της $x\mathrm{Ο}y$ και $x\mathrm{Ο}\delta=28^\circ$. Βρες την $x\mathrm{Ο}y$.

\ex Βρες τη γωνία των δεικτών του ρολογιού στις: α) $3{:}00$ \quad β) $6{:}00$ \quad γ) $9{:}00$ \quad δ) $2{:}00$ \quad ε) $5{:}00$ \quad στ) $12{:}00$. (Πάρε τη μικρότερη γωνία.)

\ex Σωστό ή λάθος; α) $1^\circ=60'$. \quad β) Η διχοτόμος χωρίζει τη γωνία σε δύο άνισα μέρη. \quad γ) Δύο γωνίες είναι ίσες αν εφαρμόζουν. \quad δ) Η διχοτόμος είναι ημιευθεία. \quad ε) $1'=60''$.

\zoni{Το χρησιμοποιώ}{κήπος, ρολόι, πίτσα}

\ex \textbf{Το μονοπάτι.} Η γωνία του φράχτη είναι $84^\circ$. Με ποια γωνία ξεκινά το μονοπάτι που τη χωρίζει σε δύο ίσα παρτέρια;

\ex \textbf{Το ρολόι.} Πότε (ποια ώρα, με τον λεπτοδείκτη στο $12$) οι δείκτες σχηματίζουν γωνία $150^\circ$; Υπάρχουν δύο ώρες;

\ex \textbf{Η πίτσα.} Μια πίτσα κόβεται σε $8$ ίσες φέτες. Πόση γωνία έχει η κορυφή κάθε φέτας; Αν κόβεται σε $6$ φέτες;

\ex \textbf{Ο ανεμόμυλος.} Ένας ανεμόμυλος έχει $3$ πτερύγια με ίσες γωνίες μεταξύ τους. Πόση γωνία σχηματίζουν δύο διαδοχικά;

\ex \textbf{Το μοιρογνωμόνιο.} Ο Στέλιος διάβασε $140^\circ$ σε μια γωνία που φαίνεται οξεία. Τι πιθανότατα έκανε; Ποιο είναι το σωστό;

\zoni{Το σκέφτομαι}{γιατί $360$;}

\ex Οι Βαβυλώνιοι χωρίζουν τον κύκλο σε $360$ μοίρες. Γιατί ο $360$ είναι βολικός αριθμός (σκέψου τους διαιρέτες του);

\ex Αν διπλασιάσεις το μήκος των πλευρών μιας γωνίας, αλλάζει το μέτρο της; Εξήγησε.

\ex Γιατί κάθε γωνία έχει μία και μοναδική διχοτόμο;

\ex Η διχοτόμος μιας ορθής γωνίας ($90^\circ$) πόσες μοίρες σχηματίζει με κάθε πλευρά;

\zoni{Για το σπίτι}{δίπλωσε το χαρτί}

\ex Κόψε μια γωνία από χαρτί και βρες τη διχοτόμο διπλώνοντάς την. Μέτρησε με μοιρογνωμόνιο.

\ex Μέτρησε με μοιρογνωμόνιο τρεις γωνίες στο σπίτι και γράψε τα μέτρα τους.

\paragrafos{Είδη γωνιών και κάθετες ευθείες}{Β.1.6, «Είδη γωνιών – Κάθετες ευθείες»}

\begin{anagke}
Η Α3 πρέπει να στήσει τον φράχτη ώστε ο κήπος να έχει «σωστές γωνίες» στις γωνιές του. «Σωστή γωνία σημαίνει τετράγωνη, όπως στο θρανίο», λέει ο Νίκος. Η Αγγελική το λέει επίσημα: ορθή γωνία. Αλλά πώς φτιάχνουν ορθή γωνία στο έδαφος, χωρίς μοιρογνωμόνιο; Ο παππούς δείχνει ένα κόλπο που ξέρουν οι χτίστες από την αρχαιότητα: ένα σχοινί με $12$ ίσα διαστήματα, με κόμπους. Το διπλώνει σε τρίγωνο με πλευρές $3$, $4$ και $5$ διαστημάτων: «Η γωνία ανάμεσα στο $3$ και στο $4$ είναι ορθή».

Η Πολυξένη το καταγράφει και ρωτάει για όλες τις άλλες γωνίες που βλέπει: οξείες, αμβλείες, όπως στο φτερό του αεροπλάνου και στη στέγη. Ο Βασίλης θέλει να ταξινομήσει τις γωνίες του κήπου: «Αυτή που ανοίγει λίγο», «αυτή που ανοίγει πολύ», «η ευθεία». Η κυρία Αλεξάνδρα του δίνει ονόματα. Και όταν δύο μονοπάτια τέμνονται σε ορθή γωνία, τα λέμε κάθετα.

Σε αυτή την παράγραφο ταξινομούμε τις γωνίες ανάλογα με το μέτρο τους και ορίζουμε την καθετότητα, μια από τις πιο χρήσιμες έννοιες στη Γεωμετρία.
\end{anagke}

\begin{lexi}{ορθή, οξεία, αμβλεία, κάθετος}{ορθή, οξεία, αμβλεία}
\lexiword{Ορθή} σημαίνει «σωστή, όρθια» (από το «ορθός»): η γωνία ενός όρθιου τοίχου με το δάπεδο. \lexiword{Οξεία} είναι η «αιχμηρή», σαν την οξεία μύτη. \lexiword{Αμβλεία} είναι η «αμβλυμένη», η στρογγυλεμένη, που ανοίγει περισσότερο από την ορθή. \lexiword{Κάθετη} είναι η ευθεία που πέφτει «κάθετα», σαν το νήμα της στάθμης.
\end{lexi}

\begin{metafrasi}
Γωνία μικρότερη από την ορθή. & οξεία: $0^\circ<\omega<90^\circ$\\
Γωνία $90^\circ$. & ορθή\\
Γωνία ανάμεσα στην ορθή και την ευθεία. & αμβλεία: $90^\circ<\omega<180^\circ$\\
Γωνία $180^\circ$. & ευθεία γωνία\\
Δύο ευθείες που τέμνονται σε ορθή γωνία. & $\varepsilon_1\perp\varepsilon_2$\\
\end{metafrasi}

\tmima{Η θεωρία}

\begin{orismos}{είδη γωνιών}
Μια γωνία $\omega$ λέγεται:
\begin{itemize}[leftmargin=1.2em,itemsep=0.3ex,topsep=2pt]
\item \textbf{οξεία} αν $0^\circ<\omega<90^\circ$,
\item \textbf{ορθή} αν $\omega=90^\circ$,
\item \textbf{αμβλεία} αν $90^\circ<\omega<180^\circ$,
\item \textbf{ευθεία} αν $\omega=180^\circ$ (οι πλευρές της είναι αντικείμενες ημιευθείες),
\item \textbf{μη κυρτή} αν $180^\circ<\omega<360^\circ$ και \textbf{πλήρης} αν $\omega=360^\circ$,
\item \textbf{μηδενική} αν $\omega=0^\circ$ (οι πλευρές συμπίπτουν).
\end{itemize}
\end{orismos}

\begin{center}
\begin{tikzpicture}[scale=0.8]
\foreach \a/\n/\x in {40/οξεία/0,90/ορθή/3.6,130/αμβλεία/7.2}{
\begin{scope}[xshift=\x cm]
\draw[thick] (2.2,0)--(0,0)--({2.2*cos(\a)},{2.2*sin(\a)});
\draw[gblue] (0.6,0) arc(0:\a:0.6);
\node at (1.1,-0.7) {\n};
\end{scope}}
\draw[gblue] (3.6+0.25,0)--(3.6+0.25,0.25)--(3.6,0.25);
\end{tikzpicture}
\end{center}

\tmimasmall{Κάθετες ευθείες}

Δύο ευθείες που τέμνονται σχηματίζουν τέσσερις γωνίες. Αν μία από αυτές είναι ορθή, τότε είναι ορθές και οι τέσσερις. Τότε οι ευθείες λέγονται \textbf{κάθετες} και γράφουμε $\varepsilon_1\perp\varepsilon_2$.

\begin{idiotita}{ύπαρξη και μοναδικότητα κάθετης}
Από ένα σημείο του επιπέδου διέρχεται \textbf{μία και μοναδική} ευθεία κάθετη προς μια δοσμένη ευθεία $(\varepsilon)$, είτε το σημείο ανήκει στην ευθεία είτε όχι.
\end{idiotita}

Τις κάθετες τις σχεδιάζουμε με \textbf{γνώμονα} (τρίγωνο με ορθή γωνία) και χάρακα.

\begin{prosoxi}
\begin{itemize}[leftmargin=1.2em,itemsep=0.3ex,topsep=2pt]
\item Η ορθή γωνία είναι ακριβώς $90^\circ$. Η $91^\circ$ είναι αμβλεία και η $89^\circ$ οξεία.
\item Οι κάθετες δεν χρειάζεται να είναι «οριζόντιες και κατακόρυφες»: δύο ευθείες που βρίσκονται σε οποιαδήποτε θέση είναι κάθετες αν σχηματίζουν ορθή γωνία.
\item Η ευθεία γωνία δεν είναι «ευθεία γραμμή» αλλά γωνία $180^\circ$ με αντικείμενες πλευρές.
\item Από σημείο εκτός ευθείας διέρχεται \emph{μία} κάθετη, αλλά άπειρες πλάγιες.
\end{itemize}
\end{prosoxi}

\tmima{Λυμένα παραδείγματα}

\paradeigma{Ταξινομώ γωνίες}
\tixero Οι γωνίες $37^\circ$, $90^\circ$, $123^\circ$, $180^\circ$, $250^\circ$, $360^\circ$.\par
\tithelo Να πω το είδος κάθε μίας.\par
\possindeo Συγκρίνω με τα $90^\circ$, $180^\circ$, $360^\circ$.\par
\apantisi $37^\circ$ οξεία, $90^\circ$ ορθή, $123^\circ$ αμβλεία, $180^\circ$ ευθεία, $250^\circ$ μη κυρτή, $360^\circ$ πλήρης.\par
\elegxos Η $250^\circ$ είναι ανάμεσα στο $180^\circ$ και στο $360^\circ$.

\paradeigma{Το ρολόι}
\tixero Οι δείκτες στις $2{:}00$, $3{:}00$, $4{:}00$.\par
\tithelo Το είδος της γωνίας.\par
\possindeo Κάθε ώρα είναι $30^\circ$. Στις $2{:}00$: $60^\circ$ οξεία. Στις $3{:}00$: $90^\circ$ ορθή. Στις $4{:}00$: $120^\circ$ αμβλεία.\par
\apantisi Οξεία, ορθή, αμβλεία.\par
\elegxos Στις $6{:}00$ θα ήταν ευθεία γωνία.

\paradeigma{Κάθετες ευθείες}
\tixero Δύο ευθείες που τέμνονται σχηματίζοντας γωνία $90^\circ$.\par
\tithelo Οι υπόλοιπες γωνίες.\par
\possindeo Οι διαδοχικές γωνίες ανήκουν σε ευθεία γωνία $180^\circ$: $180^\circ-90^\circ=90^\circ$. Έτσι και οι τέσσερις είναι ορθές.\par
\apantisi $90^\circ$ η καθεμία.\par
\elegxos Το άθροισμα και των τεσσάρων είναι $360^\circ$.

\paradeigma{Η ορθή γωνία του παππού}
\tixero Σχοινί με $12$ ίσα διαστήματα, διπλωμένο σε τρίγωνο με πλευρές $3$, $4$, $5$.\par
\tithelo Γιατί σχηματίζεται ορθή γωνία.\par
\possindeo Το τρίγωνο με πλευρές $3$, $4$, $5$ έχει ορθή γωνία ανάμεσα στις πλευρές $3$ και $4$ (το αποδεικνύουμε αργότερα). Οι αρχαίοι αρπεδονάπτες το χρησιμοποιούσαν για τις γωνίες των κτιρίων και των αγρών.\par
\apantisi Γιατί $3^2+4^2=5^2$ ($9+16=25$).\par
\elegxos Δοκίμασε με γνώμονα.

\tmima{Ασκήσεις}

\zoni{Το κατάλαβα}{είδη γωνιών και κάθετες}

\ex Ταξινόμησε τις γωνίες $37^\circ$, $90^\circ$, $123^\circ$, $180^\circ$, $250^\circ$, $360^\circ$, $0^\circ$, $89^\circ$.

\ex Σχεδίασε με μοιρογνωμόνιο μια οξεία, μια ορθή και μια αμβλεία γωνία και μέτρησε τις διχοτόμους τους.

\ex Βρες το είδος της γωνίας των δεικτών του ρολογιού στις $2{:}00$, $3{:}00$, $4{:}00$, $6{:}00$, $7{:}00$ (πάρε τη μικρότερη).

\ex Δύο ευθείες τέμνονται και μια γωνία είναι $90^\circ$. Πόσες είναι οι άλλες γωνίες;

\ex Με γνώμονα σχεδίασε μια ευθεία $(\varepsilon)$, ένα σημείο $\mathrm{Α}$ πάνω της και ένα σημείο $\mathrm{Β}$ εκτός. Φέρε από το $\mathrm{Α}$ και από το $\mathrm{Β}$ τις κάθετες προς την $(\varepsilon)$.

\ex Πόσες κάθετες διέρχονται από ένα σημείο προς μια ευθεία; Και πόσες πλάγιες;

\ex Σωστό ή λάθος; α) Μια γωνία $91^\circ$ είναι αμβλεία. \quad β) Η ευθεία γωνία είναι $180^\circ$. \quad γ) Δύο κάθετες ευθείες σχηματίζουν γωνίες που δεν είναι όλες ίσες. \quad δ) Η $0^\circ$ είναι οξεία. \quad ε) Η $250^\circ$ είναι κυρτή.

\zoni{Το χρησιμοποιώ}{γωνιές, στέγες, αναβάτες}

\ex \textbf{Ο φράχτης.} Στις γωνιές του κήπου ο παππούς θέλει ορθές γωνίες. Πώς θα το ελέγξεις με ένα φύλλο χαρτί $\mathrm{Α4}$;

\ex \textbf{Το σχοινί του παππού.} Ένα σχοινί με $12$ ίσα διαστήματα διπλώνεται σε τρίγωνο με πλευρές $3$, $4$, $5$. Τι γωνία σχηματίζεται και πού;

\ex \textbf{Τα δωμάτια.} Βρες στην τάξη τρεις ορθές, δύο οξείες και μία αμβλεία γωνία.

\ex \textbf{Η στέγη.} Μια στέγη έχει γωνία $110^\circ$ στην κορυφή. Τι είδους είναι; Αν η στέγη ήταν ορθή γωνία πόσες μοίρες θα ήταν;

\ex \textbf{Ο ανεμοδείκτης.} Ένας ανεμοδείκτης στρέφεται από τη διεύθυνση Βορράς προς Ανατολή. Τι γωνία διέγραψε; Και προς Νότο;

\zoni{Το σκέφτομαι}{όταν μία ορθή φέρνει τρεις}

\ex Δείξε ότι αν μία από τις γωνίες που σχηματίζουν δύο τεμνόμενες ευθείες είναι ορθή, είναι ορθές και οι άλλες τρεις.

\ex Είναι δυνατόν να έχει ένα τρίγωνο δύο ορθές γωνίες; Δοκίμασε να σχεδιάσεις.

\ex Γιατί ο γνώμονας φτιάχνει κάθετες; Πώς θα έλεγχες ότι είναι σωστός (σκέψου την αντιστροφή του);

\ex Αν η γωνία $\omega$ είναι οξεία, τι είδους είναι η $180^\circ-\omega$;

\zoni{Για το σπίτι}{ψάξε γωνίες}

\ex Βρες στο σπίτι πέντε αντικείμενα με ορθές γωνίες και τρία με αμβλείες.

\ex Φτιάξε με σπάγκο και κόμπους το τρίγωνο $3,\,4,\,5$ και έλεγξε με ορθογώνιο χαρτί αν η γωνία είναι ορθή.

\paragrafos{Εφεξής και διαδοχικές γωνίες: οι γωνίες που προστίθενται}{Β.1.7, «Εφεξής και διαδοχικές γωνίες – Άθροισμα γωνιών»}

\begin{anagke}
Στο κέντρο του κήπου θα υπάρχει το συντριβάνι και από αυτό θα ξεκινούν πέντε μονοπάτια προς διαφορετικές κατευθύνσεις. Η Αγγελική σχεδιάζει το σχήμα και μετράει τις γωνίες ανάμεσα στα διαδοχικά μονοπάτια: $60^\circ$, $75^\circ$, $50^\circ$, $90^\circ$. «Η πέμπτη γωνία δεν χωρούσε στο μοιρογνωμόνιο», λέει, «γιατί είναι μη κυρτή;». «Όχι», απαντάει η Πολυξένη, «απλώς πρέπει να βρούμε πόσο μένει μέχρι τον κύκλο. Ολόκληρος ο κύκλος έχει $360^\circ$».

Ο Βασίλης προσθέτει: $60+75+50+90=275$, άρα η πέμπτη είναι $360-275=85$. «Μπορεί να μην βρίσκομαι εκεί, αλλά είμαι σίγουρος!». Ο Στέλιος ρωτάει αν οι γωνίες που έχουν κοινή πλευρά και «χωρίζουν» το μέρος που έχουν κοινή κορυφή έχουν κάποιο όνομα. Η κυρία Αλεξάνδρα απαντάει: «Λέγονται εφεξής, και όταν είναι πολλές στη σειρά, διαδοχικές. Τα μέτρα τους προστίθενται».

Στην παράγραφο αυτή μαθαίνουμε πώς προσθέτουμε και αφαιρούμε γωνίες που έχουν κοινή κορυφή και κοινή πλευρά, και πώς χρησιμοποιούμε το γεγονός ότι γύρω από ένα σημείο υπάρχουν $360^\circ$ και πάνω σε μια ευθεία $180^\circ$.
\end{anagke}

\begin{lexi}{εφεξής, διαδοχικές}{εφεξής, διαδοχικές}
\lexiword{Εφεξής} σημαίνει «η μία μετά την άλλη, στη σειρά» (εφ' εξής). \lexiword{Διαδοχικές} είναι οι γωνίες που «διαδέχονται» η μία την άλλη, έτσι ώστε η κάθε μία να είναι εφεξής με την επόμενη.
\end{lexi}

\begin{metafrasi}
Δύο γωνίες με κοινή κορυφή και κοινή πλευρά, και χωριστά εσωτερικά. & εφεξής γωνίες $x\mathrm{Ο}y$, $y\mathrm{Ο}z$\\
Το άθροισμά τους. & $x\mathrm{Ο}z=x\mathrm{Ο}y+y\mathrm{Ο}z$\\
Όλες οι διαδοχικές γύρω από ένα σημείο. & άθροισμα $360^\circ$\\
Οι διαδοχικές πάνω σε ευθεία. & άθροισμα $180^\circ$\\
\end{metafrasi}

\tmima{Η θεωρία}

\begin{orismos}{εφεξής και διαδοχικές γωνίες}
Δύο γωνίες $x\mathrm{Ο}y$ και $y\mathrm{Ο}z$ λέγονται \textbf{εφεξής} όταν έχουν κοινή κορυφή $\mathrm{Ο}$, κοινή πλευρά $\mathrm{Οy}$ και τα εσωτερικά τους δεν έχουν κοινά σημεία (η μία είναι «πριν» και η άλλη «μετά» την $\mathrm{Οy}$). Πολλές γωνίες λέγονται \textbf{διαδοχικές} όταν καθεμία είναι εφεξής με την επόμενή της.
\end{orismos}

\begin{center}
\begin{tikzpicture}[scale=1.1]
\draw[thick] (3.4,0)--(0,0)--(2.4,2.4); \draw[thick,gblue] (0,0)--(0.6,3.1); \fill (0,0) circle (2pt);
\node[right] at (3.4,0) {$x$}; \node[right] at (2.4,2.4) {$y$}; \node[above] at (0.6,3.1) {$z$}; \node[below left] at (0,0) {$\mathrm{Ο}$};
\draw (0.8,0) arc(0:45:0.8); \draw[gblue] (0.55,0.55) arc(45:79:0.55);
\end{tikzpicture}
\end{center}

\tmimasmall{Άθροισμα γωνιών}

Αν οι $x\mathrm{Ο}y$ και $y\mathrm{Ο}z$ είναι εφεξής, η γωνία $x\mathrm{Ο}z$ λέγεται \textbf{άθροισμά} τους:
$$x\mathrm{Ο}z=x\mathrm{Ο}y+y\mathrm{Ο}z.$$
Και αντίστροφα η διαφορά: $y\mathrm{Ο}z=x\mathrm{Ο}z-x\mathrm{Ο}y$. Το άθροισμα των μέτρων μπορεί να ξεπεράσει τα $180^\circ$ (τότε η γωνία είναι μη κυρτή).

\begin{idiotita}{άθροισμα διαδοχικών γωνιών}
Το άθροισμα των διαδοχικών γωνιών γύρω από ένα σημείο, που καλύπτουν όλο το επίπεδο, είναι $360^\circ$. Το άθροισμα διαδοχικών γωνιών που βρίσκονται από την ίδια μεριά μιας ευθείας (πλευρές αντικείμενες ημιευθείες) είναι $180^\circ$.
\end{idiotita}

\begin{prosoxi}
\begin{itemize}[leftmargin=1.2em,itemsep=0.3ex,topsep=2pt]
\item Δύο γωνίες δεν είναι εφεξής αν δεν έχουν κοινή κορυφή ή αν τα εσωτερικά τους επικαλύπτονται.
\item Το άθροισμα γωνιών αθροίζει τα \emph{μέτρα}: $40^\circ+65^\circ=105^\circ$.
\item Στον κύκλο, όλο το άθροισμα είναι $360^\circ$, όχι $180^\circ$ (αυτό είναι για την ευθεία).
\item Για εφεξής γωνίες το άθροισμα μπορεί να ξεπεράσει τις $180^\circ$ αλλά η γωνία $x\mathrm{Ο}z$ είναι τότε μη κυρτή.
\end{itemize}
\end{prosoxi}

\tmima{Λυμένα παραδείγματα}

\paradeigma{Εφεξής γωνίες}
\tixero $x\mathrm{Ο}y=35^\circ$ και $y\mathrm{Ο}z=50^\circ$ εφεξής.\par
\tithelo Η $x\mathrm{Ο}z$.\par
\possindeo $x\mathrm{Ο}z=35^\circ+50^\circ=85^\circ$.\par
\apantisi $85^\circ$.\par
\elegxos $85^\circ-35^\circ=50^\circ$.

\paradeigma{Με αφαίρεση}
\tixero Οι $\mathrm{ΑΟΒ}$, $\mathrm{ΒΟΓ}$ εφεξής με $\mathrm{ΑΟΒ}=70^\circ$ και $\mathrm{ΑΟΓ}=130^\circ$.\par
\tithelo Η $\mathrm{ΒΟΓ}$.\par
\possindeo $\mathrm{ΒΟΓ}=\mathrm{ΑΟΓ}-\mathrm{ΑΟΒ}=130^\circ-70^\circ=60^\circ$.\par
\apantisi $60^\circ$.\par
\elegxos $70^\circ+60^\circ=130^\circ$.

\paradeigma{Τα πέντε μονοπάτια}
\tixero Τέσσερις διαδοχικές γωνίες γύρω από το συντριβάνι: $60^\circ$, $75^\circ$, $50^\circ$, $90^\circ$.\par
\tithelo Η πέμπτη γωνία.\par
\possindeo Το άθροισμα γύρω από το σημείο είναι $360^\circ$. Τέσσερις γωνίες: $60+75+50+90=275$. Άρα η πέμπτη είναι $360-275=85^\circ$.\par
\apantisi $85^\circ$.\par
\elegxos $275^\circ+85^\circ=360^\circ$.

\paradeigma{Γωνίες πάνω σε ευθεία}
\tixero Διαδοχικές γωνίες $40^\circ$, $65^\circ$ και $\omega$ πάνω από μια ευθεία (η τρίτη κλείνει στην ευθεία).\par
\tithelo Η $\omega$.\par
\possindeo Το άθροισμα είναι $180^\circ$: $40+65+\omega=180$, άρα $\omega=180-105=75^\circ$.\par
\apantisi $75^\circ$.\par
\elegxos $40+65+75=180$.

\paradeigma{Διχοτόμος και άθροισμα}
\tixero Η $\mathrm{Οδ}$ είναι διχοτόμος της $x\mathrm{Ο}y$ και $\delta\mathrm{Ο}y=37^\circ$. Η $y\mathrm{Ο}z=48^\circ$ είναι εφεξής στην $x\mathrm{Ο}y$.\par
\tithelo Η $x\mathrm{Ο}z$.\par
\possindeo $x\mathrm{Ο}y=2\cdot37^\circ=74^\circ$. $x\mathrm{Ο}z=74^\circ+48^\circ=122^\circ$.\par
\apantisi $122^\circ$.\par
\elegxos Η $\mathrm{Οδ}$ πλευρά της $x\mathrm{Ο}z$: $37+37+48=122$.

\tmima{Ασκήσεις}

\zoni{Το κατάλαβα}{άθροισμα και διαφορά γωνιών}

\ex Οι $x\mathrm{Ο}y=35^\circ$ και $y\mathrm{Ο}z=50^\circ$ είναι εφεξής. Βρες την $x\mathrm{Ο}z$.

\ex Οι $\mathrm{ΑΟΒ}$ και $\mathrm{ΒΟΓ}$ είναι εφεξής, $\mathrm{ΑΟΓ}=130^\circ$ και $\mathrm{ΑΟΒ}=70^\circ$. Βρες την $\mathrm{ΒΟΓ}$.

\ex Τέσσερις διαδοχικές γωνίες γύρω από σημείο είναι $60^\circ$, $75^\circ$, $50^\circ$, $90^\circ$. Βρες την πέμπτη.

\ex Οι $40^\circ$, $65^\circ$ και $\omega$ είναι διαδοχικές πάνω από ευθεία. Βρες την $\omega$.

\ex Τρεις διαδοχικές γωνίες γύρω από σημείο είναι ίσες. Πόσες μοίρες είναι η καθεμία;

\ex Οι γωνίες $x\mathrm{Ο}y=24^\circ30'$ και $y\mathrm{Ο}z=37^\circ45'$ είναι εφεξής. Βρες την $x\mathrm{Ο}z$.

\ex Σωστό ή λάθος; α) Δύο γωνίες με κοινή κορυφή είναι πάντα εφεξής. \quad β) Οι εφεξής γωνίες έχουν κοινή πλευρά. \quad γ) Διαδοχικές γωνίες γύρω από σημείο έχουν άθροισμα $360^\circ$. \quad δ) Διαδοχικές γωνίες πάνω από ευθεία έχουν άθροισμα $180^\circ$.

\zoni{Το χρησιμοποιώ}{κήπος, πίτσα, ρολόι}

\ex \textbf{Τα πέντε μονοπάτια.} Στο συντριβάνι καταλήγουν πέντε μονοπάτια και οι γωνίες είναι $60^\circ$, $75^\circ$, $50^\circ$, $90^\circ$ και $\omega$. Βρες την $\omega$.

\ex \textbf{Η πίτσα.} Μια πίτσα κόβεται σε $6$ φέτες: δύο από $75^\circ$ και τρεις από $50^\circ$. Πόσο είναι η έκτη;

\ex \textbf{Ο ήλιος.} Ο Ήλιος σηκώνεται από τον ορίζοντα $12^\circ$ κάθε $48$ λεπτά. Σε πόση ώρα θα φτάσει να ανεβεί $60^\circ$;

\ex \textbf{Το ρολόι.} Στις $4{:}30$ ο ωροδείκτης βρίσκεται στις $135^\circ$ από το $12$ και ο λεπτοδείκτης στο $6$ ($180^\circ$ από το $12$). Πόση γωνία σχηματίζουν;

\ex \textbf{Οι στίβοι.} Τρεις στίβοι ξεκινούν από το ίδιο σημείο σχηματίζοντας ίσες γωνίες. Πόση είναι η γωνία ανάμεσα σε δύο διαδοχικούς;

\zoni{Το σκέφτομαι}{όταν το σύνολο είναι σταθερό}

\ex Γιατί το άθροισμα των διαδοχικών γωνιών γύρω από σημείο είναι $360^\circ$ ανεξάρτητα από τον αριθμό των γωνιών;

\ex Είναι δυνατόν δύο εφεξής γωνίες να είναι και οι δύο αμβλείες; Εξήγησε.

\ex Τρεις διαδοχικές γωνίες γύρω από σημείο είναι $x$, $2x$, $3x$. Βρες το $x$.

\ex Η $\mathrm{Οδ}$ είναι διχοτόμος της $x\mathrm{Ο}y=60^\circ$ και η $y\mathrm{Ο}z=20^\circ$ είναι εφεξής προς αυτήν. Βρες τις $x\mathrm{Ο}z$ και $\delta\mathrm{Ο}z$.

\zoni{Για το σπίτι}{άθροισμα στο χαρτί}

\ex Σχεδίασε τέσσερις διαδοχικές γωνίες γύρω από σημείο και μέτρησέ τις. Πόσο είναι το άθροισμά τους;

\ex Βρες στο σπίτι ένα αντικείμενο με διαδοχικές γωνίες γύρω από σημείο (π.χ. ένα ρολόι ή μια ανεμογεννήτρια).

\paragrafos{Παραπληρωματικές, συμπληρωματικές και κατακορυφήν γωνίες}{Β.1.8, «Παραπληρωματικές και Συμπληρωματικές γωνίες – Κατακορυφήν γωνίες»}

\begin{anagke}
Δύο μονοπάτια του κήπου διασταυρώνονται σαν ένα μεγάλο «Χ». Η Αγγελική μετράει τη μία γωνία και βρίσκει $48^\circ$. «Πόσες είναι οι άλλες τρεις;» ρωτάει. Ο Νίκος μετράει όλες: $48^\circ$, $132^\circ$, $48^\circ$, $132^\circ$. «Οι απέναντι είναι ίσες! Και οι διπλανές κάνουν μαζί $180^\circ$». «Είναι τυχαίο;» ρωτάει η Πολυξένη. «Μήπως αν κάνω το μονοπάτι πιο λοξό, θα αλλάξει;». Δοκιμάζουν με άλλη γωνία, $70^\circ$, και βρίσκουν $70^\circ$, $110^\circ$, $70^\circ$, $110^\circ$.

Η Ρήνη ρωτάει για μια γωνία που τη συμπληρώνουμε μέχρι την ορθή: ένα παρτέρι σχήματος γωνίας $35^\circ$ και το υπόλοιπο μέχρι τη γωνία του φράχτη των $90^\circ$ θα φυτευτεί με τριαντάφυλλα: «Πόσες μοίρες μένουν;». Ο Στέλιος υπολογίζει: «Εδώ το άθροισμα είναι $90^\circ$ και όχι $180^\circ$».

Έτσι η παρέα συναντά τρεις έννοιες: τις παραπληρωματικές γωνίες, που αθροίζουν $180^\circ$, τις συμπληρωματικές, που αθροίζουν $90^\circ$, και τις κατακορυφήν, που βρίσκονται απέναντι όταν δύο ευθείες τέμνονται. Θα βρουν και γιατί οι κατακορυφήν γωνίες είναι ίσες, χωρίς να χρειάζονται μοιρογνωμόνιο.
\end{anagke}

\begin{lexi}{παραπληρωματική, συμπληρωματική, κατακορυφήν}{παραπληρωματική, συμπληρωματική, κατακορυφήν}
Η \lexiword{συμπληρωματική} γωνία «συμπληρώνει» μια γωνία μέχρι την ορθή. Η \lexiword{παραπληρωματική} «παρα-πληρώνει», δηλαδή «συμπληρώνει ως τη μισή στροφή»: ως τις $180^\circ$. Οι \lexiword{κατακορυφήν} γωνίες είναι αυτές που βρίσκονται «κατά κορυφήν», δηλαδή απέναντι, με κοινή κορυφή.
\end{lexi}

\begin{metafrasi}
Δύο γωνίες που αθροίζουν $90^\circ$. & συμπληρωματικές: $\omega+\varphi=90^\circ$\\
Δύο γωνίες που αθροίζουν $180^\circ$. & παραπληρωματικές: $\omega+\varphi=180^\circ$\\
Οι απέναντι γωνίες δύο τεμνόμενων ευθειών. & κατακορυφήν: ίσες\\
Η συμπληρωματική της $35^\circ$. & $90^\circ-35^\circ=55^\circ$\\
Η παραπληρωματική της $130^\circ$. & $180^\circ-130^\circ=50^\circ$\\
\end{metafrasi}

\tmima{Η θεωρία}

\begin{orismos}{συμπληρωματικές, παραπληρωματικές}
Δύο γωνίες λέγονται \textbf{συμπληρωματικές} αν το άθροισμά τους είναι $90^\circ$ και \textbf{παραπληρωματικές} αν το άθροισμά τους είναι $180^\circ$. Η καθεμία λέγεται συμπλήρωμα (ή παραπλήρωμα) της άλλης.
\end{orismos}

Για παράδειγμα η συμπληρωματική της $35^\circ$ είναι η $55^\circ$ και η παραπληρωματική της $130^\circ$ η $50^\circ$. Δύο εφεξής παραπληρωματικές γωνίες έχουν τις μη κοινές πλευρές τους αντικείμενες ημιευθείες. Δύο συμπληρωματικές γωνίες είναι πάντα οξείες, ενώ από δύο παραπληρωματικές η μία είναι αμβλεία και η άλλη οξεία (ή και οι δύο ορθές).

\tmimasmall{Κατακορυφήν γωνίες}

Όταν δύο ευθείες τέμνονται σχηματίζονται τέσσερις γωνίες. Δύο από αυτές που δεν είναι διπλανές λέγονται \textbf{κατακορυφήν}.

\begin{center}
\begin{tikzpicture}[scale=1.1]
\draw[thick,{Stealth}-{Stealth}] (-2.5,-1)--(2.5,1); \draw[thick,{Stealth}-{Stealth}] (-1.2,-1.8)--(1.2,1.8);
\fill (0,0) circle (2pt); \node[below right] at (0.05,0) {$\mathrm{Ο}$};
\node[gblue] at (1.1,0.18) {$\omega$}; \node[gblue] at (-1.1,-0.18) {$\omega$}; \node at (0.18,0.9) {$\varphi$}; \node at (-0.18,-0.9) {$\varphi$};
\end{tikzpicture}
\end{center}

\begin{idiotita}{κατακορυφήν γωνίες}
Οι κατακορυφήν γωνίες είναι \textbf{ίσες}.
\end{idiotita}

Γιατί: η $\omega$ και η διπλανή της $\varphi$ είναι παραπληρωματικές ($\omega+\varphi=180^\circ$). Η κατακορυφήν της $\omega$, έστω $\omega'$, είναι κι αυτή παραπληρωματική της $\varphi$ ($\omega'+\varphi=180^\circ$). Άρα $\omega=180^\circ-\varphi=\omega'$.

\begin{prosoxi}
\begin{itemize}[leftmargin=1.2em,itemsep=0.3ex,topsep=2pt]
\item Συμπληρωματικές είναι $90^\circ$, παραπληρωματικές $180^\circ$. Μην τις μπερδεύεις: η «παραπληρωματική» είναι η μεγαλύτερη λέξη και το μεγαλύτερο άθροισμα ($180^\circ$).
\item Δεν υπάρχει συμπληρωματική γωνία για γωνία μεγαλύτερη ή ίση με $90^\circ$ (θα έβγαινε μηδενική ή αρνητική).
\item Οι διπλανές γωνίες δύο τεμνόμενων ευθειών δεν είναι ίσες (εκτός αν είναι ορθές): είναι παραπληρωματικές.
\item Δύο γωνίες που αθροίζουν $180^\circ$ δεν είναι αναγκαστικά εφεξής: αρκεί το μέτρο.
\end{itemize}
\end{prosoxi}

\tmima{Λυμένα παραδείγματα}

\paradeigma{Συμπλήρωμα και παραπλήρωμα}
\tixero Η γωνία $35^\circ$ και η γωνία $130^\circ$.\par
\tithelo Η συμπληρωματική της πρώτης και η παραπληρωματική της δεύτερης.\par
\possindeo $90^\circ-35^\circ=55^\circ$. $180^\circ-130^\circ=50^\circ$.\par
\apantisi $55^\circ$ και $50^\circ$.\par
\elegxos $35+55=90$ και $130+50=180$.

\paradeigma{Δύο τεμνόμενες ευθείες}
\tixero Δύο ευθείες τέμνονται και μία γωνία είναι $48^\circ$.\par
\tithelo Οι άλλες τρεις γωνίες.\par
\possindeo Η κατακορυφήν της είναι $48^\circ$. Οι διπλανές είναι παραπληρωματικές της: $180^\circ-48^\circ=132^\circ$, και η απέναντί τους επίσης $132^\circ$.\par
\apantisi $48^\circ$, $132^\circ$, $132^\circ$.\par
\elegxos Το άθροισμα και των τεσσάρων: $48+132+48+132=360$.

\paradeigma{Παραπληρωματικές με λόγο}
\tixero Δύο παραπληρωματικές γωνίες έχουν λόγο $2:3$.\par
\tithelo Να τις βρω.\par
\possindeo Αν η μία είναι $2x$ και η άλλη $3x$, $2x+3x=180$, $5x=180$, $x=36$. Οι γωνίες είναι $72^\circ$ και $108^\circ$.\par
\apantisi $72^\circ$ και $108^\circ$.\par
\elegxos $72+108=180$ και $\frac{72}{108}=\frac{2}{3}$.

\paradeigma{Εξίσωση με κατακορυφήν}
\tixero Οι κατακορυφήν γωνίες είναι $2x+10$ και $3x-20$ μοίρες.\par
\tithelo Το $x$ και οι γωνίες.\par
\possindeo Ίσες: $2x+10=3x-20$, άρα $x=30$ και οι γωνίες είναι $2\cdot30+10=70^\circ$.\par
\apantisi $x=30$ και γωνίες $70^\circ$.\par
\elegxos $3\cdot30-20=70$.

\tmima{Ασκήσεις}

\zoni{Το κατάλαβα}{συμπληρωματικές, παραπληρωματικές, κατακορυφήν}

\ex Βρες το συμπλήρωμα των γωνιών: $35^\circ$, $62^\circ$, $8^\circ$, $44^\circ30'$.

\ex Βρες το παραπλήρωμα των γωνιών: $130^\circ$, $90^\circ$, $5^\circ$, $73^\circ15'$.

\ex Δύο ευθείες τέμνονται και μία γωνία είναι $48^\circ$. Βρες τις άλλες τρεις.

\ex Δύο ευθείες τέμνονται και μία γωνία είναι ορθή. Βρες τις άλλες.

\ex Βρες γωνία που είναι ίση με το συμπλήρωμά της. Βρες γωνία ίση με το παραπλήρωμά της.

\ex Δύο παραπληρωματικές γωνίες έχουν λόγο $2:3$. Βρες τες.

\ex Σωστό ή λάθος; α) Οι κατακορυφήν γωνίες είναι ίσες. \quad β) Οι διπλανές γωνίες δύο τεμνόμενων ευθειών είναι συμπληρωματικές. \quad γ) Υπάρχει συμπληρωματική της $100^\circ$. \quad δ) Δύο συμπληρωματικές γωνίες είναι οξείες. \quad ε) Οι παραπληρωματικές γωνίες έχουν άθροισμα $180^\circ$.

\zoni{Το χρησιμοποιώ}{διασταυρώσεις, παρτέρια, καθρέφτες}

\ex \textbf{Η διασταύρωση.} Δύο μονοπάτια του κήπου τέμνονται και μια γωνία είναι $70^\circ$. Πόσες μοίρες έχουν οι άλλες γωνίες; Πόσες διαφορετικές τιμές υπάρχουν;

\ex \textbf{Τα τριαντάφυλλα.} Στη γωνία $90^\circ$ του φράχτη φυτεύεται ένα παρτέρι με γωνία $35^\circ$. Πόση γωνία μένει για τα τριαντάφυλλα;

\ex \textbf{Ο καθρέφτης.} Ένα χέρι του κήπου σχηματίζει γωνία $\omega=2x+10$ και η κατακορυφήν της είναι $3x-20$. Βρες το $x$ και την $\omega$.

\ex \textbf{Η σκάλα.} Μια σκάλα ακουμπά σε έναν τοίχο σχηματίζοντας γωνία $72^\circ$ με το έδαφος. Πόση γωνία σχηματίζει με τον τοίχο; (Ο τοίχος είναι κάθετος στο έδαφος.)

\ex \textbf{Το ψαλίδι.} Ένα ψαλίδι ανοίγει ώστε οι λεπίδες να σχηματίζουν γωνία $40^\circ$. Πόση είναι η γωνία των λαβών (κατακορυφήν);

\zoni{Το σκέφτομαι}{γιατί ίσες}

\ex Εξήγησε γιατί οι κατακορυφήν γωνίες είναι ίσες, χρησιμοποιώντας τα παραπληρώματα.

\ex Δύο γωνίες είναι ίσες. Είναι απαραίτητα κατακορυφήν; Δώσε αντιπαράδειγμα.

\ex Αν η $\omega$ είναι οξεία, τι είδους είναι η παραπληρωματική της; Και η συμπληρωματική της;

\ex Οι διχοτόμοι δύο εφεξής παραπληρωματικών γωνιών: πόση γωνία σχηματίζουν μεταξύ τους; (δοκίμασε με $\omega=50^\circ$).

\zoni{Για το σπίτι}{ψάξε Χ}

\ex Σχεδίασε δύο τεμνόμενες ευθείες, μέτρησε και τις τέσσερις γωνίες και έλεγξε ότι οι κατακορυφήν είναι ίσες.

\ex Βρες στο σπίτι ένα ζευγάρι τεμνόμενων ευθειών (π.χ. ψαλίδι, παραθύρων) και μέτρησε τις γωνίες.

\paragrafos{Θέσεις ευθειών στο επίπεδο: τεμνόμενες και παράλληλες}{Β.1.9, «Θέσεις ευθειών στο επίπεδο»}

\begin{anagke}
Ο παππούς δείχνει πώς θα φυτευτούν οι σειρές με τα καρότα στον κήπο: «Κάθε σειρά θα είναι ευθεία και όλες θα έχουν την ίδια κατεύθυνση. Καμιά δεν θα ακουμπά την άλλη όσο κι αν τις συνεχίσετε». Η Αγγελική ξαφνικά σκέφτεται τις γραμμές στο σχολικό γήπεδο, τις γραμμές στο τετράδιο και τις ράγες του τρένου: «Δεν συναντιούνται ποτέ».

Ο Ηλίας δείχνει δύο μονοπάτια του κήπου που διασταυρώνονται: «Αυτά ακουμπούν σε ένα σημείο». «Μα αν συνεχίσεις ένα μονοπάτι που δεν διασταυρώνεται, μπορεί να διασταυρωθεί αργότερα!» διαμαρτύρεται ο Στέλιος. Η Πολυξένη προτείνει ένα τεστ: «Σκέψου τις ευθείες απεριόριστες. Δύο ευθείες του ίδιου επιπέδου είτε έχουν ένα κοινό σημείο είτε δεν έχουν κανένα».

Στην παράγραφο αυτή ονομάζουμε τις δύο δυνατές θέσεις: τεμνόμενες και παράλληλες. Και ο Ευκλείδης θα μας δώσει ένα μοναδικό αξίωμα για το πώς φέρνουμε παράλληλη από ένα σημείο.
\end{anagke}

\begin{lexi}{παράλληλες, τεμνόμενες}{παράλληλες, τεμνόμενες}
\lexiword{Παράλληλες} είναι οι ευθείες που πηγαίνουν «δίπλα-δίπλα» (παρ' άλληλες), «η μία πλάι στην άλλη», χωρίς ποτέ να συναντιούνται. \lexiword{Τεμνόμενες} είναι οι ευθείες που «τέμνονται», δηλαδή σταυρώνουν.
\end{lexi}

\begin{metafrasi}
Δύο ευθείες χωρίς κοινό σημείο. & παράλληλες: $\varepsilon_1\parallel\varepsilon_2$\\
Δύο ευθείες με ένα κοινό σημείο. & τεμνόμενες\\
Το κοινό τους σημείο. & σημείο τομής\\
Τμήματα πάνω σε παράλληλες ευθείες. & $\mathrm{ΑΒ}\parallel\Gamma\Delta$\\
\end{metafrasi}

\tmima{Η θεωρία}

\begin{orismos}{παράλληλες και τεμνόμενες ευθείες}
Δύο ευθείες του ίδιου επιπέδου λέγονται \textbf{παράλληλες} αν δεν έχουν κανένα κοινό σημείο, όσο κι αν προεκταθούν. Γράφουμε $\varepsilon_1\parallel\varepsilon_2$. Λέγονται \textbf{τεμνόμενες} αν έχουν ένα (και μόνο ένα) κοινό σημείο, που λέγεται \textbf{σημείο τομής}.
\end{orismos}

Δύο ευθείες του ίδιου επιπέδου ή είναι παράλληλες ή τέμνονται (αν έχουν δύο κοινά σημεία ταυτίζονται, αφού από δύο σημεία περνά μία ευθεία). Δύο τμήματα λέγονται \textbf{παράλληλα} όταν βρίσκονται σε παράλληλες ευθείες.

\begin{center}
\begin{tikzpicture}[scale=0.9]
\draw[thick,{Stealth}-{Stealth}] (0,0)--(4,0) node[right] {$\varepsilon_1$}; \draw[thick,{Stealth}-{Stealth}] (0,1.3)--(4,1.3) node[right] {$\varepsilon_2$};
\node at (2,-0.6) {$\varepsilon_1\parallel\varepsilon_2$};
\begin{scope}[xshift=7cm]
\draw[thick,{Stealth}-{Stealth}] (0,0)--(4,1.6); \draw[thick,{Stealth}-{Stealth}] (0,1.5)--(4,-0.1); \fill (2,0.77) circle (2pt);
\node at (2,-0.6) {τεμνόμενες};
\end{scope}
\end{tikzpicture}
\end{center}

\begin{idiotita}{αξίωμα παραλλήλων (Ευκλείδη)}
Από ένα σημείο $\mathrm{Α}$ εκτός ευθείας $(\varepsilon)$ διέρχεται \textbf{μία και μοναδική} ευθεία παράλληλη προς την $(\varepsilon)$.
\end{idiotita}

Το αξίωμα αυτό ανήκει στο έργο του Ευκλείδη, τα «Στοιχεία», που γράφτηκε στην Αλεξάνδρεια περίπου το $300$ π.Χ. Για πολλούς αιώνες οι μαθηματικοί προσπαθούσαν να το αποδείξουν από τα άλλα, και τελικά κατάλαβαν ότι είναι αναγκαίο: υπάρχουν και γεωμετρίες όπου δεν ισχύει.

Μπορούμε να φέρουμε παράλληλη με \textbf{γνώμονα και χάρακα}: ακουμπάμε τον χάρακα στη δοσμένη ευθεία $(\varepsilon)$ και τον γνώμονα πάνω στον χάρακα. Ολισθαίνουμε τον γνώμονα κατά μήκος του χάρακα μέχρι να περάσει από το σημείο $\mathrm{Α}$ και χαράσσουμε. Η ευθεία που προκύπτει είναι παράλληλη.

\begin{prosoxi}
\begin{itemize}[leftmargin=1.2em,itemsep=0.3ex,topsep=2pt]
\item Για να είναι παράλληλες δύο ευθείες δεν αρκεί να μην συναντιούνται στο σχέδιο: πρέπει να μην έχουν κοινό σημείο όσο κι αν προεκταθούν.
\item Δύο παράλληλα τμήματα δεν χρειάζεται να έχουν το ίδιο μήκος.
\item Στο επίπεδο δεν υπάρχουν δύο ευθείες που να μην είναι ούτε τεμνόμενες ούτε παράλληλες (στον χώρο υπάρχουν, οι ασύμβατες).
\item Αν μια ευθεία είναι παράλληλη σε δύο άλλες, αυτές είναι παράλληλες μεταξύ τους.
\end{itemize}
\end{prosoxi}

\tmima{Λυμένα παραδείγματα}

\paradeigma{Παράλληλες ή τεμνόμενες;}
\tixero Οι ευθείες: τα όρια του γηπέδου, οι ράγες, δύο μονοπάτια που διασταυρώνονται.\par
\tithelo Να πω ποιες είναι παράλληλες.\par
\possindeo Παράλληλες: οι ράγες και τα απέναντι όρια του γηπέδου. Τεμνόμενες: τα μονοπάτια.\par
\apantisi Ράγες και απέναντι πλευρές: παράλληλες. Μονοπάτια: τεμνόμενες.\par
\elegxos Οι παράλληλες δεν έχουν κοινό σημείο.

\paradeigma{Φέρνω παράλληλη}
\tixero Μια ευθεία $(\varepsilon)$ και ένα σημείο $\mathrm{Α}$ εκτός της.\par
\tithelo Να φέρω την παράλληλη.\par
\possindeo (1) Τοποθετώ τον χάρακα πάνω στην $(\varepsilon)$. (2) Τοποθετώ τον γνώμονα πάνω στον χάρακα, με μια πλευρά κατά μήκος του. (3) Σύρω τον γνώμονα μέχρι η άλλη του πλευρά να περάσει από το $\mathrm{Α}$. (4) Χαράσσω.\par
\apantisi Η ευθεία από το $\mathrm{Α}$ είναι η μοναδική παράλληλη προς την $(\varepsilon)$.\par
\elegxos Μετρώ δύο φορές την απόσταση των ευθειών: βγαίνει ίδια.

\paradeigma{Πόσα σημεία τομής;}
\tixero Τρεις διαφορετικές ευθείες του ίδιου επιπέδου.\par
\tithelo Πόσα το πολύ και πόσα το λιγότερο σημεία τομής μπορούν να έχουν.\par
\possindeo Αν είναι όλες παράλληλες, $0$. Αν δύο είναι παράλληλες και η τρίτη τέμνει και τις δύο: $2$. Αν περνούν και οι τρεις από το ίδιο σημείο: $1$. Αν ανά δύο τέμνονται σε διαφορετικά σημεία: $3$.\par
\apantisi $0$, $1$, $2$ ή $3$ σημεία.\par
\elegxos Δεν μπορεί να υπάρχει περισσότερα από $\frac{3\cdot2}{2}=3$ σημεία (ένα ανά ζευγάρι).

\paradeigma{Παράλληλα τμήματα}
\tixero Οι ευθείες $\mathrm{ΑΒ}$ και $\Gamma\Delta$ είναι παράλληλες, με $\mathrm{ΑΒ}=3$ και $\Gamma\Delta=7$.\par
\tithelo Τι λέμε για τα τμήματα.\par
\possindeo Δύο τμήματα σε παράλληλες ευθείες είναι παράλληλα.\par
\apantisi $\mathrm{ΑΒ}\parallel\Gamma\Delta$, παρότι έχουν διαφορετικά μήκη.\par
\elegxos Η παραλληλία εξαρτάται από τις ευθείες και όχι από τα μήκη.

\tmima{Ασκήσεις}

\zoni{Το κατάλαβα}{παράλληλες και τεμνόμενες}

\ex Ποιες από τις παρακάτω ευθείες είναι παράλληλες και ποιες τεμνόμενες: α) οι ράγες του τρένου \quad β) τα δύο μονοπάτια του κήπου που διασταυρώνονται \quad γ) οι αράδες του τετραδίου \quad δ) δύο ευθείες που έχουν ένα κοινό σημείο.

\ex Σχεδίασε μια ευθεία $(\varepsilon)$ και ένα σημείο $\mathrm{Α}$ εκτός της. Φέρε από το $\mathrm{Α}$ την παράλληλη με γνώμονα και χάρακα.

\ex Πόσες παράλληλες προς μια ευθεία μπορούν να φέρουμε από ένα σημείο εκτός της; Από ένα σημείο πάνω της;

\ex Πόσα σημεία τομής μπορεί να έχουν $3$ διαφορετικές ευθείες; Σχεδίασε κάθε περίπτωση.

\ex Πόσα το πολύ σημεία τομής έχουν $4$ ευθείες; Και $\nu$ ευθείες;

\ex Αν $\varepsilon_1\parallel\varepsilon_2$ και $\varepsilon_2\parallel\varepsilon_3$, τι ξέρεις για τις $\varepsilon_1$, $\varepsilon_3$;

\ex Σωστό ή λάθος; α) Δύο ευθείες του επιπέδου που δεν τέμνονται είναι παράλληλες. \quad β) Δύο παράλληλα τμήματα έχουν ίσα μήκη. \quad γ) Αν δύο ευθείες έχουν δύο κοινά σημεία, ταυτίζονται. \quad δ) Από σημείο εκτός ευθείας περνούν δύο παράλληλες προς αυτήν.

\zoni{Το χρησιμοποιώ}{κήπος, ράγες, γήπεδο}

\ex \textbf{Οι σειρές.} Οι σειρές των λαχανικών είναι παράλληλες ευθείες και απέχουν $60$ cm. Πώς θα τις σημάνεις στο έδαφος με σχοινί και παλούκια;

\ex \textbf{Το γήπεδο.} Πόσα ζεύγη παράλληλων πλευρών έχει ένα ορθογώνιο γήπεδο; Πόσα ζεύγη κάθετων πλευρών;

\ex \textbf{Τα μονοπάτια.} Πέντε μονοπάτια του κήπου είναι ευθείες και δεν είναι δύο παράλληλα. Πόσα σημεία τομής έχουν το πολύ;

\ex \textbf{Οι ράγες.} Γιατί οι ράγες του τρένου πρέπει να είναι παράλληλες;

\ex \textbf{Το μοτίβο.} Σχεδίασε τρεις παράλληλες ευθείες και μια ευθεία που τις τέμνει. Πόσα σημεία τομής έχεις;

\zoni{Το σκέφτομαι}{ο Ευκλείδης και η παραλληλία}

\ex Γιατί από δύο σημεία περνά μία ευθεία ενώ από σημείο εκτός ευθείας περνά μία παράλληλη; Σε τι διαφέρουν οι δύο προτάσεις;

\ex Αν δύο ευθείες δεν έχουν κοινό σημείο στο χαρτί σου, μπορεί να τέμνονται σε σημείο πέρα από το χαρτί; Εξήγησε.

\ex Δύο διαφορετικές ευθείες είναι κάθετες σε μια τρίτη. Τι σχέση έχουν μεταξύ τους;

\ex Μπορείς να φέρεις δύο διαφορετικές παράλληλες από ένα σημείο προς μια ευθεία; Τι θα συνέβαινε;

\zoni{Για το σπίτι}{ψάξε παράλληλες}

\ex Βρες στο σπίτι πέντε ζεύγη παράλληλων ευθειών και τρία ζεύγη τεμνόμενων.

\ex Σχεδίασε με γνώμονα και χάρακα τρεις παράλληλες ευθείες και μέτρησε την απόσταση ανάμεσά τους.

\paragrafos{Απόσταση σημείου από ευθεία και απόσταση παραλλήλων}{Β.1.10, «Απόσταση σημείου από ευθεία – Απόσταση παραλλήλων»}

\begin{anagke}
Στον κήπο υπάρχει ένα ευθύγραμμο μονοπάτι και ένα δέντρο που βρίσκεται λίγο πιο πέρα, εκτός μονοπατιού. Ο Νίκος θέλει να φτάσει από το δέντρο στο μονοπάτι όσο πιο γρήγορα γίνεται. Δοκιμάζει τρεις δρόμους με ένα σχοινί: έναν λοξό προς τα αριστερά, έναν λοξό προς τα δεξιά και έναν κατευθείαν «πάνω». Ο ευθύς είναι ο μικρότερος. «Ο δρόμος που πέφτει κάθετα στο μονοπάτι είναι ο συντομότερος!» ανακαλύπτει.

Η Αγγελική προσθέτει ένα δεύτερο σχέδιο: δύο παράλληλες σειρές με δέντρα. «Πόσο απέχουν οι δύο σειρές;» Μετράει από ένα δέντρο της μιας στην άλλη, κάθετα, και βρίσκει $2{,}4$ μέτρα. Μετράει από ένα άλλο δέντρο και βρίσκει πάλι $2{,}4$. «Σε όλα τα σημεία η απόσταση είναι ίδια!». Η Πολυξένη το συμπληρώνει: «Γι' αυτό οι ράγες του τρένου έχουν σταθερή απόσταση: αλλιώς το τρένο δεν χωράει».

Σε αυτή την παράγραφο ορίζουμε την απόσταση ενός σημείου από μια ευθεία και την απόσταση δύο παραλλήλων, δύο έννοιες που θα τις χρειαζόμαστε συνεχώς.
\end{anagke}

\begin{lexi}{απόσταση, κάθετο, πλάγιο}{απόσταση, κάθετο, πλάγιο}
Η \lexiword{απόσταση} σημείου από ευθεία είναι ο «συντομότερος δρόμος» από το σημείο προς την ευθεία. Το \lexiword{κάθετο} τμήμα πέφτει με ορθή γωνία πάνω στην ευθεία, ενώ τα υπόλοιπα, τα \lexiword{πλάγια}, είναι «λοξά» και μακρύτερα.
\end{lexi}

\begin{metafrasi}
Το κάθετο τμήμα από το $\mathrm{Α}$ στην $(\varepsilon)$. & $\mathrm{ΑΜ}\perp(\varepsilon)$\\
Η απόσταση του $\mathrm{Α}$ από την $(\varepsilon)$. & $\mathrm{ΑΜ}$\\
Η απόσταση δύο παραλλήλων. & η απόσταση οποιουδήποτε σημείου της μιας από την άλλη\\
Κάθε πλάγιο τμήμα. & μεγαλύτερο από το κάθετο\\
\end{metafrasi}

\tmima{Η θεωρία}

\begin{orismos}{απόσταση σημείου από ευθεία}
\textbf{Απόσταση} ενός σημείου $\mathrm{Α}$ από μια ευθεία $(\varepsilon)$ είναι το μήκος του \textbf{κάθετου} τμήματος $\mathrm{ΑΜ}$ από το $\mathrm{Α}$ προς την $(\varepsilon)$, όπου $\mathrm{Μ}$ το ίχνος της καθέτου.
\end{orismos}

Το κάθετο τμήμα είναι το \textbf{συντομότερο} από όλα τα τμήματα που ενώνουν το $\mathrm{Α}$ με σημεία της $(\varepsilon)$. Κάθε \textbf{πλάγιο} τμήμα είναι μεγαλύτερο από το κάθετο. Αν το σημείο ανήκει στην ευθεία, η απόστασή του από αυτήν είναι $0$.

\begin{center}
\begin{tikzpicture}[scale=1.1]
\draw[thick,{Stealth}-{Stealth}] (-2.5,0)--(3.5,0) node[right] {$(\varepsilon)$};
\fill (0.5,2) circle (2pt) node[above] {$\mathrm{Α}$}; \fill (0.5,0) circle (2pt) node[below] {$\mathrm{Μ}$};
\draw[gblue,thick] (0.5,2)--(0.5,0); \draw (0.5,0.25)--(0.75,0.25)--(0.75,0);
\draw (0.5,2)--(-1.5,0); \draw (0.5,2)--(2.5,0);
\node[gblue,right] at (0.5,1) {$\mathrm{ΑΜ}$};
\end{tikzpicture}
\end{center}

\tmimasmall{Απόσταση παραλλήλων}

Αν $\varepsilon_1\parallel\varepsilon_2$, όλα τα σημεία της $\varepsilon_1$ απέχουν την ίδια απόσταση από την $\varepsilon_2$. Η απόσταση αυτή λέγεται \textbf{απόσταση των παραλλήλων}. Δύο παράλληλες ευθείες έχουν σταθερή απόσταση, δηλαδή δεν πλησιάζουν ούτε απομακρύνονται.

Το σύνολο των σημείων που απέχουν απόσταση $\delta$ από μια ευθεία $(\varepsilon)$ είναι δύο παράλληλες ευθείες προς την $(\varepsilon)$, μία σε κάθε ημιεπίπεδο, που απέχουν $\delta$ από αυτήν.

\begin{prosoxi}
\begin{itemize}[leftmargin=1.2em,itemsep=0.3ex,topsep=2pt]
\item Η απόσταση σημείου από ευθεία μετριέται πάντα κάθετα, όχι λοξά.
\item Το πλάγιο τμήμα δεν δείχνει την απόσταση: είναι μεγαλύτερο.
\item Η απόσταση παραλλήλων είναι η ίδια όπου κι αν τη μετρήσεις, αλλά πρέπει να τη μετράς κάθετα.
\item Η απόσταση ενός σημείου της ευθείας από την ευθεία είναι $0$.
\end{itemize}
\end{prosoxi}

\tmima{Λυμένα παραδείγματα}

\paradeigma{Ο συντομότερος δρόμος}
\tixero Ένα δέντρο $\mathrm{Α}$ εκτός ευθύγραμμου μονοπατιού $(\varepsilon)$.\par
\tithelo Ποιος δρόμος από το δέντρο προς το μονοπάτι είναι ο συντομότερος.\par
\possindeo Η απόσταση του $\mathrm{Α}$ από την $(\varepsilon)$ είναι το κάθετο τμήμα $\mathrm{ΑΜ}$. Κάθε πλάγιο $\mathrm{ΑΒ}$ είναι μεγαλύτερο.\par
\apantisi Ο κάθετος δρόμος $\mathrm{ΑΜ}$.\par
\elegxos Με σχοινί: ο λοξός δρόμος βγαίνει μακρύτερος.

\paradeigma{Η απόσταση παραλλήλων}
\tixero Δύο σειρές δέντρων παράλληλες και $\mathrm{ΑΜ}\perp\varepsilon_2$ με $\mathrm{ΑΜ}=2{,}4$ μ.\par
\tithelo Η απόσταση ενός άλλου δέντρου $\mathrm{Β}$ της $\varepsilon_1$ από την $\varepsilon_2$.\par
\possindeo Οι παράλληλες έχουν σταθερή απόσταση.\par
\apantisi $2{,}4$ μ.\par
\elegxos Κάθε κάθετο τμήμα ανάμεσα στις παράλληλες είναι $2{,}4$ μ.

\paradeigma{Τρεις σειρές}
\tixero Τρεις παράλληλες σειρές δέντρων σε ίσες αποστάσεις $2{,}4$ μ.\par
\tithelo Η απόσταση της πρώτης από την τρίτη.\par
\possindeo Η απόσταση από την πρώτη στη δεύτερη είναι $2{,}4$ και από τη δεύτερη στην τρίτη $2{,}4$, αφού οι παράλληλες είναι διαδοχικές. $2{,}4+2{,}4=4{,}8$.\par
\apantisi $4{,}8$ μ.\par
\elegxos Η μεσαία είναι στη μέση των δύο άλλων.

\paradeigma{Σημεία με δοσμένη απόσταση}
\tixero Μια ευθεία $(\varepsilon)$.\par
\tithelo Ποια σημεία απέχουν $3$ cm από αυτήν.\par
\possindeo Σε κάθε ημιεπίπεδο υπάρχει μία παράλληλη προς την $(\varepsilon)$ σε απόσταση $3$ cm.\par
\apantisi Δύο παράλληλες ευθείες, μία σε κάθε ημιεπίπεδο.\par
\elegxos Η απόσταση ανάμεσά τους είναι $6$ cm.

\tmima{Ασκήσεις}

\zoni{Το κατάλαβα}{κάθετα και πλάγια}

\ex Σχεδίασε μια ευθεία $(\varepsilon)$ και ένα σημείο $\mathrm{Α}$ εκτός της. Φέρε την κάθετη από το $\mathrm{Α}$ και δύο πλάγια. Μέτρησε και σύγκρινε τα μήκη τους.

\ex Μέτρησε την απόσταση του $\mathrm{Α}$ από την $(\varepsilon)$. Πόση είναι η απόσταση ενός σημείου που βρίσκεται πάνω στην $(\varepsilon)$;

\ex Σχεδίασε δύο παράλληλες ευθείες και μέτρησε την απόστασή τους σε τρία διαφορετικά σημεία. Τι παρατηρείς;

\ex Τρεις παράλληλες ευθείες απέχουν διαδοχικά $2{,}4$ cm. Πόση είναι η απόσταση της πρώτης από την τρίτη;

\ex Σχεδίασε όλα τα σημεία που απέχουν $2$ cm από μια ευθεία $(\varepsilon)$.

\ex Σωστό ή λάθος; α) Το κάθετο τμήμα είναι το μικρότερο από όλα όσα ενώνουν ένα σημείο με την ευθεία. \quad β) Η απόσταση παραλλήλων δεν είναι σταθερή. \quad γ) Η απόσταση σημείου της ευθείας από την ευθεία είναι $0$. \quad δ) Το πλάγιο είναι μικρότερο από το κάθετο.

\zoni{Το χρησιμοποιώ}{κήπος, δρόμοι, σχοινιά}

\ex \textbf{Το κοντινότερο.} Ένα δέντρο απέχει $3$ μ. κάθετα από ένα μονοπάτι. Ο Νίκος δοκιμάζει να φτάσει λοξά και βρίσκει $3{,}6$ μ. Μπορεί να έχει υπολογίσει τον κάθετο δρόμο λάθος; Εξήγησε.

\ex \textbf{Οι σειρές.} Δύο παράλληλες σειρές λαχανικών απέχουν $60$ cm. Μια τρίτη σειρά, παράλληλη προς αυτές και προς την άλλη πλευρά της δεύτερης, απέχει $60$ cm από τη δεύτερη. Πόσο απέχει από την πρώτη;

\ex \textbf{Το πλάγιο.} Ένα σημείο $\mathrm{Α}$ απέχει $4$ cm από μια ευθεία $(\varepsilon)$. Ένα σημείο $\mathrm{Β}$ της $(\varepsilon)$ μπορεί να απέχει από το $\mathrm{Α}$ $3$ cm; $4$ cm; $6$ cm;

\ex \textbf{Το ράφι.} Ένα ράφι είναι παράλληλο προς το πάτωμα και απέχει $1{,}2$ μ. Πόσο απέχει κάθε σημείο του από το πάτωμα;

\ex \textbf{Το χωράφι.} Ένα ορθογώνιο χωράφι έχει διαστάσεις $40$ και $25$ μ. Πόση είναι η απόσταση ανάμεσα στις μεγάλες απέναντι πλευρές του; Και ανάμεσα στις μικρές;

\zoni{Το σκέφτομαι}{ο συντομότερος δρόμος}

\ex Εξήγησε γιατί το κάθετο τμήμα είναι το συντομότερο (σκέψου ένα τεντωμένο νήμα και μια ράβδο που στρέφεται).

\ex Σε ποια περίπτωση η απόσταση δύο παραλλήλων είναι $0$;

\ex Βρες τα σημεία που απέχουν $3$ cm από την $(\varepsilon)$ και $5$ cm από μια παράλληλη προς αυτήν σε απόσταση $8$ cm.

\ex Δύο παράλληλες ευθείες απέχουν $6$ cm. Πόσες ευθείες παράλληλες προς αυτές απέχουν $3$ cm και από τις δύο;

\zoni{Για το σπίτι}{μέτρησε κάθετα}

\ex Μέτρησε την απόσταση ανάμεσα σε δύο παράλληλες ακμές του τραπεζιού σου σε δύο διαφορετικά σημεία.

\ex Βρες τη μικρότερη απόσταση ενός σημείου του δωματίου από έναν τοίχο και εξήγησε τι μέτρησες.

\paragrafos{Ο κύκλος και τα στοιχεία του}{Β.1.11, «Κύκλος και στοιχεία του κύκλου»}

\begin{anagke}
Το πιο όμορφο σημείο του κήπου θα είναι ένα στρογγυλό παρτέρι με ένα συντριβάνι στο κέντρο. Ο παππούς Γιάννης δείχνει πώς θα το χαράξουν στο χώμα: καρφώνει ένα παλούκι στο κέντρο, δένει ένα σχοινί και, κρατώντας το τεντωμένο, γυρίζει και χαράσσει τη γραμμή. «Όλα τα σημεία της γραμμής έχουν την ίδια απόσταση από το παλούκι, όσο είναι το σχοινί», λέει.

Η Αγγελική το αναγνωρίζει: «Αυτό είναι ο κύκλος! Και το σχοινί είναι η ακτίνα!». Ο Νίκος ρωτάει πώς λέγεται το πιο μακρύ τμήμα που μπορούν να τεντώσουν από ένα άκρο του κύκλου ως το άλλο: «Η διάμετρος», απαντά η Πολυξένη, «περνάει από το κέντρο και είναι διπλάσια από την ακτίνα». Ο Βασίλης στέκεται κάπου και ρωτάει αν είναι «μέσα» ή «έξω» από τον κύκλο: «Αν η απόστασή σου από το παλούκι είναι μικρότερη από το σχοινί, είσαι μέσα».

Ο κύκλος είναι το πιο συμμετρικό σχήμα και όλα τα στοιχεία του προκύπτουν από τη μέτρηση αποστάσεων από το κέντρο.
\end{anagke}

\begin{lexi}{κύκλος, ακτίνα, διάμετρος, χορδή}{κύκλος, ακτίνα, διάμετρος, χορδή}
Ο \lexiword{κύκλος} είναι ό,τι «κυλάει», γυρίζει. Η \lexiword{ακτίνα} είναι η «αχτίδα», η ακτίνα του φωτός που ξεκινά από το κέντρο. Η \lexiword{διάμετρος} είναι το «δια-μέτρο», αυτό που μετράει «δια μέσου» του κέντρου. Η \lexiword{χορδή} είναι η «χορδή» του μουσικού οργάνου: ένα τεντωμένο νήμα ανάμεσα σε δύο σημεία. Το \lexiword{ημικύκλιο} είναι το μισό του κύκλου.
\end{lexi}

\begin{metafrasi}
Ο κύκλος με κέντρο $\mathrm{Ο}$ και ακτίνα $\rho$. & $(\mathrm{Ο},\rho)$\\
Τα σημεία πάνω στον κύκλο. & $\mathrm{ΟΜ}=\rho$\\
Τα σημεία μέσα στον κύκλο. & $\mathrm{ΟΜ}<\rho$\\
Τα σημεία έξω από τον κύκλο. & $\mathrm{ΟΜ}>\rho$\\
Η διάμετρος. & $\delta=2\rho$\\
\end{metafrasi}

\tmima{Η θεωρία}

\begin{orismos}{κύκλος}
\textbf{Κύκλος} με κέντρο $\mathrm{Ο}$ και ακτίνα $\rho$ είναι το σύνολο όλων των σημείων του επιπέδου που απέχουν από το $\mathrm{Ο}$ απόσταση ίση με $\rho$. Συμβολίζεται $(\mathrm{Ο},\rho)$. Η απόσταση $\rho$ λέγεται \textbf{ακτίνα} και το $\mathrm{Ο}$ \textbf{κέντρο}.
\end{orismos}

Για να σχεδιάσουμε κύκλο χρησιμοποιούμε \textbf{διαβήτη}. Δύο κύκλοι με ίσες ακτίνες είναι ίσοι.

\begin{center}
\begin{tikzpicture}[scale=1.15]
\draw[thick] (0,0) circle (1.6);
\fill (0,0) circle (2pt) node[below left] {$\mathrm{Ο}$};
\draw[gblue,thick] (0,0)--(1.13,1.13) node[midway,above left,gblue] {$\rho$}; \fill (1.13,1.13) circle (2pt) node[above right] {$\mathrm{Α}$};
\draw (-1.6,0)--(1.6,0); \fill (-1.6,0) circle (2pt) node[left] {$\mathrm{Β}$}; \fill (1.6,0) circle (2pt) node[right] {$\mathrm{Γ}$};
\node[below] at (0,-0.1) {\small διάμετρος $\mathrm{ΒΓ}$};
\draw[dashed] (-1.2,1.06)--(0.8,-1.38); \fill (-1.2,1.06) circle (2pt); \fill (0.8,-1.38) circle (2pt);
\node at (-1.2,-0.9) {\small χορδή};
\end{tikzpicture}
\end{center}

\begin{itemize}[leftmargin=1.2em,itemsep=0.3ex,topsep=2pt]
\item \textbf{Χορδή} είναι το ευθύγραμμο τμήμα που ενώνει δύο σημεία του κύκλου.
\item \textbf{Διάμετρος} είναι η χορδή που περνά από το κέντρο. Είναι η μεγαλύτερη χορδή, ίση με $\delta=2\rho$, και χωρίζει τον κύκλο σε δύο ίσα μέρη, τα \textbf{ημικύκλια}.
\item Ένα σημείο $\mathrm{Μ}$ είναι \textbf{πάνω} στον κύκλο αν $\mathrm{ΟΜ}=\rho$, \textbf{μέσα} σε αυτόν αν $\mathrm{ΟΜ}<\rho$ και \textbf{έξω} αν $\mathrm{ΟΜ}>\rho$.
\item Ο κύκλος μαζί με το εσωτερικό του λέγεται \textbf{κυκλικός δίσκος}.
\end{itemize}

\begin{prosoxi}
\begin{itemize}[leftmargin=1.2em,itemsep=0.3ex,topsep=2pt]
\item Ο κύκλος είναι μόνο η γραμμή, όχι το εσωτερικό: το εσωτερικό μαζί με τη γραμμή είναι ο κυκλικός δίσκος.
\item Η ακτίνα είναι τμήμα (ή μήκος), όχι ολόκληρη ευθεία.
\item Η διάμετρος είναι πάντα $2\rho$, και κάθε διάμετρος είναι χορδή, αλλά όχι κάθε χορδή είναι διάμετρος.
\item Σε έναν κύκλο όλες οι ακτίνες είναι ίσες, όλες οι διάμετροι είναι ίσες.
\end{itemize}
\end{prosoxi}

\tmima{Λυμένα παραδείγματα}

\paradeigma{Το κυκλικό παρτέρι}
\tixero Το σχοινί του παππού είναι $2{,}5$ μ.\par
\tithelo Η ακτίνα και η διάμετρος του παρτεριού.\par
\possindeo Το σχοινί είναι η ακτίνα: $\rho=2{,}5$ μ. Η διάμετρος $\delta=2\cdot2{,}5=5$ μ.\par
\apantisi $\rho=2{,}5$ μ. και $\delta=5$ μ.\par
\elegxos Το παρτέρι έχει πλάτος $5$ μ.

\paradeigma{Μέσα, πάνω ή έξω;}
\tixero Κύκλος $(\mathrm{Ο},4)$ και σημεία με $\mathrm{ΟΑ}=3$, $\mathrm{ΟΒ}=4$, $\mathrm{ΟΓ}=5$.\par
\tithelo Η θέση κάθε σημείου.\par
\possindeo Συγκρίνω την απόσταση με την ακτίνα $4$. $\mathrm{ΟΑ}<4$: μέσα. $\mathrm{ΟΒ}=4$: πάνω. $\mathrm{ΟΓ}>4$: έξω.\par
\apantisi Το $\mathrm{Α}$ μέσα, το $\mathrm{Β}$ πάνω, το $\mathrm{Γ}$ έξω.\par
\elegxos Η απόσταση από το κέντρο καθορίζει τη θέση.

\paradeigma{Μία χορδή}
\tixero Ο κύκλος $(\mathrm{Ο},5)$ και μια χορδή $\mathrm{ΑΒ}$.\par
\tithelo Πόσο μεγάλη μπορεί να είναι η χορδή.\par
\possindeo Η μεγαλύτερη χορδή είναι η διάμετρος, $2\cdot5=10$.\par
\apantisi $\mathrm{ΑΒ}\le10$, ίσο με $10$ μόνο όταν περνά από το κέντρο.\par
\elegxos Αν $\mathrm{ΑΒ}=10$ τότε το κέντρο είναι το μέσο του $\mathrm{ΑΒ}$.

\paradeigma{Ίσοι κύκλοι}
\tixero Δύο κύκλοι: $(\mathrm{Ο},3)$ και $(\mathrm{Κ},3)$.\par
\tithelo Είναι ίσοι;\par
\possindeo Έχουν ίσες ακτίνες, άρα μπορούν να εφαρμοστούν ο ένας πάνω στον άλλο μετακινώντας το κέντρο.\par
\apantisi Ναι.\par
\elegxos Ίσες ακτίνες δίνουν ίσους κύκλους.

\tmima{Ασκήσεις}

\zoni{Το κατάλαβα}{ακτίνα, διάμετρος, χορδή}

\ex Σχεδίασε με διαβήτη κύκλο $(\mathrm{Ο},3\ \mathrm{cm})$. Φέρε μια ακτίνα, μια διάμετρο και μια χορδή που δεν περνά από το κέντρο.

\ex Συμπλήρωσε: α) $\rho=4$, $\delta=\ldots$ \quad β) $\delta=9$, $\rho=\ldots$ \quad γ) $\rho=2{,}5$, $\delta=\ldots$ \quad δ) $\delta=7{,}2$, $\rho=\ldots$.

\ex Κύκλος $(\mathrm{Ο},4)$. Πες αν είναι μέσα, πάνω ή έξω τα σημεία με $\mathrm{ΟΑ}=3$, $\mathrm{ΟΒ}=4$, $\mathrm{ΟΓ}=5$, $\mathrm{ΟΔ}=0$.

\ex Στον κύκλο $(\mathrm{Ο},5)$ πόσο μεγάλη είναι η μεγαλύτερη χορδή; Πόσες διάμετροι υπάρχουν;

\ex Σχεδίασε δύο κύκλους με την ίδια ακτίνα και δείξε ότι είναι ίσοι με διαφανές χαρτί.

\ex Σωστό ή λάθος; α) Η διάμετρος είναι χορδή. \quad β) Η μεγαλύτερη χορδή είναι η ακτίνα. \quad γ) $\delta=2\rho$. \quad δ) Ένα σημείο με $\mathrm{ΟΜ}<\rho$ είναι πάνω στον κύκλο. \quad ε) Όλες οι ακτίνες ενός κύκλου είναι ίσες.

\zoni{Το χρησιμοποιώ}{παρτέρια, ρολόγια, ρόδες}

\ex \textbf{Το παρτέρι.} Το κυκλικό παρτέρι χαράσσεται με σχοινί $2{,}5$ μ. Πόσο είναι το πλάτος του;

\ex \textbf{Το ποδήλατο.} Ο τροχός ενός ποδηλάτου έχει διάμετρο $66$ cm. Πόση είναι η ακτίνα του; Σε ύψος πόσο βρίσκεται ο άξονάς του από το έδαφος;

\ex \textbf{Το ρολόι.} Ο λεπτοδείκτης έχει μήκος $8$ cm. Σε ποιο κύκλο κινείται η άκρη του και πόση διάμετρο έχει;

\ex \textbf{Η σέντρα.} Ο κύκλος της σέντρας σε ένα γήπεδο ποδοσφαίρου έχει ακτίνα $9{,}15$ μ. Πόση είναι η διάμετρός του;

\ex \textbf{Η στρογγυλή πίτα.} Μια πίτα έχει διάμετρο $26$ cm. Θα χωρέσει σε ταψί διαστάσεων $25\times25$ cm; Εξήγησε.

\zoni{Το σκέφτομαι}{όλα ίσα από το κέντρο}

\ex Γιατί η διάμετρος είναι η μεγαλύτερη χορδή; (Σκέψου το τμήμα ανάμεσα σε δύο σημεία και την απόστασή τους από το κέντρο.)

\ex Αν σε έναν κύκλο οι ακτίνες είναι $5$ cm, πόσα σημεία του κύκλου απέχουν $5$ cm από ένα δοσμένο σημείο του κύκλου; Πρόβλεψε και σχεδίασε.

\ex Πόσοι κύκλοι μπορούν να περάσουν από δύο δοσμένα σημεία $\mathrm{Α}$ και $\mathrm{Β}$; Πού βρίσκονται τα κέντρα τους;

\ex Μπορούν δύο διαφορετικοί κύκλοι με το ίδιο κέντρο να έχουν κοινό σημείο; Εξήγησε.

\zoni{Για το σπίτι}{στρογγυλά αντικείμενα}

\ex Βρες πέντε κυκλικά αντικείμενα στο σπίτι και μέτρησε διάμετρο και ακτίνα.

\ex Σχεδίασε έναν κύκλο με νήμα και καρφίτσα στο χαρτόνι.

\paragrafos{Επίκεντρη γωνία, τόξο και μέτρηση τόξου}{Β.1.12, «Επίκεντρη γωνία – Σχέση επίκεντρης γωνίας και του αντίστοιχου τόξου – Μέτρηση τόξου»}

\begin{anagke}
Στο πάρκο δίπλα στο σχολείο υπάρχει μια ρόδα λούνα παρκ με $8$ ίσες καμπίνες. Η παρέα την κοιτάζει από ψηλά: οι καμπίνες είναι στερεωμένες σε ίσες αποστάσεις πάνω σε ένα κύκλο. «Πόσες μοίρες γυρίζει ο τροχός για να πάει η μία καμπίνα στη θέση της επόμενης;» ρωτάει ο Ηλίας. Η Πολυξένη σκέφτεται: «Ο πλήρης κύκλος είναι $360^\circ$ και υπάρχουν $8$ ίσα διαστήματα: $360:8=45^\circ$».

Η Αγγελική σχεδιάζει δύο ακτίνες προς δύο διαδοχικές καμπίνες και παρατηρεί τη γωνία ανάμεσά τους με κορυφή το κέντρο: «Είναι $45^\circ$, και το κομμάτι του κύκλου ανάμεσα στις δύο καμπίνες το μετράμε και αυτό $45^\circ$». Ο Στέλιος ρωτάει: «Το μήκος του κύκλου και οι μοίρες είναι το ίδιο;». Η κυρία Αλεξάνδρα απαντάει: «Όχι. Άλλο το μήκος σε μέτρα και άλλο το άνοιγμα σε μοίρες. Σήμερα μιλάμε για το άνοιγμα, το λεγόμενο \emph{μέτρο του τόξου}».

Στην παράγραφο αυτή ορίζουμε την επίκεντρη γωνία και το αντίστοιχο τόξο, και μετράμε τα τόξα με μοίρες.
\end{anagke}

\begin{lexi}{επίκεντρη γωνία, τόξο}{επίκεντρη γωνία, τόξο}
Η \lexiword{επίκεντρη} γωνία είναι η γωνία που έχει την κορυφή της «επί του κέντρου» του κύκλου. Το \lexiword{τόξο} είναι το κομμάτι του κύκλου, σαν το τόξο του τοξοβόλου, καμπυλωμένο. Το \lexiword{ημικύκλιο} είναι το μισό του κύκλου.
\end{lexi}

\begin{metafrasi}
Γωνία με κορυφή το κέντρο. & επίκεντρη γωνία $\widehat{\mathrm{ΑΟΒ}}$\\
Το κομμάτι του κύκλου μέσα στη γωνία. & αντίστοιχο τόξο $\overset{\frown}{\mathrm{ΑΒ}}$\\
Το μέτρο του τόξου. & ίσο με το μέτρο της επίκεντρης γωνίας\\
Ο πλήρης κύκλος. & $360^\circ$\\
Το ημικύκλιο. & $180^\circ$\\
\end{metafrasi}

\tmima{Η θεωρία}

\begin{orismos}{επίκεντρη γωνία και τόξο}
\textbf{Επίκεντρη} γωνία ενός κύκλου $(\mathrm{Ο},\rho)$ είναι κάθε γωνία $x\mathrm{Ο}y$ με κορυφή το κέντρο $\mathrm{Ο}$. Αν οι πλευρές της τέμνουν τον κύκλο στα $\mathrm{Α}$ και $\mathrm{Β}$, το τμήμα του κύκλου που βρίσκεται στο εσωτερικό της γωνίας λέγεται \textbf{αντίστοιχο τόξο} της και γράφεται $\overset{\frown}{\mathrm{ΑΒ}}$.
\end{orismos}

\begin{center}
\begin{tikzpicture}[scale=1.2]
\draw (0,0) circle (1.6);
\fill (0,0) circle (2pt) node[below left] {$\mathrm{Ο}$};
\draw[thick] (0,0)--(1.6,0) (0,0)--({1.6*cos(50)},{1.6*sin(50)});
\fill (1.6,0) circle (2pt) node[right] {$\mathrm{Α}$}; \fill ({1.6*cos(50)},{1.6*sin(50)}) circle (2pt) node[above right] {$\mathrm{Β}$};
\draw[gblue,very thick] (1.6,0) arc(0:50:1.6);
\draw[gblue] (0.5,0) arc(0:50:0.5); \node[gblue] at (0.8,0.2) {\small$\omega$};
\end{tikzpicture}
\end{center}

\begin{kanonas}{μέτρο τόξου}
Το \textbf{μέτρο ενός τόξου} σε μοίρες είναι ίσο με το μέτρο της αντίστοιχης επίκεντρης γωνίας. Ολόκληρος ο κύκλος είναι τόξο $360^\circ$ και το ημικύκλιο $180^\circ$. Ίσες επίκεντρες γωνίες αντιστοιχούν σε ίσα τόξα, και αντίστροφα.
\end{kanonas}

Έτσι το τόξο $\overset{\frown}{\mathrm{ΑΒ}}$ του σχήματος είναι $50^\circ$. Όταν η γωνία είναι μη κυρτή, το τόξο που αντιστοιχεί σε αυτήν είναι το μεγάλο: $360^\circ-50^\circ=310^\circ$.

Το μέτρο τόξου (σε μοίρες) είναι διαφορετικό από το \emph{μήκος} του τόξου (σε μέτρα): δύο ομόκεντροι κύκλοι έχουν τόξα με ίδιο μέτρο αλλά διαφορετικό μήκος.

\begin{prosoxi}
\begin{itemize}[leftmargin=1.2em,itemsep=0.3ex,topsep=2pt]
\item Το μέτρο του τόξου σε μοίρες δεν είναι το μήκος του σε μέτρα.
\item Κάθε χορδή $\mathrm{ΑΒ}$ ορίζει δύο τόξα (ένα μικρό και ένα μεγάλο), εκτός αν είναι διάμετρος (δύο ημικύκλια).
\item Η επίκεντρη γωνία έχει την κορυφή της στο κέντρο: αν η κορυφή είναι αλλού, δεν είναι επίκεντρη.
\item Σε διαφορετικούς κύκλους, ίσα τόξα (σε μοίρες) δεν σημαίνει ίσα μήκη.
\end{itemize}
\end{prosoxi}

\tmima{Λυμένα παραδείγματα}

\paradeigma{Η ρόδα με τις $8$ καμπίνες}
\tixero Οκτώ ίσες καμπίνες πάνω σε κύκλο.\par
\tithelo Το μέτρο του τόξου ανάμεσα σε δύο διαδοχικές.\par
\possindeo Ο κύκλος είναι $360^\circ$ και τα $8$ τόξα είναι ίσα: $360^\circ:8=45^\circ$.\par
\apantisi $45^\circ$.\par
\elegxos $8\cdot45=360$.

\paradeigma{Τόξο και επίκεντρη}
\tixero Η επίκεντρη γωνία $\widehat{\mathrm{ΑΟΒ}}=50^\circ$.\par
\tithelo Το μικρό και το μεγάλο τόξο $\overset{\frown}{\mathrm{ΑΒ}}$.\par
\possindeo Το μικρό τόξο είναι $50^\circ$. Το μεγάλο είναι $360^\circ-50^\circ=310^\circ$.\par
\apantisi $50^\circ$ και $310^\circ$.\par
\elegxos $50+310=360$.

\paradeigma{Διάμετρος και τόξα}
\tixero Η $\mathrm{ΑΒ}$ είναι διάμετρος του κύκλου.\par
\tithelo Το μέτρο των τόξων που ορίζει.\par
\possindeo Η διάμετρος χωρίζει τον κύκλο σε δύο ημικύκλια. Κάθε ημικύκλιο είναι $180^\circ$.\par
\apantisi $180^\circ$ το καθένα.\par
\elegxos Η επίκεντρη γωνία $\widehat{\mathrm{ΑΟΒ}}$ είναι ευθεία.

\paradeigma{Ένα ρολόι}
\tixero Ο λεπτοδείκτης κινείται από το $12$ στο $4$.\par
\tithelo Το τόξο που διέγραψε η άκρη του.\par
\possindeo Κάθε λεπτό ($1/60$ του κύκλου) είναι $6^\circ$, ή κάθε πεντάλεπτο $30^\circ$. Από το $12$ στο $4$ είναι $4$ πεντάλεπτα, $4\cdot30^\circ=120^\circ$.\par
\apantisi $120^\circ$.\par
\elegxos Στα $20$ λεπτά, $20\cdot6^\circ=120^\circ$.

\tmima{Ασκήσεις}

\zoni{Το κατάλαβα}{τόξα και μοίρες}

\ex Σχεδίασε κύκλο και επίκεντρη γωνία $70^\circ$. Ποιο είναι το μέτρο του μικρού και του μεγάλου τόξου που αντιστοιχούν;

\ex Ένας κύκλος χωρίζεται σε $5$ ίσα τόξα. Πόσες μοίρες είναι το καθένα; Και σε $9$; Σε $12$;

\ex Πόσες μοίρες είναι το $\frac{1}{4}$, το $\frac{1}{8}$, το $\frac{3}{4}$ και το $\frac{5}{6}$ του κύκλου;

\ex Ένα τόξο έχει μέτρο $60^\circ$. Τι μέρος του κύκλου είναι;

\ex Μια διάμετρος $\mathrm{ΑΒ}$ ορίζει δύο τόξα. Πόσες μοίρες είναι το καθένα;

\ex Σχεδίασε κύκλο και τρεις επίκεντρες διαδοχικές γωνίες $100^\circ$, $120^\circ$, $140^\circ$. Πόσο είναι τα τόξα και πόσο το άθροισμά τους;

\ex Σωστό ή λάθος; α) Το μέτρο του τόξου είναι ίσο με της επίκεντρης γωνίας. \quad β) Το ημικύκλιο είναι $180^\circ$. \quad γ) Ίσα τόξα έχουν πάντα ίσα μήκη σε οποιουσδήποτε κύκλους. \quad δ) Ο κύκλος είναι τόξο $360^\circ$.

\zoni{Το χρησιμοποιώ}{ρόδες, ρολόγια, πίτες}

\ex \textbf{Η ρόδα.} Μια ρόδα λούνα παρκ έχει $12$ ίσες καμπίνες. Πόσες μοίρες γυρίζει για να έρθει η επόμενη καμπίνα στη θέση της;

\ex \textbf{Το ρολόι.} Πόσες μοίρες διαγράφει η άκρη του λεπτοδείκτη σε $20$ λεπτά; Σε $45$ λεπτά; Σε μία ώρα;

\ex \textbf{Η πίτα.} Μια πίτα κόβεται σε φέτες με επίκεντρες γωνίες $45^\circ$. Πόσες φέτες βγαίνουν; Τι μέρος της πίτας είναι $3$ φέτες;

\ex \textbf{Το ραντάρ.} Ένα ραντάρ περιστρέφεται και σαρώνει ένα τόξο $75^\circ$. Τι μέρος του κύκλου καλύπτει;

\ex \textbf{Ο ανεμόμυλος.} Ένας ανεμόμυλος με $4$ πτερύγια γυρίζει $90^\circ$. Πόσα τέτοια τόξα χρειάζονται για μία πλήρη περιστροφή;

\zoni{Το σκέφτομαι}{μέτρο ή μήκος;}

\ex Δύο κύκλοι έχουν ίδιο κέντρο και διαφορετικές ακτίνες. Οι ίδιες δύο ακτίνες ορίζουν στον καθένα ένα τόξο. Είναι ίσα τα μέτρα τους; Είναι ίσα τα μήκη τους;

\ex Γιατί ένας πλήρης κύκλος είναι $360^\circ$ και όχι $100^\circ$; Θα άλλαζε τίποτα αν το διαλέγαμε $100$;

\ex Μια χορδή $\mathrm{ΑΒ}$ ορίζει τόξο $120^\circ$. Πόσο είναι το άλλο τόξο της;

\ex Αν η επίκεντρη γωνία διπλασιαστεί, τι γίνεται το τόξο;

\zoni{Για το σπίτι}{ο κύκλος και οι μοίρες}

\ex Σχεδίασε έναν κύκλο και χώρισέ τον σε $6$ ίσα τόξα με διαβήτη και μοιρογνωμόνιο. Πόσες μοίρες είναι το καθένα;

\ex Βρες κάτι στο σπίτι (ρολόι, πίτα, ρόδα) που χωρίζεται σε ίσα τόξα και υπολόγισε τις μοίρες κάθε τόξου.

\paragrafos{Θέσεις ευθείας και κύκλου: εξωτερική, εφαπτομένη, τέμνουσα}{Β.1.13, «Θέσεις ευθείας και κύκλου»}

\begin{anagke}
Γύρω από το στρογγυλό παρτέρι θα περάσει ένα μονοπάτι, και η Α3 συζητάει πού θα το τοποθετήσει. Ο Νίκος προτείνει να περνά ακριβώς πάνω από το παρτέρι. «Θα πατάμε τα λουλούδια!», διαμαρτύρεται η Ρήνη. Ο Βασίλης προτείνει ένα μονοπάτι που περνά μακριά από το παρτέρι, χωρίς να το αγγίζει. «Τότε δεν θα το βλέπουμε από κοντά», λέει η Αγγελική. Η Πολυξένη προτείνει το τρίτο: ένα μονοπάτι που ακουμπά το παρτέρι σε ένα μόνο σημείο, σαν τον δρόμο που ακουμπά μια ρόδα.

Η κυρία Αλεξάνδρα το ονομάζει: τρεις θέσεις μιας ευθείας ως προς έναν κύκλο. Δείχνει πώς καταλαβαίνουμε ποια είναι από την απόσταση του κέντρου από την ευθεία: «Μετρήστε την απόσταση του κέντρου από το μονοπάτι και συγκρίνετέ την με την ακτίνα».

Στο τέλος του κεφαλαίου δένουμε την απόσταση σημείου από ευθεία (1.10) με τον κύκλο (1.11): τρεις θέσεις ευθείας και κύκλου, και τι ισχύει όταν η ευθεία εφάπτεται.
\end{anagke}

\begin{lexi}{εφαπτομένη, τέμνουσα}{εφαπτομένη, τέμνουσα, εξωτερική}
Η \lexiword{εφαπτομένη} είναι αυτή που «εφάπτεται», δηλαδή «ακουμπά», μόνο σε ένα σημείο. Η \lexiword{τέμνουσα} είναι αυτή που «τέμνει», περνά μέσα από τον κύκλο και τον τέμνει σε δύο σημεία. Η \lexiword{εξωτερική} ευθεία βρίσκεται «έξω» από τον κύκλο και δεν τον ακουμπά.
\end{lexi}

\begin{metafrasi}
Η απόσταση του κέντρου από την ευθεία. & $\mathrm{ΟΜ}=d$\\
Η ευθεία είναι μακριά από τον κύκλο. & $d>\rho$: εξωτερική\\
Η ευθεία ακουμπά τον κύκλο σε ένα σημείο. & $d=\rho$: εφαπτομένη\\
Η ευθεία περνά μέσα στον κύκλο. & $d<\rho$: τέμνουσα\\
\end{metafrasi}

\tmima{Η θεωρία}

Έστω κύκλος $(\mathrm{Ο},\rho)$ και ευθεία $(\varepsilon)$ και $d=\mathrm{ΟΜ}$ η απόσταση του κέντρου από την ευθεία (με $\mathrm{ΟΜ}\perp\varepsilon$). Υπάρχουν τρεις περιπτώσεις:

\begin{center}
\begin{tikzpicture}[scale=0.85]
\begin{scope}
\draw (0,0) circle (1.2); \fill (0,0) circle (1.5pt) node[above left,font=\small] {$\mathrm{Ο}$};
\draw[thick,{Stealth}-{Stealth}] (-1.8,1.9)--(1.8,1.9); \draw[gblue] (0,0)--(0,1.9);
\node at (0,-1.7) {\small $d>\rho$: εξωτερική};
\end{scope}
\begin{scope}[xshift=4.5cm]
\draw (0,0) circle (1.2); \fill (0,0) circle (1.5pt);
\draw[thick,{Stealth}-{Stealth}] (-1.8,1.2)--(1.8,1.2); \draw[gblue] (0,0)--(0,1.2); \fill (0,1.2) circle (2pt) node[above right,font=\small] {$\mathrm{Μ}$};
\node at (0,-1.7) {\small $d=\rho$: εφαπτομένη};
\end{scope}
\begin{scope}[xshift=9cm]
\draw (0,0) circle (1.2); \fill (0,0) circle (1.5pt);
\draw[thick,{Stealth}-{Stealth}] (-1.8,0.6)--(1.8,0.6); \draw[gblue] (0,0)--(0,0.6);
\fill (-1.04,0.6) circle (2pt) node[above left,font=\small] {$\mathrm{Α}$}; \fill (1.04,0.6) circle (2pt) node[above right,font=\small] {$\mathrm{Β}$};
\node at (0,-1.7) {\small $d<\rho$: τέμνουσα};
\end{scope}
\end{tikzpicture}
\end{center}

\begin{itemize}[leftmargin=1.2em,itemsep=0.3ex,topsep=2pt]
\item Αν $d>\rho$, η ευθεία δεν έχει κανένα κοινό σημείο με τον κύκλο: είναι \textbf{εξωτερική} του κύκλου.
\item Αν $d=\rho$, η ευθεία έχει \textbf{ένα} κοινό σημείο $\mathrm{Μ}$ με τον κύκλο και λέγεται \textbf{εφαπτομένη} του κύκλου στο $\mathrm{Μ}$. Το $\mathrm{Μ}$ λέγεται σημείο επαφής και η εφαπτομένη είναι \textbf{κάθετη} στην ακτίνα $\mathrm{ΟΜ}$.
\item Αν $d<\rho$, η ευθεία έχει \textbf{δύο} κοινά σημεία $\mathrm{Α}$, $\mathrm{Β}$ με τον κύκλο και λέγεται \textbf{τέμνουσα}.
\end{itemize}

Αντίστροφα, από το πλήθος των κοινών σημείων προκύπτει η θέση: $0$ κοινά σημεία σημαίνει εξωτερική, $1$ εφαπτομένη, $2$ τέμνουσα. Μια ευθεία δεν μπορεί να έχει με κύκλο περισσότερα από δύο κοινά σημεία.

\begin{prosoxi}
\begin{itemize}[leftmargin=1.2em,itemsep=0.3ex,topsep=2pt]
\item Η απόσταση του κέντρου από την ευθεία είναι η κάθετη απόσταση, όχι οποιοδήποτε πλάγιο τμήμα.
\item Η εφαπτομένη έχει ένα μόνο κοινό σημείο με τον κύκλο. Αν έχει δύο, είναι τέμνουσα.
\item Η εφαπτομένη είναι κάθετη στην ακτίνα του σημείου επαφής, όχι σε οποιαδήποτε ακτίνα.
\item Σε ένα σημείο του κύκλου υπάρχει μία και μοναδική εφαπτομένη.
\end{itemize}
\end{prosoxi}

\tmima{Λυμένα παραδείγματα}

\paradeigma{Η θέση του μονοπατιού}
\tixero Κύκλος $(\mathrm{Ο},3)$ και ευθείες με αποστάσεις $d$ από το κέντρο: $5$, $3$, $1{,}5$.\par
\tithelo Η θέση καθεμίας.\par
\possindeo Συγκρίνω το $d$ με το $\rho=3$. $5>3$: εξωτερική. $3=3$: εφαπτομένη. $1{,}5<3$: τέμνουσα.\par
\apantisi Εξωτερική, εφαπτομένη, τέμνουσα.\par
\elegxos Στην τέμνουσα υπάρχουν δύο κοινά σημεία.

\paradeigma{Κοινά σημεία}
\tixero Μια ευθεία με απόσταση $d=4$ από το κέντρο κύκλου ακτίνας $\rho=4$.\par
\tithelo Πόσα κοινά σημεία έχει.\par
\possindeo $d=\rho$, εφαπτομένη, άρα ένα κοινό σημείο.\par
\apantisi Ένα.\par
\elegxos Το σημείο επαφής είναι το ίχνος της καθέτου από το κέντρο.

\paradeigma{Η εφαπτομένη σε σημείο}
\tixero Ένα σημείο $\mathrm{Μ}$ του κύκλου $(\mathrm{Ο},\rho)$.\par
\tithelo Πώς σχεδιάζω την εφαπτομένη στο $\mathrm{Μ}$;\par
\possindeo Φέρνω την ακτίνα $\mathrm{ΟΜ}$ και την κάθετη στην $\mathrm{ΟΜ}$ στο σημείο $\mathrm{Μ}$.\par
\apantisi Η εφαπτομένη είναι η κάθετη στην ακτίνα $\mathrm{ΟΜ}$ στο $\mathrm{Μ}$.\par
\elegxos Η απόσταση του $\mathrm{Ο}$ από την ευθεία είναι $\mathrm{ΟΜ}=\rho$.

\paradeigma{Πόση απόσταση;}
\tixero Κύκλος $(\mathrm{Ο},5)$ και ευθεία που τον τέμνει σε δύο σημεία.\par
\tithelo Τι ξέρουμε για την απόσταση $d$ του κέντρου από την ευθεία.\par
\possindeo Αφού είναι τέμνουσα, $d<5$.\par
\apantisi $0\le d<5$.\par
\elegxos Αν $d=0$ η ευθεία περνά από το κέντρο και είναι διάμετρος.

\tmima{Ασκήσεις}

\zoni{Το κατάλαβα}{τρεις θέσεις}

\ex Κύκλος $(\mathrm{Ο},3)$. Πες τη θέση της ευθείας όταν η απόσταση του κέντρου από αυτήν είναι: α) $5$ \quad β) $3$ \quad γ) $1{,}5$ \quad δ) $0$ \quad ε) $3{,}01$.

\ex Πόσα κοινά σημεία έχει ευθεία με κύκλο όταν είναι εξωτερική, εφαπτομένη, τέμνουσα;

\ex Σχεδίασε κύκλο με ακτίνα $3$ cm και φέρε μια εξωτερική, μια εφαπτομένη και μια τέμνουσα. Μέτρησε την απόσταση του κέντρου από την καθεμία.

\ex Σχεδίασε την εφαπτομένη ενός κύκλου σε ένα σημείο του $\mathrm{Μ}$.

\ex Κύκλος $(\mathrm{Ο},5)$ και ευθεία που τον τέμνει σε δύο σημεία. Μεταξύ ποιων τιμών βρίσκεται η απόσταση $d$ του κέντρου από την ευθεία;

\ex Σωστό ή λάθος; α) Η εφαπτομένη έχει ένα κοινό σημείο με τον κύκλο. \quad β) Η τέμνουσα έχει τρία κοινά σημεία. \quad γ) Η εφαπτομένη είναι κάθετη στην ακτίνα του σημείου επαφής. \quad δ) Αν $d>\rho$ η ευθεία τέμνει τον κύκλο.

\zoni{Το χρησιμοποιώ}{μονοπάτια και ρόδες}

\ex \textbf{Το μονοπάτι.} Το παρτέρι έχει ακτίνα $2{,}5$ μ. Ένα ευθύγραμμο μονοπάτι περνά σε απόσταση $2{,}5$ μ. από το κέντρο. Πατάμε τα λουλούδια;

\ex \textbf{Τρία μονοπάτια.} Τρία ευθύγραμμα μονοπάτια περνούν σε αποστάσεις $1{,}8$ μ., $2{,}5$ μ. και $3{,}4$ μ. από το κέντρο του παρτεριού ακτίνας $2{,}5$ μ. Ποιο τέμνει το παρτέρι, ποιο το εφάπτεται και ποιο δεν το ακουμπά;

\ex \textbf{Ο τροχός.} Ένας τροχός ακτίνας $33$ cm κυλά σε ευθύγραμμο δρόμο. Πόσο απέχει το κέντρο του από το δρόμο; Πόσα κοινά σημεία έχει ο τροχός με τον δρόμο;

\ex \textbf{Ο φακός.} Ένας φακός φωτίζει ευθεία γραμμή που εφάπτεται σε ένα κυκλικό δωμάτιο. Πού ακουμπά το φως στον τοίχο;

\ex \textbf{Η ρόδα.} Γιατί μια ρόδα κυλά ομαλά πάνω σε επίπεδο δρόμο; (σκέψου ότι ο δρόμος εφάπτεται στη ρόδα).

\zoni{Το σκέφτομαι}{όταν η ευθεία ακουμπά}

\ex Γιατί η εφαπτομένη είναι κάθετη στην ακτίνα του σημείου επαφής; (Σκέψου ότι είναι η μικρότερη απόσταση του κέντρου από την ευθεία.)

\ex Ευθεία περνά από δύο διαφορετικά σημεία του κύκλου. Μπορεί να έχει και τρίτο κοινό σημείο; Εξήγησε.

\ex Πόσες εφαπτομένες διέρχονται από σημείο του κύκλου; Και από σημείο εκτός κύκλου; (Σχεδίασε.)

\ex Δύο παράλληλες ευθείες εφάπτονται σε κύκλο ακτίνας $\rho$. Πόση είναι η απόσταση ανάμεσά τους;

\zoni{Για το σπίτι}{εφαπτόμενα σχήματα}

\ex Βρες στο σπίτι πέντε περιπτώσεις ευθείας και κύκλου με κοινό ένα σημείο (π.χ. τροχός πάνω στο πάτωμα).

\ex Σχεδίασε έναν κύκλο και τρεις εφαπτομένες του σε διαφορετικά σημεία.

\epanalipsi{Β1}

\tmima{Λεξικό του κεφαλαίου}

\renewcommand{\arraystretch}{1.25}
\begin{tabular*}{\linewidth}{@{\extracolsep{\fill}} >{\raggedright\arraybackslash}p{0.28\linewidth} >{\raggedright\arraybackslash}p{0.62\linewidth}}
\toprule
\textbf{Λέξη ή σύμβολο} & \textbf{Με μία πρόταση}\\
\midrule
ευθύγραμμο τμήμα, ευθεία, ημιευθεία & Με δύο άκρα, χωρίς άκρα, με μία αρχή. Από δύο σημεία περνά μία ευθεία.\\
μέσο & Το σημείο $\mathrm{Μ}$ του $\mathrm{ΑΒ}$ με $\mathrm{ΑΜ}=\mathrm{ΜΒ}$.\\
περίμετρος & Το άθροισμα των πλευρών ενός πολυγώνου.\\
γωνία, διχοτόμος & Δύο ημιευθείες με κοινή αρχή· η διχοτόμος τη χωρίζει σε δύο ίσες.\\
οξεία, ορθή, αμβλεία & $<90^\circ$, $=90^\circ$, ανάμεσα σε $90^\circ$ και $180^\circ$.\\
συμπληρωματικές, παραπληρωματικές & Άθροισμα $90^\circ$ και $180^\circ$ αντίστοιχα.\\
κατακορυφήν & Οι απέναντι γωνίες τεμνόμενων ευθειών: ίσες.\\
$\varepsilon_1\parallel\varepsilon_2$, $\varepsilon_1\perp\varepsilon_2$ & Παράλληλες: χωρίς κοινό σημείο. Κάθετες: σχηματίζουν ορθή.\\
κύκλος $(\mathrm{Ο},\rho)$ & Τα σημεία που απέχουν $\rho$ από το $\mathrm{Ο}$. Διάμετρος $\delta=2\rho$.\\
τόξο & Το μέτρο του (σε μοίρες) ισούται με της επίκεντρης γωνίας.\\
εφαπτομένη & Ευθεία με ένα κοινό σημείο με τον κύκλο, κάθετη στην ακτίνα: $d=\rho$.\\
\bottomrule
\end{tabular*}

\tmima{Ασκήσεις επανάληψης}

\ex Πόσες ευθείες ορίζουν τρία μη συνευθειακά σημεία και πόσες τέσσερα σημεία, αν κανένα τριάδα δεν είναι συνευθειακή;

\ex Πόσες διαγώνιους έχουν το πεντάγωνο και το εξάγωνο;

\ex Ένα τμήμα $\mathrm{ΑΒ}=7{,}2$ cm. Βρες το μέσο του και σύγκρινέ το με τμήμα $70$ mm.

\ex Τα $\mathrm{Α},\,\mathrm{Β},\,\Gamma,\,\Delta$ είναι στη σειρά με $\mathrm{ΑΒ}=2{,}5$, $\mathrm{ΒΓ}=3$, $\Gamma\Delta=4$. Βρες τα $\mathrm{ΑΔ}$ και $\mathrm{ΒΔ}$.

\ex Βρες την περίμετρο ενός ορθογωνίου με πλευρές $6{,}5$ και $4$.

\ex Για τη γωνία $124^\circ$ βρες τη διχοτόμο (μέτρο μισής γωνίας) και το παραπλήρωμά της. Βρες το συμπλήρωμα της γωνίας $36^\circ$.

\ex Δύο ευθείες τέμνονται και μία γωνία είναι $65^\circ$. Βρες τις άλλες τρεις.

\ex Τρεις διαδοχικές γωνίες γύρω από σημείο είναι $90^\circ$, $80^\circ$, $110^\circ$. Βρες την τέταρτη.

\ex Τι είδους γωνία σχηματίζουν οι δείκτες του ρολογιού στις $5{:}00$;

\ex Πόσα σημεία τομής έχουν το πολύ $4$ ευθείες του επιπέδου; Από σημείο εκτός ευθείας πόσες παράλληλες διέρχονται;

\ex Τρεις παράλληλες ευθείες απέχουν διαδοχικά $1{,}8$ cm. Πόση απόσταση έχει η πρώτη από την τρίτη;

\ex Στον κύκλο $(\mathrm{Ο},6)$: πόση είναι η διάμετρος; Πες αν τα σημεία με $\mathrm{ΟΑ}=4$ και $\mathrm{ΟΒ}=7$ είναι μέσα ή έξω.

\ex Ένας κύκλος χωρίζεται σε $9$ ίσα τόξα. Πόσες μοίρες είναι το καθένα; Πόσο είναι ένα τόξο που αντιστοιχεί σε επίκεντρη γωνία $150^\circ$;

\ex Κύκλος $(\mathrm{Ο},5)$. Πες τη θέση ευθείας με απόσταση από το κέντρο: α) $5{,}3$ \quad β) $5$ \quad γ) $2$. Πόσα κοινά σημεία έχει καθεμία;

\ex \textbf{Ερωτήσεις αυτοαξιολόγησης (τμήματα, ευθείες).} Σωστό ή λάθος; α) Από δύο σημεία περνά μία ευθεία. \quad β) Η ημιευθεία έχει δύο άκρα. \quad γ) Το μέσο ενός τμήματος είναι μοναδικό. \quad δ) Η περίμετρος ισούται με το άθροισμα των πλευρών. \quad ε) Το ευθύγραμμο τμήμα είναι συντομότερο από κάθε τεθλασμένη με τα ίδια άκρα.

\ex \textbf{Ερωτήσεις αυτοαξιολόγησης (γωνίες, κύκλος).} Σωστό ή λάθος; α) Οι κατακορυφήν γωνίες είναι ίσες. \quad β) Οι συμπληρωματικές γωνίες έχουν άθροισμα $180^\circ$. \quad γ) Δύο ευθείες του επιπέδου ή τέμνονται ή είναι παράλληλες. \quad δ) Η εφαπτομένη είναι κάθετη στην ακτίνα του σημείου επαφής. \quad ε) Η διάμετρος είναι η μεγαλύτερη χορδή. \quad στ) Το κάθετο τμήμα από σημείο σε ευθεία είναι το μεγαλύτερο.

\tmima{Ο χάρτης της επανάληψης}

Μετά τις λύσεις, δες σε ποιες ασκήσεις δυσκολεύτηκες. Ο πίνακας σου λέει ποια παράγραφο να ξαναδιαβάσεις.

\renewcommand{\arraystretch}{1.5}
\begin{tabular*}{\linewidth}{@{\extracolsep{\fill}} >{\raggedright\arraybackslash}p{0.38\linewidth} >{\raggedright\arraybackslash}p{0.52\linewidth}}
\toprule
\textbf{Αν δυσκολεύτηκες στις} & \textbf{ξαναδιάβασε τις παραγράφους}\\
\midrule
ΕΒ1.1 & 1.1 (σημείο, ευθεία, ημιευθεία)\\
ΕΒ1.2 & 1.2 (γωνία, πολύγωνα)\\
ΕΒ1.3, ΕΒ1.4, ΕΒ1.5 & 1.3 και 1.4 (μήκη, μέσο, περίμετρος)\\
ΕΒ1.6 & 1.5 (μοίρες, διχοτόμος) και 1.8\\
ΕΒ1.7 & 1.8 (κατακορυφήν γωνίες)\\
ΕΒ1.8, ΕΒ1.9 & 1.6 και 1.7 (είδη γωνιών, άθροισμα)\\
ΕΒ1.10, ΕΒ1.11 & 1.9 και 1.10 (παράλληλες, αποστάσεις)\\
ΕΒ1.12 & 1.11 (κύκλος)\\
ΕΒ1.13 & 1.12 (τόξα)\\
ΕΒ1.14 & 1.13 (ευθεία και κύκλος)\\
ΕΒ1.15, ΕΒ1.16 & 1.1 έως 1.13\\
\bottomrule
\end{tabular*}

